\chapter{Brave-New Formal Lie Groupoids and Infinitesimal Quotients}
\label{ch:formal-lie-groupoids}

\section{From differential cohesion to effective infinitesimal quotients}
\label{sec:flg-purpose}

The nonlinear infinitesimal geometry used in this chapter has already been
constructed.  Chapter~2 formed the finite spectral infinitesimal Zariski
$\infty$-topos
\[
    \mathcal H_{\inf}
\]
and its endpoint adjoint string
\[
    i_!\dashv i^*\dashv i_*\dashv i^!,
\]
then recognized the induced endofunctors
\[
    \Re:=i_!i^*,
    \qquad
    \Im:=i_*i^*,
    \qquad
    \&:=i_*i^!
\]
as the reduction, infinitesimal-shape, and infinitesimal-flat operations of the
differential-cohesive endpoint semantics.  In particular, for a morphism
$g:E\to B$ in $\mathcal H_{\inf}$, Chapter~2 constructed the relative
infinitesimal shape
\[
    \Im_BE
    :=
    B\times_{\Im B}\Im E
\]
and proved that
\[
    E\longrightarrow \Im_BE\longrightarrow B
\]
is the relative formally \'etale factorization of $g$.  For an ordinary spectral morphism
$X\to S$, the same chapter defined infinitesimal equality by the \v{C}ech
groupoid of the relative infinitesimal-shape map and proved the precise
pointwise statement: on a finite infinitesimal test $j:U\to T$, two
$S$-points of $X$ over $T$ are infinitesimally neighbouring exactly when
their restrictions to $U$ agree, with the specified compatibility over $S$.

Chapter~4 then returned this infinitesimal calculus to the native
spectral Zariski $\infty$-topos.  It constructed the relative de Rham
reflection and, for every
\[
    p:X\longrightarrow S,
\]
its effective-image factorization
\[
    X
    \xrightarrow{e_{X/S}}
    Q_{X/S}
    \xrightarrow{m_{X/S}}
    X_{\mathrm{dR}/S}.
\]
Finite-infinitesimal invariance gives the precise comparison with the
endpoint semantics of Chapter~2: a spectral object is de Rham-local exactly
when it is invariant under every finite spectral infinitesimal substitution,
equivalently when its two endpoint interpretations agree after passage to
the infinitesimal Zariski topos.  Thus relative de Rham reflection is the
native spectral localization characterized by the same finite-infinitesimal
invariance detected by the differential-cohesive endpoint comparison of
Chapter~2.  This is a recognition of a common infinitesimal invariance
condition; it is not an identification of the Chapter~2 infinitesimal-shape
modality and the Chapter~4 de Rham reflector as endofunctors of a single
ambient category.

For an arbitrary morphism
\[
    u:U\longrightarrow Y
\]
over $S$, Chapter~4 also constructed the de Rham--formally \'etale
factorization
\[
    U
    \xrightarrow{\ell_u}
    P^{\mathrm{dR}}(u)
    \xrightarrow{r_u}
    Y,
    \qquad
    P^{\mathrm{dR}}(u)
    :=
    Y\times_{Y_{\mathrm{dR}/S}}U_{\mathrm{dR}/S},
\]
where $\ell_u$ is a relative de Rham equivalence and $r_u$ is formally
\'etale.  Relative de Rham equivalences and formally \'etale morphisms form a
pullback-stable orthogonal factorization system.  This factorization, rather
than a new completion axiom, is the basic operation from which formal
groupoids will be constructed below.

The maximal effective infinitesimal equality relation on $X/S$ is
\[
    \mathcal R^{\mathrm{dR}}_{X/S,\bullet}
    \simeq
    \check C_\bullet
    \left(
        X\xrightarrow{e_{X/S}}Q_{X/S}
    \right),
\]
with
\[
    \mathcal R^{\mathrm{dR}}_{X/S}
    \simeq
    X\times_{Q_{X/S}}X
    \simeq
    X\times_{X_{\mathrm{dR}/S}}X,
    \qquad
    \left|\mathcal R^{\mathrm{dR}}_{X/S,\bullet}\right|
    \simeq Q_{X/S}.
\]
Chapter~5 extracted from this effective relation the crystalline transport
calculus: relative de Rham formal disks, the adjoint pair
\[
    \mathsf{Disk}^{\mathrm{dR},\infty}_{X/S}
    \dashv
    J^\infty_{X/S},
\]
the equivalent algebraic and coalgebraic presentations of effective descent,
and the resulting co-Kleisli differential-operator calculus.

The present chapter inserts one new nonlinear layer into this sequence.  It
studies coherent effective infinitesimal equality relations on $X$ which map
canonically to the maximal de Rham equality relation.  No monomorphism of
relations is required: the homotopy fibres of the comparison carry the
isotropy which a groupoid is meant to retain.

The terminology \emph{brave-new formal Lie groupoid} used below is
anticipatory.  At this intrinsic stage it means a unitwise de Rham-formal
fixed-object groupoid and does not include representability, smoothness,
local finite presentation, tangent perfectness, or a Lie-algebraic
presentation.  Those geometric and linear recognition conditions enter only
on explicitly stated theorem domains later in this chapter and in the
differentiation chapters which follow.  Until such a recognition theorem is
invoked, the word ``Lie'' therefore names the intended differential-geometric
role of the construction rather than an additional field of structure.

With this convention, the resulting unitwise de Rham-formal fixed-object
groupoids will be called brave-new formal Lie groupoids.

The first construction is the fixed-centre specialization of the
de Rham--formally \'etale factorization.  For a morphism
\[
    u:X\longrightarrow Y
\]
over $S$, retain the Chapter~4 factorization object
\[
    P^{\mathrm{dR}}(u)
    =
    Y\times_{Y_{\mathrm{dR}/S}}X_{\mathrm{dR}/S}.
\]
When this construction is regarded functorially on the fixed-source
undercategory, write
\[
    P_X^{\mathrm{dR}}:
    \bigl((\mathcal H_{\Sp})_{/S}\bigr)_{X/}
    \longrightarrow
    \bigl((\mathcal H_{\Sp})_{/S}\bigr)_{X/}
\]
for the endofunctor sending \(u:X\to Y\) to the centred arrow
\[
    X\xrightarrow{\ell_u}P^{\mathrm{dR}}(u).
\]
The map \(\ell_u\) is a relative de Rham equivalence and
\(P^{\mathrm{dR}}(u)\xrightarrow{r_u}Y\) is formally
\'etale.  Orthogonality makes $P_X^{\mathrm{dR}}$ the
coreflection onto centred objects whose centre is already a de Rham
equivalence.  Applying this factorization degreewise along the fully
degenerate unit simplices of a fixed-object groupoid preserves the Segal
pullbacks and produces the idempotent unitwise de Rham formalization used
throughout the chapter.

A formal Lie groupoid is therefore equivalently a fixed-object groupoid
\[
    \mathfrak G_\bullet\rightrightarrows X
\]
whose unit
\[
    u:X\longrightarrow\mathfrak G_1
\]
is a relative de Rham equivalence.  The fixed-point description by unitwise
formalization and the de Rham-equivalence description are proved equivalent
below.  The maximal relation
$\mathcal R^{\mathrm{dR}}_{X/S,\bullet}$ is terminal among such groupoids, hence every
formal Lie groupoid carries a contractibly unique canonical infinitesimal
anchor
\[
    \phi_{\mathfrak G}:
    \mathfrak G_\bullet
    \longrightarrow
    \mathcal R^{\mathrm{dR}}_{X/S,\bullet}.
\]
The anchor is thus a theorem of differential-cohesive formality rather than
additional structure.

Passing to effective quotients realizes the same anchor as
\[
    X
    \xrightarrow{q_{\mathfrak G}}
    Q_{\mathfrak G}
    \xrightarrow{\bar\phi_{\mathfrak G}}
    Q_{X/S},
    \qquad
    \bar\phi_{\mathfrak G}q_{\mathfrak G}\simeq e_{X/S}.
\]
The quotient map $q_{\mathfrak G}$ is itself a relative de Rham equivalence.
Consequently a formal Lie groupoid is equivalently an effective infinitesimal
quotient of $X$: an effective epimorphism $q:X\twoheadrightarrow Q$ inverted
by relative de Rham reflection.

This quotient presentation permits the internal differential-cohesive
manipulations used below.  The chapter will use
\[
    x\sim_{\mathfrak G}y
\]
as internal-language shorthand for an arrow of
$\mathfrak G_1\simeq X\times_{Q_{\mathfrak G}}X$ from $x$ to $y$, and
\[
    x\sim_{\mathrm{dR}/S}y
\]
for an arrow of
$\mathcal R^{\mathrm{dR}}_{X/S}\simeq X\times_{Q_{X/S}}X$.
These symbols introduce no new relation objects: they are dependent notation
for the already constructed effective kernel pairs.  In this notation the
canonical anchor sends
\[
    \gamma:x\sim_{\mathfrak G}y
    \qquad\longmapsto\qquad
    \phi_{\mathfrak G}(\gamma):
    x\sim_{\mathrm{dR}/S}y.
\]
Likewise, formal orbit disks become dependent infinitesimal neighbourhoods,
formal isotropy becomes the kernel of this map on automorphisms, and crystalline
transport becomes substitution in a family descended to
$Q_{\mathfrak G}$.

The chapter also extends the dependent-adjoint calculus of Chapter~5.  For
$q=q_{\mathfrak G}:X\twoheadrightarrow Q_{\mathfrak G}$,
\[
    \Sigma_q\dashv q^*\dashv\Pi_q
\]
produces
\[
    \mathsf{Disk}^{\infty}_{\mathfrak G}:=q^*\Sigma_q,
    \qquad
    J^\infty_{\mathfrak G}:=q^*\Pi_q,
    \qquad
    \mathsf{Disk}^{\infty}_{\mathfrak G}\dashv J^\infty_{\mathfrak G}.
\]
If $E\to X$ is read as a dependent family, their relation-correspondence
formulas will be expressed internally as dependent sum and dependent product
over the infinitesimal equality witnesses
$x\sim_{\mathfrak G}y$.  Effective descent then identifies the resulting
algebra, coalgebra, and transport presentations with objects over
$Q_{\mathfrak G}$.

The stable differential-cohomological reconstruction of Chapter~2 remains a
coefficient-level structure.  Its BNV-type deformation, cycle, and
characteristic sectors are not inserted into the nonlinear definition of a
formal Lie groupoid.  Later in this chapter we record precisely how the crystalline
descent calculus interfaces with that stabilized endpoint semantics, while
keeping separate the additional sliced stabilization needed for a genuinely
relative differential-cohomological refinement.

\paragraph{Standing ambient convention.}
Write
\[
    \mathcal H_{\Sp}
\]
for the spectral Zariski $\infty$-topos fixed in Chapter~1 and fix
\[
    S\in\mathcal H_{\Sp}.
\]
All limits, colimits, groupoid objects, effective epimorphisms,
monomorphisms, realizations, relative de Rham objects, and dependent adjoints
in the intrinsic part of this chapter are formed in
\[
    (\mathcal H_{\Sp})_{/S}
\]
unless explicitly indicated otherwise.  We retain the Chapter~4 notation
\[
    \operatorname{dR}_S(Y)=Y_{\mathrm{dR}/S},
    \qquad
    \eta_{Y/S}^{\mathrm{dR}}:Y\longrightarrow Y_{\mathrm{dR}/S}
\]
for relative de Rham reflection and its unit.  The universe conventions of the
preceding chapters remain in force.

\paragraph{Differential-cohesive internal notation.}
The symbols
\[
    x\sim_{\mathfrak G}y,
    \qquad
    x\sim_{\mathrm{dR}/S}y
\]
are used only after the corresponding relation objects have been constructed.
They abbreviate the dependent fibres of
\[
    (s,t):\mathfrak G_1\to X\times_SX,
    \qquad
    \mathcal R^{\mathrm{dR}}_{X/S}\to X\times_SX.
\]
Dependent sums, dependent products, substitution, and transport written using
these symbols are semantic shorthand for the finite-limit and
$\Sigma\dashv(-)^*\dashv\Pi$ constructions in the ambient
$\infty$-topos.  No separate type-theoretic axiom, Kock--Lawvere principle,
microlinearity hypothesis, or additional cohesive operation is assumed.

This internal style is the same general synthetic pattern used for formal
disks and jet comonads in \cite{KhavkineSchreiber2017}, but every relation,
adjunction, and formal \'etale factorization appearing here is the spectral
construction already established in Chapters~2--5.

\paragraph{Terminology and theorem scope.}
The phrase \emph{formal Lie groupoid} is anticipatory: no representability,
smoothness, finite-presentation condition, tangent perfectness, or
Lie-algebraic presentation is included in the definition.  Representability
enters only in the final first-order interface, where the cotangent complex
and the structured first infinitesimal neighbourhood constructed in
Chapter~2 are available.  Finite local freeness is imposed there only in the
final dualizable tangent range.  A named smoothness condition may be
substituted only through an external recognition theorem stated on its own
theorem domain.  Brackets, spectral Lie objects, partition-Lie algebroids,
Spencer normalization, PBW comparison, and Chevalley--Eilenberg constructions
belong to the differentiation and operation chapters which follow.

\paragraph{No-extra-data convention.}
The standing higher-coherence convention of Chapter~1 remains active.  A
structure produced by a functor, adjunction, finite limit, effective-image
factorization, geometric realization, Beck--Chevalley transformation, mate
correspondence, or descent totalization is retained only through the coherent
structure supplied by that construction.  In particular, the words
``formal'', ``anchor'', ``leaf'', ``bitorsor'', ``adjoint'', ``principal'',
and ``gauge'' below do not name independent admissibility or
presentation fields.  When an auxiliary map is genuinely supplied from
outside the intrinsic construction, it is explicitly identified as external
input and is not included in the definition of a formal Lie groupoid.
Likewise, composition and inversion below refer to the operations canonically
encoded by an internal groupoid object, not to separately chosen fields; and
the word ``represented'' in the final interface means membership in the full
spectral-scheme subcategory, not a choice of representing presentation.

The logical progression is
\[
\boxed{
\begin{aligned}
\text{differential cohesion and infinitesimal shape}
&\longrightarrow
\text{native de Rham reflection and formal \'etalification}
\\
&\longrightarrow
\text{effective infinitesimal equality and crystalline transport}
\\
&\longrightarrow
\text{unitwise de Rham formal Lie groupoids}
\\
&\longrightarrow
\text{effective infinitesimal quotients and canonical anchors}
\\
&\longrightarrow
\text{formal leaves, isotropy, dependent transport, and crystalline descent}.
\end{aligned}}
\]
Differentiation begins only after this nonlinear interface is complete.

\section{Effective groupoids in the ambient $\infty$-topos}
\label{sec:flg-effective-groupoids}

\subsection{Groupoid objects with fixed object space}

Let
\[
    \operatorname{Gpd}((\mathcal H_{\Sp})_{/S})
    \subset
    \operatorname{Fun}(\Delta^{\mathrm{op}},(\mathcal H_{\Sp})_{/S})
\]
denote the full $\infty$-category of groupoid objects in the standard internal
sense \cite[Section~6.1.2]{LurieHTT}.  Thus the Segal pullbacks and the
groupoid condition are properties of the simplicial object itself.  For
$\mathfrak G_\bullet\in\operatorname{Gpd}((\mathcal H_{\Sp})_{/S})$ we write
\[
    \mathfrak G_1
    \mathop{\rightrightarrows}^{s}_{t}
    \mathfrak G_0
\]
for the source and target maps and
\[
    u:\mathfrak G_0\longrightarrow\mathfrak G_1
\]
for the unit.  The induced composition and inversion will be denoted
\[
    \mathsf{comp}_{\mathfrak G}:
    \mathfrak G_1\times_{t,\mathfrak G_0,s}\mathfrak G_1
    \longrightarrow
    \mathfrak G_1,
    \qquad
    \mathsf{inv}_{\mathfrak G}:\mathfrak G_1\xrightarrow{\simeq}\mathfrak G_1.
\]
These are the canonical operations encoded by the groupoid object; they are
not additional fields of structure.  Their identities and all higher
coherences are those supplied by the simplicial object.

\begin{definition}[Fixed-object groupoids]
\label{def:fixed-object-groupoids}
For $X\in(\mathcal H_{\Sp})_{/S}$, let
\[
    \operatorname{Gpd}_X((\mathcal H_{\Sp})_{/S})
\]
denote the homotopy fibre of
\[
    \operatorname{ev}_0:
    \operatorname{Gpd}((\mathcal H_{\Sp})_{/S})
    \longrightarrow
    (\mathcal H_{\Sp})_{/S}
\]
over $X$.  Thus an object carries the fibre identification
$\mathfrak G_0\simeq X$, and a morphism has degree-zero component identified
with $\operatorname{id}_X$ through that fibre data.
\end{definition}

\begin{remark}
The fixed-object formulation is useful because the later differentiated
object is infinitesimal structure over the fixed geometric object $X$.  It
also makes the one-object case transparent: when $X=S$, a group object in
$(\mathcal H_{\Sp})_{/S}$ is precisely a one-object
groupoid in $\operatorname{Gpd}_S((\mathcal H_{\Sp})_{/S})$.
\end{remark}

\subsection{Effectivity and quotient objects}

Write
\[
    \operatorname{EffEpi}((\mathcal H_{\Sp})_{/S})
    \subset
    \operatorname{Fun}(\Delta^1,(\mathcal H_{\Sp})_{/S})
\]
for the full $\infty$-subcategory spanned by effective epimorphisms.

\begin{theorem}[Effectivity equivalence for internal groupoids]
\label{thm:groupoid-effectivity}
Realization and \v{C}ech nerve induce inverse equivalences
\[
    \operatorname{Gpd}((\mathcal H_{\Sp})_{/S})
    \xrightarrow{\ \simeq\ }
    \operatorname{EffEpi}((\mathcal H_{\Sp})_{/S}),
    \qquad
    \mathfrak G_\bullet
    \longmapsto
    \left(
        \mathfrak G_0\longrightarrow
        \lvert\mathfrak G_\bullet\rvert
    \right),
\]
and
\[
    \operatorname{EffEpi}((\mathcal H_{\Sp})_{/S})
    \xrightarrow{\ \simeq\ }
    \operatorname{Gpd}((\mathcal H_{\Sp})_{/S}),
    \qquad
    (q:U\twoheadrightarrow Q)
    \longmapsto
    \check C_\bullet(q).
\]
The equivalence is canonically over the degree-zero/domain object:
\[
\begin{tikzcd}
\operatorname{Gpd}((\mathcal H_{\Sp})_{/S})
    \arrow[r,"\simeq"]
    \arrow[d,"\operatorname{ev}_0"']
&
\operatorname{EffEpi}((\mathcal H_{\Sp})_{/S})
    \arrow[d,"\operatorname{dom}"]
\\
(\mathcal H_{\Sp})_{/S} \arrow[r,equal] & (\mathcal H_{\Sp})_{/S} .
\end{tikzcd}
\]
Consequently, for every groupoid object,
\[
    q_{\mathfrak G}:
    \mathfrak G_0\twoheadrightarrow
    Q_{\mathfrak G}:=\lvert\mathfrak G_\bullet\rvert
\]
is an effective epimorphism and
\[
    \mathfrak G_\bullet
    \simeq
    \check C_\bullet(q_{\mathfrak G}).
\]
\end{theorem}

\begin{proof}
This is the effectivity theorem for groupoid objects in an $\infty$-topos
\cite[Section~6.1.2]{LurieHTT}.  Every internal groupoid in an
$\infty$-topos is effective, so its augmentation to its geometric realization
is an effective epimorphism and the original simplicial object is canonically
its \v{C}ech nerve.  Conversely, the \v{C}ech nerve of an effective
epimorphism is a groupoid object and realizes to its target.  These
constructions are functorial on the corresponding $\infty$-categories, hence
identify mapping spaces and all higher coherence, not merely objects up to
equivalence.  Their degree-zero/domain compatibility is tautological from the
two constructions.
\end{proof}

\begin{definition}[Effective quotient of a groupoid]
\label{def:groupoid-quotient}
For $\mathfrak G_\bullet\in\operatorname{Gpd}_X((\mathcal H_{\Sp})_{/S})$, define
\[
    Q_{\mathfrak G}:=\lvert\mathfrak G_\bullet\rvert,
    \qquad
    q_{\mathfrak G}:X\twoheadrightarrow Q_{\mathfrak G}.
\]
We call $Q_{\mathfrak G}$ its effective quotient and
$q_{\mathfrak G}$ its quotient presentation.  The word ``presentation'' here
denotes only this canonically constructed effective epimorphism; no atlas,
choice of representatives, or additional quotient datum is retained.
\end{definition}

\begin{proposition}[Fixed-object groupoids and effective quotients]
\label{prop:groupoids-effective-quotients-equivalence}
Let
\[
    \operatorname{EffEpi}(X)
    \subset
    \bigl((\mathcal H_{\Sp})_{/S}\bigr)_{X/}
\]
be the full $\infty$-subcategory spanned by effective epimorphisms
$q:X\twoheadrightarrow Q$.  Taking the homotopy fibre over $X$ in
Theorem~\ref{thm:groupoid-effectivity} gives a canonical equivalence
\[
    \boxed{
    \operatorname{Gpd}_X((\mathcal H_{\Sp})_{/S})
    \simeq
    \operatorname{EffEpi}(X).}
\]
Under this equivalence,
\[
    \mathfrak G_n
    \simeq
    \underbrace{
        X\times_{Q_{\mathfrak G}}\cdots
        \times_{Q_{\mathfrak G}}X
    }_{n+1\text{ factors}}.
\]
\end{proposition}

\begin{proof}
Take the homotopy fibre over $X$ of the square in
Theorem~\ref{thm:groupoid-effectivity}.  The fibre of
$\operatorname{ev}_0$ is $\operatorname{Gpd}_X((\mathcal H_{\Sp})_{/S})$, while the fibre
of the domain functor is exactly the full effective-epimorphism subcategory
$\operatorname{EffEpi}(X)\subset\bigl((\mathcal H_{\Sp})_{/S}\bigr)_{X/}$.  Equivalences of
$\infty$-categories are preserved by homotopy fibres.  The degreewise formula
is the \v{C}ech-nerve description.
\end{proof}

\begin{remark}[Why quotient language is primary]
\label{rem:quotient-language-primary}
The quotient presentation
\[
    X\twoheadrightarrow Q_{\mathfrak G}
\]
does not truncate or coarse-grain the internal groupoid: it is an equivalent
presentation of the entire fixed-object groupoid in the ambient
$\infty$-topos, including its mapping spaces and higher coherence.
\end{remark}


\section{Centred formal \'etalification and unitwise de Rham formalization}
\label{sec:flg-unitwise-formalization}

The formalization of groupoids is controlled by the de Rham--formally
\'etale factorization of Chapter~4, specialized to maps with a fixed source.
We therefore work first in the undercategory
\[
    \bigl((\mathcal H_{\Sp})_{/S}\bigr)_{X/}
    =
    ((\mathcal H_{\Sp})_{/S})_{X/}.
\]
The point of the fixed centre is not to introduce a second completion
procedure.  It is to retain the source $X$ while applying the already
constructed relative formal \'etalification functorially to the target.

\subsection{Centred de Rham formal \'etalification}

The coreflection used below is a formal consequence of orthogonal
factorization once the source object is fixed.

\begin{lemma}[Fixed-source coreflection from an orthogonal factorization system]
\label{lem:fixed-source-ofs-coreflection}
Let $\mathcal C$ be an $\infty$-category equipped with a functorial
orthogonal factorization system $(\mathcal L,\mathcal R)$, fix
$X\in\mathcal C$, and let
\[
    (\mathcal C_{X/})_{\mathcal L}
    \subset
    \mathcal C_{X/}
\]
be the full subcategory spanned by the arrows $v:X\to F$ which lie in
$\mathcal L$.  Write the functorial factorization of
$u:X\to Y$ as
\[
    X\xrightarrow{\ell_u}P_X(u)\xrightarrow{r_u}Y,
    \qquad
    \ell_u\in\mathcal L,\quad r_u\in\mathcal R.
\]
Then
\[
    P_X:\mathcal C_{X/}\longrightarrow
    (\mathcal C_{X/})_{\mathcal L}
\]
is right adjoint to the fully faithful inclusion.  Equivalently, for
$v:X\to F$ in $\mathcal L$, composition with $r_u$ induces a canonical
equivalence
\[
    \Map_{\mathcal C_{X/}}
    (F,P_X(u))
    \xrightarrow{\simeq}
    \Map_{\mathcal C_{X/}}(F,Y).
\]
The induced comonad on $\mathcal C_{X/}$ is idempotent.
\end{lemma}

\begin{proof}
Fix a morphism $f:F\to Y$ under $X$.  The homotopy fibre over $f$ of
\[
    \Map_{\mathcal C_{X/}}(F,P_X(u))
    \longrightarrow
    \Map_{\mathcal C_{X/}}(F,Y)
\]
is the space of diagonal fillers in
\[
\begin{tikzcd}
X \arrow[r,"\ell_u"] \arrow[d,"v"'] &
P_X(u) \arrow[d,"r_u"]\\
F \arrow[r,"f"'] &
Y .
\end{tikzcd}
\]
Because $v\in\mathcal L$ and $r_u\in\mathcal R$, orthogonality makes this
filler space contractible.  Hence the displayed mapping-space map is an
equivalence, naturally in $F$ and $Y$, which is the required adjunction.
The inclusion is fully faithful, so the unit of this adjunction is an
equivalence on its source; therefore the associated comonad is idempotent.
All coherence is the coherence of the functorial factorization and the
orthogonality equivalence.
\end{proof}

\begin{definition}[Centred relative de Rham formal \'etalification]
\label{def:centred-dr-formalization}
Let
\[
    u:X\longrightarrow Y
\]
be an object of $\bigl((\mathcal H_{\Sp})_{/S}\bigr)_{X/}$.  Retain the Chapter~4 notation
\[
    P^{\mathrm{dR}}(u)
    =
    Y\times_{Y_{\mathrm{dR}/S}}X_{\mathrm{dR}/S}.
\]
Let
\[
    P_X^{\mathrm{dR}}:
    \bigl((\mathcal H_{\Sp})_{/S}\bigr)_{X/}
    \longrightarrow
    \bigl((\mathcal H_{\Sp})_{/S}\bigr)_{X/}
\]
denote the fixed-source endofunctor whose value at
$u:X\to Y$ is the centred arrow
\[
    X\xrightarrow{\ell_u}P^{\mathrm{dR}}(u).
\]
The established maps
\[
    X
    \xrightarrow{\ell_u}
    P^{\mathrm{dR}}(u)
    \xrightarrow{r_u}
    Y,
    \qquad
    r_u\ell_u\simeq u,
\]
give the canonical de Rham--formally \'etale factorization.
\end{definition}

\begin{proposition}[Differential-cohesive factorization of a centred object]
\label{prop:centred-formalization-recognition}
For every centred object $u:X\to Y$:
\begin{enumerate}
    \item the formalized centre
    \[
        \ell_u:X\longrightarrow P^{\mathrm{dR}}(u)
    \]
    is a relative de Rham equivalence;
    \item the projection
    \[
        r_u:P^{\mathrm{dR}}(u)\longrightarrow Y
    \]
    is formally \'etale;
    \item there is a canonical equivalence
    \[
        \operatorname{dR}_S(P^{\mathrm{dR}}(u))
        \xrightarrow{\simeq}
        \operatorname{dR}_S(X)
    \]
    carrying $\operatorname{dR}_S(\ell_u)$ to the identity;
    \item the following are equivalent:
    \begin{enumerate}
        \item $u:X\to Y$ is a relative de Rham equivalence;
        \item $r_u:P^{\mathrm{dR}}(u)\to Y$ is an equivalence.
    \end{enumerate}
\end{enumerate}
\end{proposition}

\begin{proof}
The first two statements are exactly the two classes in the
de Rham--formally \'etale factorization system of Chapter~4 applied to
$u$.  For completeness, left exactness and idempotence give
\[
\begin{aligned}
    \operatorname{dR}_S(P^{\mathrm{dR}}(u))
    &\simeq
    \operatorname{dR}_S(Y)
    \times_{\operatorname{dR}_S(\operatorname{dR}_S(Y))}
    \operatorname{dR}_S(\operatorname{dR}_S(X))\\
    &\simeq
    \operatorname{dR}_S(Y)
    \times_{\operatorname{dR}_S(Y)}
    \operatorname{dR}_S(X)\\
    &\simeq
    \operatorname{dR}_S(X),
\end{aligned}
\]
and tracing the centre identifies
$\operatorname{dR}_S(\ell_u)$ with the identity.

If $u$ is a de Rham equivalence, then
$\operatorname{dR}_S(u)$ is an equivalence and $r_u$ is its base change along
$\eta_{Y/S}^{\mathrm{dR}}:Y\to Y_{\mathrm{dR}/S}$, hence is an equivalence.  Conversely, if
$r_u$ is an equivalence, then
\[
    u\simeq r_u\ell_u
\]
is a composite of a de Rham equivalence with an equivalence and is therefore a
de Rham equivalence.
\end{proof}

\begin{definition}[Centred de Rham-formal object]
\label{def:centred-formal-object}
A centred object $u:X\to Y$ is \emph{de Rham-formal along $X$} when its centre
is a relative de Rham equivalence.  Equivalently, by
Proposition~\ref{prop:centred-formalization-recognition},
\[
    r_u:P^{\mathrm{dR}}(u)\xrightarrow{\simeq}Y.
\]
Write
\[
    \bigl((\mathcal H_{\Sp})_{/S}\bigr)_{X/}^{\mathrm{for}}
    \subset
    \bigl((\mathcal H_{\Sp})_{/S}\bigr)_{X/}
\]
for the full subcategory of de Rham-formal centred objects.
\end{definition}

\begin{proposition}[Pullback preservation]
\label{prop:centred-formalization-pullbacks}
The assignment
\[
    P_X^{\mathrm{dR}}:\bigl((\mathcal H_{\Sp})_{/S}\bigr)_{X/}\longrightarrow\bigl((\mathcal H_{\Sp})_{/S}\bigr)_{X/}
\]
is functorial, preserves pullbacks, and fixes the initial centred object:
\[
    P_X^{\mathrm{dR}}(\operatorname{id}_X)
    \simeq
    \operatorname{id}_X
\]
in \(\bigl((\mathcal H_{\Sp})_{/S}\bigr)_{X/}\).  Equivalently,
\[
    P^{\mathrm{dR}}(\operatorname{id}_X)\simeq X.
\]
The morphisms $r_u:P^{\mathrm{dR}}(u)\to Y$ assemble into a natural transformation
\[
    r:P_X^{\mathrm{dR}}\longrightarrow\operatorname{id}_{\bigl((\mathcal H_{\Sp})_{/S}\bigr)_{X/}}.
\]
\end{proposition}

\begin{proof}
A morphism
\[
    f:(X\xrightarrow{u}Y)\longrightarrow(X\xrightarrow{v}Z)
\]
in the undercategory satisfies $fu\simeq v$.  Naturality of relative de Rham
reflection gives a morphism between the defining pullback diagrams and hence
\[
    P_X^{\mathrm{dR}}(f):P^{\mathrm{dR}}(u)\longrightarrow P^{\mathrm{dR}}(v).
\]
Functoriality and all higher composition coherence are supplied by
functoriality of finite limits.  Moreover
\[
    P^{\mathrm{dR}}(\operatorname{id}_X)
    =
    X\times_{\operatorname{dR}_S(X)}\operatorname{dR}_S(X)
    \simeq X.
\]

Let
\[
\begin{tikzcd}
Y \arrow[r] \arrow[d] &
Y_1 \arrow[d]\\
Y_2 \arrow[r] &
Y_0
\end{tikzcd}
\]
be a pullback square in $\bigl((\mathcal H_{\Sp})_{/S}\bigr)_{X/}$.
Write
\[
    u:X\to Y,
    \qquad
    u_i:X\to Y_i\quad(0\leq i\leq2)
\]
for the corresponding centre maps.  Its underlying square is Cartesian and
all four centres come from the same object $X$.  Left exactness of relative
de Rham reflection gives
\[
    \operatorname{dR}_S(Y)
    \simeq
    \operatorname{dR}_S(Y_1)
    \times_{\operatorname{dR}_S(Y_0)}
    \operatorname{dR}_S(Y_2).
\]
Pullback pasting then yields
\[
\begin{aligned}
    P^{\mathrm{dR}}(u)
    &=
    Y\times_{\operatorname{dR}_S(Y)}\operatorname{dR}_S(X)\\
    &\simeq
    P^{\mathrm{dR}}(u_1)
    \times_{P^{\mathrm{dR}}(u_0)}
    P^{\mathrm{dR}}(u_2).
\end{aligned}
\]
Naturality of $r$ follows from the same diagrams.
\end{proof}

\begin{remark}[Pullback exactness versus left exactness]
\label{rem:centred-formalization-not-left-exact}
The preceding proposition is intentionally a pullback statement.  The
endofunctor $P_X^{\mathrm{dR}}$ need not preserve the terminal object of the undercategory
and therefore is not, in general, left exact as an endofunctor of
$\bigl((\mathcal H_{\Sp})_{/S}\bigr)_{X/}$.  Indeed, if
\[
    p:X\to S
\]
is the terminal centred object, then
\[
    P^{\mathrm{dR}}(p)
    \simeq
    S\times_{\operatorname{dR}_S(S)}\operatorname{dR}_S(X)
    \simeq
    \operatorname{dR}_S(X),
\]
which need not be equivalent to $S$.  The groupoid formalization below uses
only pullback preservation together with
$P^{\mathrm{dR}}(\operatorname{id}_X)\simeq X$, exactly the properties needed to
preserve the Segal pullbacks.
\end{remark}

\begin{theorem}[Centred formalization coreflection]
\label{thm:centred-formalization-coreflection}
The inclusion
\[
    \bigl((\mathcal H_{\Sp})_{/S}\bigr)_{X/}^{\mathrm{for}}
    \hookrightarrow
    \bigl((\mathcal H_{\Sp})_{/S}\bigr)_{X/}
\]
admits $P_X^{\mathrm{dR}}$ as a right adjoint.  Equivalently, if
$v:X\to F$ is a de Rham-formal centred object, then composition with
$r_u$ induces a canonical equivalence
\[
    \Map_{\bigl((\mathcal H_{\Sp})_{/S}\bigr)_{X/}}
    (F,P^{\mathrm{dR}}(u))
    \xrightarrow{\simeq}
    \Map_{\bigl((\mathcal H_{\Sp})_{/S}\bigr)_{X/}}
    (F,Y).
\]
\end{theorem}

\begin{proof}
Apply Lemma~\ref{lem:fixed-source-ofs-coreflection} in
$\mathcal C=(\mathcal H_{\Sp})_{/S}$ to the pullback-stable orthogonal factorization
system of Chapter~4 whose left class is the relative de Rham equivalences and
whose right class is the formally \'etale morphisms.  For
$u:X\to Y$ its functorial factorization object is exactly
\[
    P_X(u)
    =
    P^{\mathrm{dR}}(u),
\]
by Definition~\ref{def:centred-dr-formalization}.  The full left-class
subcategory of the undercategory is precisely
$\bigl((\mathcal H_{\Sp})_{/S}\bigr)_{X/}^{\mathrm{for}}$.  The fixed-source lemma therefore gives the
displayed mapping-space equivalence and the asserted right adjoint with
counit $r$.
\end{proof}

\begin{remark}[Relation with relative infinitesimal shape]
The formula
\[
    P^{\mathrm{dR}}(u)
    =
    Y\times_{Y_{\mathrm{dR}/S}}X_{\mathrm{dR}/S}
\]
is the native spectral counterpart of the relative infinitesimal-shape object
\[
    \Im_YX
    =
    Y\times_{\Im Y}\Im X
\]
constructed in Chapter~2.  The present chapter uses the de Rham realization
because all groupoids, quotients, and dependent adjoints are formed in the
native spectral Zariski $\infty$-topos.  The comparison with
$\mathcal H_{\inf}$ is the finite-infinitesimal endpoint recognition
of Chapter~4, not an identification of the two ambient categories.
\end{remark}

\begin{theorem}[Idempotent centred formalization]
\label{thm:centred-formalization-idempotent}
The coreflection of
Theorem~\ref{thm:centred-formalization-coreflection} equips $P_X^{\mathrm{dR}}$ with a
canonical idempotent comonad structure whose counit is
$r:P_X^{\mathrm{dR}}\to\operatorname{id}$.  In particular there is a canonical natural
equivalence
\[
    P_X^{\mathrm{dR}}P_X^{\mathrm{dR}}
    \xrightarrow{\simeq}
    P_X^{\mathrm{dR}}.
\]
\end{theorem}

\begin{proof}
Let
\[
    i:
    \bigl((\mathcal H_{\Sp})_{/S}\bigr)_{X/}^{\mathrm{for}}
    \hookrightarrow
    \bigl((\mathcal H_{\Sp})_{/S}\bigr)_{X/}
\]
denote the fully faithful inclusion.  The preceding theorem identifies
$P_X^{\mathrm{dR}}$ with the composite endofunctor of the adjunction
\[
    i\dashv P_X^{\mathrm{dR}}
\]
after regarding the right adjoint as landing in
$\bigl((\mathcal H_{\Sp})_{/S}\bigr)_{X/}^{\mathrm{for}}$.  The associated adjunction therefore
supplies the comonad structure on the endofunctor $P_X^{\mathrm{dR}}$, with counit
$r$.

Because $i$ is fully faithful, the unit of the adjunction is an equivalence.
Equivalently, every object in the image of $P_X^{\mathrm{dR}}$ is fixed by $P_X^{\mathrm{dR}}$.  Hence
the comultiplication is an equivalence and the comonad is idempotent, giving
the displayed natural equivalence.
\end{proof}

\begin{proposition}[Functoriality in the centre]
\label{prop:centred-formalization-varying-centre}
Suppose
\[
\begin{tikzcd}
X \arrow[r,"u"] \arrow[d,"f"'] &
Y \arrow[d,"g"]\\
X' \arrow[r,"u'"'] &
Y'
\end{tikzcd}
\]
is a commutative square in $(\mathcal H_{\Sp})_{/S}$.  Then there is a canonical morphism
\[
    P^{\mathrm{dR}}(u)
    \longrightarrow
    P^{\mathrm{dR}}(u')
\]
whose components are $g:Y\to Y'$ and
$\operatorname{dR}_S(f):\operatorname{dR}_S(X)\to \operatorname{dR}_S(X')$.  These morphisms preserve identities,
composition, and all higher functorial coherences.
\end{proposition}

\begin{proof}
Naturality of $\eta$ and commutativity of the given square supply a morphism
between the two defining pullback diagrams.  Functoriality of finite limits
gives the asserted map and all composition coherences.
\end{proof}

\begin{proposition}[Base change of centred formal \'etalification]
\label{prop:centred-formalization-base-change}
Let $c:S'\to S$, put
\[
    X':=X\times_SS',
    \qquad
    Y':=Y\times_SS',
\]
and let $u':X'\to Y'$ be the pullback of $u$.  Then there is a canonical
equivalence
\[
    P^{\mathrm{dR}}(u')
    \simeq
    P^{\mathrm{dR}}(u)\times_SS',
\]
compatible with $r$, the comonad structure, identities, iterated base
change, and higher pullback-pasting coherences.
\end{proposition}

\begin{proof}
The relative de Rham base-change theorem gives
\[
    \operatorname{dR}_{S'}(X')\simeq \operatorname{dR}_S(X)\times_SS',
    \qquad
    \operatorname{dR}_{S'}(Y')\simeq \operatorname{dR}_S(Y)\times_SS'.
\]
Substitute these into the defining pullback and commute finite limits with
base change.  Every displayed structure is constructed from these same
pullbacks, so the compatibility statements follow formally.
\end{proof}

\subsection{Unitwise de Rham formalization of fixed-object groupoids}

Fix
\[
    \mathfrak G_\bullet
    \in
    \operatorname{Gpd}_X((\mathcal H_{\Sp})_{/S}).
\]
For every $n\geq0$, let
\[
    u_n:X\longrightarrow\mathfrak G_n
\]
be the fully degenerate unit simplex.

\begin{definition}[Unitwise relative de Rham formalization]
\label{def:unitwise-groupoid-formalization}
Define
\[
    \bigl(P_X^{\mathrm{dR}}\mathfrak G_\bullet\bigr)_n
    :=
    P^{\mathrm{dR}}(u_n)
    =
    \mathfrak G_n
    \times_{(\mathfrak G_n)_{\mathrm{dR}/S}}
    X_{\mathrm{dR}/S}.
\]
Write
\[
    r_{\mathfrak G,n}
    :=
    r_{u_n}:
    \bigl(P_X^{\mathrm{dR}}\mathfrak G_\bullet\bigr)_n
    \longrightarrow
    \mathfrak G_n
\]
for the first projection.  In accordance with the inherited fixed-source
notation, we use the same symbol \(P_X^{\mathrm{dR}}\) for this degreewise
extension to fixed-object simplicial objects and groupoids; no separate
completion notation is introduced.
\end{definition}

\begin{proposition}[Unitwise formalization is a groupoid]
\label{prop:unitwise-formalization-groupoid}
The objects
$\bigl(P_X^{\mathrm{dR}}\mathfrak G_\bullet\bigr)_n$ assemble canonically into a
groupoid object
\[
    P_X^{\mathrm{dR}}(\mathfrak G_\bullet)
    \rightrightarrows X.
\]
The projections assemble into a natural fixed-object groupoid morphism
\[
    r_{\mathfrak G}:
    P_X^{\mathrm{dR}}(\mathfrak G_\bullet)
    \longrightarrow
    \mathfrak G_\bullet.
\]
Thus unitwise de Rham formalization defines an endofunctor
\[
    P_X^{\mathrm{dR}}:
    \operatorname{Gpd}_X((\mathcal H_{\Sp})_{/S})
    \longrightarrow
    \operatorname{Gpd}_X((\mathcal H_{\Sp})_{/S}).
\]
For $n\ge2$ there are canonical Segal equivalences
\[
    \bigl(P_X^{\mathrm{dR}}\mathfrak G_\bullet\bigr)_n
    \simeq
    \underbrace{
    \bigl(P_X^{\mathrm{dR}}\mathfrak G_\bullet\bigr)_1
    \times_X\cdots\times_X
    \bigl(P_X^{\mathrm{dR}}\mathfrak G_\bullet\bigr)_1
    }_{n\text{ factors}}.
\]
\end{proposition}

\begin{proof}
For every simplicial operator
\[
    \alpha:[m]\longrightarrow[n]
\]
the fully degenerate simplices satisfy
\[
    \alpha^*u_n\simeq u_m.
\]
Proposition~\ref{prop:centred-formalization-varying-centre}, applied with centre
map $\operatorname{id}_X$, therefore gives a canonical morphism
\[
    P^{\mathrm{dR}}(u_n)
    \longrightarrow
    P^{\mathrm{dR}}(u_m).
\]
Functoriality in $\alpha$, including identities, composition, and all higher
simplicial coherence, is inherited from the functoriality of centred
formalization.  Hence the objects
$\bigl(P_X^{\mathrm{dR}}\mathfrak G_\bullet\bigr)_n$ form a simplicial object, and
naturality of the projections $r_{\mathfrak G,n}$ makes
\[
    r_{\mathfrak G}:
    P_X^{\mathrm{dR}}(\mathfrak G_\bullet)
    \longrightarrow
    \mathfrak G_\bullet
\]
a simplicial morphism.  Since
\[
    P^{\mathrm{dR}}(\operatorname{id}_X)\simeq X,
\]
its degree-zero object is canonically $X$ and the degree-zero component of
$r_{\mathfrak G}$ is $\operatorname{id}_X$.

For $n\ge2$, the Segal equivalence for $\mathfrak G_\bullet$ is
\[
    \mathfrak G_n
    \simeq
    \underbrace{
        \mathfrak G_1\times_X\cdots\times_X\mathfrak G_1
    }_{n\text{ factors}},
\]
where every term is regarded as a centred object through the corresponding
fully degenerate unit simplex.  Proposition~
\ref{prop:centred-formalization-pullbacks} preserves these pullbacks and fixes
$(X,\operatorname{id}_X)$.  Therefore
\[
\begin{aligned}
    \bigl(P_X^{\mathrm{dR}}\mathfrak G_\bullet\bigr)_n
    &=
    P^{\mathrm{dR}}(u_n)\\
    &\simeq
    \underbrace{
        P^{\mathrm{dR}}(u_1)
        \times_X\cdots\times_X
        P^{\mathrm{dR}}(u_1)
    }_{n\text{ factors}}\\
    &=
    \underbrace{
        \bigl(P_X^{\mathrm{dR}}\mathfrak G_\bullet\bigr)_1
        \times_X\cdots\times_X
        \bigl(P_X^{\mathrm{dR}}\mathfrak G_\bullet\bigr)_1
    }_{n\text{ factors}}.
\end{aligned}
\]
Thus the formalized simplicial object satisfies the Segal conditions.

It remains to verify invertibility of its arrows together with the coherent
groupoid identities.  Let
\[
    \mathsf{comp}_{\mathfrak G}:
    \mathfrak G_1\times_{t,X,s}\mathfrak G_1
    \longrightarrow
    \mathfrak G_1,
    \qquad
    \mathsf{inv}_{\mathfrak G}:
    \mathfrak G_1\longrightarrow\mathfrak G_1
\]
denote the canonical composition and inversion operations encoded by the
groupoid object.  The composable-arrow object is centred by
\[
    u_2:X\longrightarrow
    \mathfrak G_1\times_{t,X,s}\mathfrak G_1,
    \qquad
    x\longmapsto(u(x),u(x)),
\]
and both $\mathsf{comp}_{\mathfrak G}$ and $\mathsf{inv}_{\mathfrak G}$ preserve the corresponding unit centres.  Applying
$P_X^{\mathrm{dR}}$ therefore gives morphisms
\[
    P^{\mathrm{dR}}(u_2)
    \longrightarrow
    P^{\mathrm{dR}}(u)
\]
and
\[
    P^{\mathrm{dR}}(u)
    \longrightarrow
    P^{\mathrm{dR}}(u).
\]
Pullback preservation identifies the first source canonically with
\[
    \bigl(P_X^{\mathrm{dR}}\mathfrak G_\bullet\bigr)_1
    \times_{t,X,s}
    \bigl(P_X^{\mathrm{dR}}\mathfrak G_\bullet\bigr)_1.
\]
Hence these maps define the formalized composition and inversion.

The unit, associativity, left-inverse, right-inverse, and all higher groupoid
coherences of $\mathfrak G_\bullet$ are homotopies between composites of the
same finite-pullback diagrams.  Since $P_X^{\mathrm{dR}}$ is a functor and preserves those
pullbacks, it carries every such homotopy and every higher compatibility to
the corresponding formalized diagram.  Thus the formalized composition,
unit, and inversion satisfy the complete coherent groupoid identities.

Consequently
$P_X^{\mathrm{dR}}(\mathfrak G_\bullet)$ is a groupoid object over
the fixed object space $X$.  The construction is functorial in
fixed-object groupoid morphisms because every step is functorial in the
centred simplicial diagram.  No composition, inverse, or coherence datum is
retained independently of the resulting groupoid object.
\end{proof}

\begin{theorem}[Idempotence of unitwise de Rham formalization]
\label{thm:unitwise-formalization-idempotent}
Unitwise de Rham formalization is an idempotent comonad on
$\operatorname{Gpd}_X((\mathcal H_{\Sp})_{/S})$.  In particular there is a canonical
natural equivalence
\[
    P_X^{\mathrm{dR}}\!\left(
        P_X^{\mathrm{dR}}(\mathfrak G_\bullet)
    \right)
    \simeq
    P_X^{\mathrm{dR}}(\mathfrak G_\bullet),
\]
and
\[
    \left(
        \bigl(P_X^{\mathrm{dR}}\mathfrak G_\bullet\bigr)_n
    \right)_{\mathrm{dR}/S}
    \simeq
    X_{\mathrm{dR}/S}
\]
for every $n$.
\end{theorem}

\begin{proof}
Apply Theorem~\ref{thm:centred-formalization-idempotent} degreewise.  Since the
comonad maps are natural for morphisms of centred objects, they commute with
all simplicial operators and therefore define the stated groupoid-level
comonad.
\end{proof}

\begin{definition}[Unitwise de Rham-formal groupoid]
\label{def:unitwise-formal-groupoid}
A fixed-object groupoid
\[
    \mathfrak G_\bullet\rightrightarrows X
\]
is \emph{unitwise de Rham-formal}, or simply \emph{formal}, if its
degree-one unit
\[
    u_1:X\longrightarrow\mathfrak G_1
\]
is a relative de Rham equivalence.  Write
\[
    \operatorname{FormLieGpd}(X/S)
    \subset
    \operatorname{Gpd}_X((\mathcal H_{\Sp})_{/S})
\]
for the full subcategory of formal fixed-object groupoids.  This is the
single category notation retained for formal groupoids in the chapter.
\end{definition}

\begin{proposition}[Degree-one and fixed-point recognition of formality]
\label{prop:degree-one-formality}
For a fixed-object groupoid $\mathfrak G_\bullet$, the following are
equivalent:
\begin{enumerate}
    \item $\mathfrak G_\bullet$ is formal;
    \item the unit
    \[
        u_1:X\longrightarrow\mathfrak G_1
    \]
    is a relative de Rham equivalence;
    \item the localized unit
    \[
        \operatorname{dR}_S(u_1):
        X_{\mathrm{dR}/S}
        \xrightarrow{\simeq}
        (\mathfrak G_1)_{\mathrm{dR}/S}
    \]
    is an equivalence;
    \item the degree-one formalization counit
    \[
        \mathfrak G_1
        \times_{(\mathfrak G_1)_{\mathrm{dR}/S}}
        X_{\mathrm{dR}/S}
        \longrightarrow
        \mathfrak G_1
    \]
    is an equivalence;
    \item the simplicial formalization counit
    \[
        r_{\mathfrak G}:
        P_X^{\mathrm{dR}}(\mathfrak G_\bullet)
        \xrightarrow{\simeq}
        \mathfrak G_\bullet
    \]
    is an equivalence.
\end{enumerate}
\end{proposition}

\begin{proof}
The equivalence of (1)--(3) is the definition of a relative de Rham
equivalence.  Proposition~\ref{prop:centred-formalization-recognition} identifies
(3) with (4).  Both $\mathfrak G_\bullet$ and
$P_X^{\mathrm{dR}}(\mathfrak G_\bullet)$ are Segal groupoids with
degree zero $X$, and under the Segal equivalences
$r_{\mathfrak G,n}$ is the $n$-fold fibre product over $X$ of
$r_{\mathfrak G,1}$.  Hence (4) and (5) are equivalent.
\end{proof}

\begin{theorem}[Formalization coreflection]
\label{thm:formalization-coreflection}
The inclusion
\[
    \operatorname{FormLieGpd}(X/S)
    \hookrightarrow
    \operatorname{Gpd}_X((\mathcal H_{\Sp})_{/S})
\]
admits
\[
    P_X^{\mathrm{dR}}
\]
as a right adjoint.  Equivalently, for every formal
$\mathfrak F_\bullet$ and every fixed-object groupoid
$\mathfrak G_\bullet$, composition with $r_{\mathfrak G}$ induces a
canonical equivalence
\[
    \Map
    \left(
        \mathfrak F_\bullet,
        P_X^{\mathrm{dR}}(\mathfrak G_\bullet)
    \right)
    \xrightarrow{\simeq}
    \Map
    \left(
        \mathfrak F_\bullet,
        \mathfrak G_\bullet
    \right).
\]
\end{theorem}

\begin{proof}
Associate to every fixed-object simplicial object
$\mathfrak G_\bullet$ its simplicial diagram of centred objects
\[
    [n]
    \longmapsto
    \left(
        X\xrightarrow{u_n}\mathfrak G_n
    \right)
    \in
    \bigl((\mathcal H_{\Sp})_{/S}\bigr)_{X/},
\]
where $u_n$ is the fully degenerate unit simplex.  Every simplicial operator
preserves the fully degenerate simplex, so this is canonically an object of
\[
    \operatorname{Fun}
    \left(
        \Delta^{\mathrm{op}},
        \bigl((\mathcal H_{\Sp})_{/S}\bigr)_{X/}
    \right).
\]
A fixed-object simplicial morphism determines and is determined by a
morphism of these centred diagrams.

Let
\[
    i:
    \bigl((\mathcal H_{\Sp})_{/S}\bigr)_{X/}^{\mathrm{for}}
    \hookrightarrow
    \bigl((\mathcal H_{\Sp})_{/S}\bigr)_{X/}
\]
be the fully faithful inclusion.  By
Theorem~\ref{thm:centred-formalization-coreflection},
\[
    i\dashv P_X^{\mathrm{dR}}.
\]
Passing pointwise to simplicial functor categories gives an adjunction
\[
    \operatorname{Fun}
    \left(
        \Delta^{\mathrm{op}},
        \bigl((\mathcal H_{\Sp})_{/S}\bigr)_{X/}^{\mathrm{for}}
    \right)
    \;\underset{\bigl(P_X^{\mathrm{dR}}\bigr)^{\Delta^{\mathrm{op}}}}{\overset{i^{\Delta^{\mathrm{op}}}}{\rightleftarrows}}\;
    \operatorname{Fun}
    \left(
        \Delta^{\mathrm{op}},
        \bigl((\mathcal H_{\Sp})_{/S}\bigr)_{X/}
    \right),
    \qquad
    i^{\Delta^{\mathrm{op}}}\dashv
    \bigl(P_X^{\mathrm{dR}}\bigr)^{\Delta^{\mathrm{op}}},
\]
whose counit is degreewise
\[
    r_{\mathfrak G,n}:
    P^{\mathrm{dR}}(u_n)
    \longrightarrow
    \mathfrak G_n.
\]

We now identify the relevant fixed points.  Suppose
$\mathfrak F_\bullet$ is formal.  By
Proposition~\ref{prop:degree-one-formality},
\[
    r_{\mathfrak F,1}:
    P^{\mathrm{dR}}(u_1)
    \xrightarrow{\simeq}
    \mathfrak F_1.
\]
For $n\ge2$, the Segal equivalences identify
$r_{\mathfrak F,n}$ with the $n$-fold fibre product over $X$ of
$r_{\mathfrak F,1}$.  Hence
\[
    r_{\mathfrak F,n}
    :
    P^{\mathrm{dR}}(u_n)
    \xrightarrow{\simeq}
    \mathfrak F_n
\]
for every $n$, while in degree zero the assertion is
$P^{\mathrm{dR}}(\operatorname{id}_X)\simeq X$.  Thus every centred degree of a
formal fixed-object groupoid lies in
$\bigl((\mathcal H_{\Sp})_{/S}\bigr)_{X/}^{\mathrm{for}}$.

Conversely, if the centred simplicial diagram is pointwise in
$\bigl((\mathcal H_{\Sp})_{/S}\bigr)_{X/}^{\mathrm{for}}$, then in particular
\[
    u_1:X\longrightarrow\mathfrak F_1
\]
is a relative de Rham equivalence, so
$\mathfrak F_\bullet$ is formal.  Therefore the formal fixed-object
groupoids are exactly the fixed-object groupoid objects whose associated
centred simplicial diagrams are pointwise fixed by $P_X^{\mathrm{dR}}$.

By Proposition~\ref{prop:unitwise-formalization-groupoid}, the pointwise right
adjoint $\bigl(P_X^{\mathrm{dR}}\bigr)^{\Delta^{\mathrm{op}}}$ carries fixed-object groupoid objects
to fixed-object groupoid objects.  It therefore restricts to
\[
    P_X^{\mathrm{dR}}:
    \operatorname{Gpd}_X((\mathcal H_{\Sp})_{/S})
    \longrightarrow
    \operatorname{FormLieGpd}(X/S).
\]
Since the groupoid categories are full subcategories of the corresponding
simplicial functor categories, the pointwise adjunction gives, for formal
$\mathfrak F_\bullet$,
\[
\begin{aligned}
    \Map_{\operatorname{Gpd}_X}
    \left(
        \mathfrak F_\bullet,
        P_X^{\mathrm{dR}}(\mathfrak G_\bullet)
    \right)
    &\simeq
    \Map_{\operatorname{Fun}
    (\Delta^{\mathrm{op}},\bigl((\mathcal H_{\Sp})_{/S}\bigr)_{X/})}
    \left(
        \mathfrak F_\bullet,
        \bigl(P_X^{\mathrm{dR}}\bigr)^{\Delta^{\mathrm{op}}}\mathfrak G_\bullet
    \right)\\
    &\simeq
    \Map_{\operatorname{Fun}
    (\Delta^{\mathrm{op}},\bigl((\mathcal H_{\Sp})_{/S}\bigr)_{X/})}
    \left(
        \mathfrak F_\bullet,
        \mathfrak G_\bullet
    \right)\\
    &\simeq
    \Map_{\operatorname{Gpd}_X}
    \left(
        \mathfrak F_\bullet,
        \mathfrak G_\bullet
    \right).
\end{aligned}
\]
The middle equivalence is the pointwise centred coreflection adjunction, and
its counit is precisely
$r_{\mathfrak G}$.  This proves the asserted right adjoint.
\end{proof}

\begin{remark}
The natural transformation points from the formalized groupoid to the original
one, so formalization is a coreflection.  The comonad and its coreflection are
not additional axioms: they are forced by the previously constructed
left-exact relative de Rham reflector and the defining pullbacks.
\end{remark}


\section{The canonical relative infinitesimal groupoid}
\label{sec:flg-canonical-relation}

Retain
\[
    \eta_{X/S}^{\mathrm{dR}}:X\to X_{\mathrm{dR}/S}
\]
and its effective-image factorization
\[
    X\xrightarrow{e_{X/S}}Q_{X/S}
      \xrightarrow{m_{X/S}}X_{\mathrm{dR}/S}.
\]

\begin{definition}[Canonical relative infinitesimal groupoid]
\label{def:canonical-relative-infinitesimal-groupoid}
Retain the Chapter~4 notation
\[
    \mathcal R^{\mathrm{dR}}_{X/S,\bullet}
    \simeq
    \check C_\bullet(e_{X/S}),
    \qquad
    \mathcal R^{\mathrm{dR}}_{X/S}:=\mathcal R^{\mathrm{dR}}_{X/S,1}.
\]
This introduces no new relation object; it only fixes the inherited notation
for use in the present chapter.
\end{definition}

\begin{lemma}[The effective infinitesimal quotient is a relative de Rham equivalence]
\label{lem:effective-quotient-dr-equivalence}
The effective epimorphism
\[
    e_{X/S}:X\twoheadrightarrow Q_{X/S}
\]
is inverted by the relative de Rham localization:
\[
    \operatorname{dR}_S(e_{X/S}):
    \operatorname{dR}_S(X)\xrightarrow{\simeq}\operatorname{dR}_S(Q_{X/S}).
\]
Moreover
\[
    \operatorname{dR}_S(m_{X/S}):
    \operatorname{dR}_S(Q_{X/S})
    \xrightarrow{\simeq}
    \operatorname{dR}_S(\operatorname{dR}_S(X))
    \simeq
    \operatorname{dR}_S(X)
\]
is an equivalence.
\end{lemma}

\begin{proof}
Apply $\operatorname{dR}_S$ to the effective-epimorphism--monomorphism factorization
\[
    X\xrightarrow{e_{X/S}}Q_{X/S}
      \xrightarrow{m_{X/S}}\operatorname{dR}_S(X)
\]
of the localization unit.  Chapter~4 proves
that relative de Rham localization carries effective epimorphisms to
effective epimorphisms, so
\[
    \operatorname{dR}_S(e_{X/S}):
    \operatorname{dR}_S(X)\longrightarrow \operatorname{dR}_S(Q_{X/S})
\]
is an effective epimorphism.  Since $\operatorname{dR}_S$ is left exact, it preserves
monomorphisms, and therefore
\[
    \operatorname{dR}_S(m_{X/S}):
    \operatorname{dR}_S(Q_{X/S})\longrightarrow \operatorname{dR}_S(\operatorname{dR}_S(X))
\]
is a monomorphism.

Their composite is
\[
    \operatorname{dR}_S(\eta_{X/S}^{\mathrm{dR}}):
    \operatorname{dR}_S(X)\longrightarrow \operatorname{dR}_S(\operatorname{dR}_S(X)),
\]
which is an equivalence by idempotence of the localization.  Hence the two
displayed morphisms form the effective-epimorphism--monomorphism
factorization of an equivalence.  By uniqueness of that factorization, both
are equivalences.
\end{proof}

Chapter~4 gives canonical equivalences
\[
    \mathcal R^{\mathrm{dR}}_{X/S,n}
    \simeq
    \underbrace{
    X\times_{Q_{X/S}}\cdots\times_{Q_{X/S}}X
    }_{n+1}
    \simeq
    \underbrace{
    X\times_{X_{\mathrm{dR}/S}}\cdots
      \times_{X_{\mathrm{dR}/S}}X
    }_{n+1}.
\]
In degree one we retain the source--target convention of Chapter~4:
\[
    \mathsf s=\operatorname{pr}_1,
    \qquad
    \mathsf t=\operatorname{pr}_2.
\]


\paragraph{Infinitesimal equality interpretation.}
In the internal notation fixed above, the degree-one relation classifies
relative de Rham infinitesimal equality:
\[
    x\sim_{\mathrm{dR}/S}y
    \qquad\Longleftrightarrow\qquad
    (x,y)\ \text{lifts to}\ \mathcal R^{\mathrm{dR}}_{X/S}.
\]
Equivalently, it is the equality relation induced by the effective
infinitesimal quotient
\[
    e_{X/S}:X\twoheadrightarrow Q_{X/S}.
\]
Thus composition, inversion, and units in
$\mathcal R^{\mathrm{dR}}_{X/S,\bullet}$ are the coherent composition, reversal, and
reflexivity operations of relative infinitesimal equality.  This is the
native de Rham realization of the infinitesimal-equality groupoid constructed
from relative infinitesimal shape in Chapter~2.

\subsection{The pair groupoid and maximal formal relation}

Define the relative pair groupoid by
\[
    \operatorname{Pair}(X/S)_n
    :=
    X^{\times_S(n+1)}.
\]
It is the relative $0$-coskeleton of $X$ and hence is terminal among
simplicial objects with specified degree-zero object $X$.

\begin{theorem}[Formalization of the pair groupoid]
\label{thm:pair-groupoid-formalization}
There is a canonical equivalence of fixed-object formal groupoids
\[
    P_X^{\mathrm{dR}}\bigl(\operatorname{Pair}(X/S)_\bullet\bigr)
    \xrightarrow{\simeq}
    \mathcal R^{\mathrm{dR}}_{X/S,\bullet}.
\]
\end{theorem}

\begin{proof}
In degree $n$, the centre of the pair groupoid is the $(n+1)$-fold diagonal
\[
    \Delta_n:X\longrightarrow X^{\times_S(n+1)}.
\]
Left exactness gives
\[
    \operatorname{dR}_S\!\left(X^{\times_S(n+1)}\right)
    \simeq
    \operatorname{dR}_S(X)^{\times_S(n+1)}.
\]
Therefore
\[
\begin{aligned}
    P^{\mathrm{dR}}(\Delta_n)
    &\simeq
    X^{\times_S(n+1)}
    \times_{\operatorname{dR}_S(X)^{\times_S(n+1)}}
    \operatorname{dR}_S(X)\\
    &\simeq
    \underbrace{
    X\times_{\operatorname{dR}_S(X)}\cdots\times_{\operatorname{dR}_S(X)}X
    }_{n+1}\\
    &\simeq
    \mathcal R^{\mathrm{dR}}_{X/S,n}.
\end{aligned}
\]
The construction uses only the product projections and diagonal, so the
equivalences commute with every face and degeneracy map.
\end{proof}

\begin{theorem}[Terminality of the maximal formal relation]
\label{thm:maximal-formal-terminal}
The groupoid $\mathcal R^{\mathrm{dR}}_{X/S,\bullet}$ is terminal in
\[
    \operatorname{FormLieGpd}(X/S).
\]
Equivalently, for every formal fixed-object groupoid $\mathfrak F_\bullet$,
the mapping space
\[
    \Map_{\operatorname{Gpd}_X}
    \left(
        \mathfrak F_\bullet,
        \mathcal R^{\mathrm{dR}}_{X/S,\bullet}
    \right)
\]
is contractible.
\end{theorem}

\begin{proof}
The relative pair groupoid
$\operatorname{Pair}(X/S)_\bullet$ is the relative $0$-coskeleton of $X$ and
is therefore terminal in
$\operatorname{Gpd}_X((\mathcal H_{\Sp})_{/S})$.
Theorem~\ref{thm:formalization-coreflection} identifies unitwise de Rham formalization
as the right adjoint to the inclusion of formal fixed-object groupoids.
A right adjoint carries terminal objects to terminal objects.  Hence
\[
    P_X^{\mathrm{dR}}\bigl(\operatorname{Pair}(X/S)_\bullet\bigr)
\]
is terminal in
$\operatorname{FormLieGpd}(X/S)$.
By Theorem~\ref{thm:pair-groupoid-formalization} this formalized pair groupoid is
canonically equivalent to $\mathcal R^{\mathrm{dR}}_{X/S,\bullet}$.
\end{proof}


\begin{corollary}[Canonical infinitesimal anchor of a formal groupoid]
\label{cor:canonical-formal-anchor}
Every formal fixed-object groupoid
$\mathfrak F_\bullet\in
\operatorname{FormLieGpd}(X/S)$
admits a contractibly unique morphism
\[
    \phi_{\mathfrak F}:
    \mathfrak F_\bullet
    \longrightarrow
    \mathcal R^{\mathrm{dR}}_{X/S,\bullet}.
\]
We call it the \emph{canonical infinitesimal anchor}.
\end{corollary}

\begin{proof}
This is Theorem~\ref{thm:maximal-formal-terminal}.
\end{proof}

\begin{proposition}[The canonical anchor is the formalized source--target map]
\label{prop:canonical-anchor-formalized-source-target}
Let
\[
    \mathfrak G_\bullet
    \in
    \operatorname{FormLieGpd}(X/S).
\]
In degree one, the canonical anchor of
Corollary~\ref{cor:canonical-formal-anchor} is canonically equivalent to the
formalization of the source--target map:
\[
\begin{tikzcd}
P^{\mathrm{dR}}(u)
    \arrow[r,"{P_X^{\mathrm{dR}}(s,t)}"]
    \arrow[d,"{r_{\mathfrak G,1}}"',"\simeq"]
&
P^{\mathrm{dR}}(\Delta_X)
    \arrow[d,"\simeq"]
\\
\mathfrak G_1
    \arrow[r,"{\phi_{\mathfrak G,1}}"']
&
\mathcal R^{\mathrm{dR}}_{X/S}.
\end{tikzcd}
\]
Thus, after the canonical identifications supplied by formality and
Theorem~\ref{thm:pair-groupoid-formalization},
\[
    \phi_{\mathfrak G,1}
    \simeq
    P_X^{\mathrm{dR}}(s,t).
\]
\end{proposition}

\begin{proof}
The canonical vertex morphism
\[
    \mathfrak G_\bullet
    \longrightarrow
    \operatorname{Pair}(X/S)_\bullet
\]
has degree-one component $(s,t):\mathfrak G_1\to X\times_SX$ and carries the
unit centre to the diagonal.  Applying unitwise de Rham formalization gives
\[
    P_X^{\mathrm{dR}}(\mathfrak G_\bullet)
    \longrightarrow
    P_X^{\mathrm{dR}}\bigl(\operatorname{Pair}(X/S)_\bullet\bigr).
\]
The source identifies with $\mathfrak G_\bullet$ because $\mathfrak G$ is
formal, and the target identifies with $\mathcal R^{\mathrm{dR}}_{X/S,\bullet}$ by
Theorem~\ref{thm:pair-groupoid-formalization}.  This is the canonical anchor
constructed in Theorem~\ref{thm:maximal-formal-terminal}; its degree-one
component is therefore the displayed formalized source--target map.
\end{proof}

\begin{proposition}[Extremal formal groupoids]
\label{prop:extremal-infinitesimal-groupoids}
There are canonical formal groupoids on $X$:
\begin{enumerate}
    \item the identity or zero groupoid
    \[
        0_{X,\bullet}
        :=
        \check C_\bullet(\operatorname{id}_X),
    \]
    whose quotient is $X$;
    \item the maximal infinitesimal groupoid
    $\mathcal R^{\mathrm{dR}}_{X/S,\bullet}$, whose quotient is $Q_{X/S}$.
\end{enumerate}
For every formal fixed-object groupoid $\mathfrak G_\bullet$ there are
canonical morphisms
\[
    0_{X,\bullet}
    \longrightarrow
    \mathfrak G_\bullet
    \longrightarrow
    \mathcal R^{\mathrm{dR}}_{X/S,\bullet}.
\]
The first is initial and the second terminal in the formal fixed-object
category.
\end{proposition}

\begin{proof}
The zero groupoid is formal because its degree-one centre is
$\operatorname{id}_X$.  Under
Proposition~\ref{prop:groupoids-effective-quotients-equivalence} it
corresponds to
\[
    \operatorname{id}_X:X\longrightarrow X,
\]
which is initial in the undercategory $\bigl((\mathcal H_{\Sp})_{/S}\bigr)_{X/}$ and hence initial in
its full effective-epimorphism subcategory $\operatorname{EffEpi}(X)$.  Therefore
$0_{X,\bullet}$ is initial in $\operatorname{Gpd}_X((\mathcal H_{\Sp})_{/S})$ and in its
full formal subcategory.  The maximal relation is formal by
Theorem~\ref{thm:pair-groupoid-formalization} and idempotence of formalization,
and is terminal by Theorem~\ref{thm:maximal-formal-terminal}.
\end{proof}

\begin{remark}
The word ``zero'' refers to the absence of additional relation arrows between
distinct objects.  Because the ambient category is an $\infty$-topos, this
terminology must not be read as a claim that all homotopy-theoretic inertia of
the object $X$ itself vanishes.  The raw and formal isotropy objects below
make this distinction explicit.
\end{remark}


\section{Formal Lie groupoids and their quotient classification}
\label{sec:flg-definition-classification}

\subsection{Anchored and formal groupoids}

\begin{definition}[Infinitesimally anchored groupoid]
\label{def:infinitesimally-anchored-groupoid}
An \emph{infinitesimally anchored groupoid on $X/S$} is an object
\[
    \mathfrak G_\bullet
    \in
    \operatorname{Gpd}_X((\mathcal H_{\Sp})_{/S})
\]
equipped with a morphism
\[
    \phi_{\mathfrak G}:
    \mathfrak G_\bullet
    \longrightarrow
    \mathcal R^{\mathrm{dR}}_{X/S,\bullet}.
\]
The $\infty$-category of anchored groupoids is
\[
    \operatorname{AnchGpd}(X/S)
    :=
    \left(
        \operatorname{Gpd}_X((\mathcal H_{\Sp})_{/S})
    \right)_{/\mathcal R^{\mathrm{dR}}_{X/S,\bullet}}.
\]
\end{definition}

\begin{definition}[Brave-new formal Lie groupoid]
\label{def:formal-lie-groupoid}
A \emph{brave-new formal Lie groupoid on $X/S$} is an object
\[
    \mathfrak G_\bullet\in\operatorname{FormLieGpd}(X/S).
\]
Equivalently, it is a fixed-object groupoid whose unit
\[
    u:X\longrightarrow\mathfrak G_1
\]
is a relative de Rham equivalence, or whose unitwise formalization counit
\[
    P_X^{\mathrm{dR}}(\mathfrak G_\bullet)
    \longrightarrow
    \mathfrak G_\bullet
\]
is an equivalence by Proposition~\ref{prop:degree-one-formality}.  No
infinitesimal anchor is part of this definition.
\end{definition}

\begin{theorem}[The anchor is constructed, not additional formal data]
\label{thm:formal-anchor-forgetful-equivalence}
The canonical anchor of
Corollary~\ref{cor:canonical-formal-anchor} defines a fully faithful functor
\[
    \operatorname{can}_{X/S}:
    \operatorname{FormLieGpd}(X/S)
    \longrightarrow
    \operatorname{AnchGpd}(X/S).
\]
Its essential image is the full subcategory of anchored groupoids whose
underlying fixed-object groupoid is formal.  Forgetting the anchor on this
essential image is inverse to $\operatorname{can}_{X/S}$.
\end{theorem}

\begin{proof}
For every formal fixed-object groupoid $\mathfrak F_\bullet$,
Theorem~\ref{thm:maximal-formal-terminal} gives
\[
    \Map_{\operatorname{Gpd}_X}
    \left(
        \mathfrak F_\bullet,
        \mathcal R^{\mathrm{dR}}_{X/S,\bullet}
    \right)
    \simeq *.
\]
Hence the space of anchors on $\mathfrak F_\bullet$ is contractible.

Let $\mathfrak F_\bullet,\mathfrak G_\bullet$ be formal.  Any fixed-object
groupoid morphism
\[
    F:\mathfrak F_\bullet\longrightarrow\mathfrak G_\bullet
\]
gives two maps
\[
    \phi_{\mathfrak G}F,\ \phi_{\mathfrak F}:
    \mathfrak F_\bullet\rightrightarrows \mathcal R^{\mathrm{dR}}_{X/S,\bullet}.
\]
The mapping space from $\mathfrak F_\bullet$ to $\mathcal R^{\mathrm{dR}}_{X/S,\bullet}$ is
contractible, so the space of compatible homotopies
\[
    \phi_{\mathfrak G}F\simeq\phi_{\mathfrak F}
\]
is contractible.  Therefore passing from fixed-object morphisms to anchored
morphisms changes no mapping space.  This proves full faithfulness and shows
that every anchored formal object lies, through a contractible space of
identifications, in the essential image.
\end{proof}

\begin{remark}[No monomorphism condition]
\label{rem:no-mono-formal-gpd}
The canonical anchor
\[
    \mathfrak G_1\longrightarrow \mathcal R^{\mathrm{dR}}_{X/S}
\]
is not required to be a monomorphism.  Its homotopy fibres carry genuine
isotropy information.  Requiring a monomorphism would erase those fibres and,
in particular, would exclude one-object formal groups and formal gauge
groupoids.  The formality condition is instead that the unit section be a relative de Rham
equivalence.
\end{remark}

\subsection{Canonical formalization of arbitrary internal groupoids}

\begin{theorem}[Canonical Unitwise Formalization Theorem]
\label{thm:canonical-unitwise-formalization}
Let
\[
    \mathcal H_\bullet\rightrightarrows X
\]
be any groupoid object in $(\mathcal H_{\Sp})_{/S}$.  Its unitwise de Rham formalization
\[
    P_X^{\mathrm{dR}}(\mathcal H_\bullet)
\]
is formal and admits a canonical anchor
\[
    \phi_{P_X^{\mathrm{dR}}\mathcal H}:
    P_X^{\mathrm{dR}}(\mathcal H_\bullet)
    \longrightarrow
    \mathcal R^{\mathrm{dR}}_{X/S,\bullet}.
\]
The construction is functorial in fixed-object groupoid morphisms and
compatible with arbitrary base change in $S$.
\end{theorem}

\begin{proof}
There is a contractibly unique fixed-object simplicial morphism
\[
    \mathcal H_\bullet
    \longrightarrow
    \operatorname{Pair}(X/S)_\bullet
\]
given degreewise by the ordered tuple of vertex maps.  Applying unitwise
formalization and Theorem~\ref{thm:pair-groupoid-formalization} gives
\[
    P_X^{\mathrm{dR}}(\mathcal H_\bullet)
    \longrightarrow
    P_X^{\mathrm{dR}}\bigl(\operatorname{Pair}(X/S)_\bullet\bigr)
    \simeq
    \mathcal R^{\mathrm{dR}}_{X/S,\bullet}.
\]
The source is formal by
Theorem~\ref{thm:unitwise-formalization-idempotent}.  Functoriality follows from
functoriality of the map to the relative pair groupoid and of formalization.
Base-change compatibility follows from
Proposition~\ref{prop:centred-formalization-base-change}.
\end{proof}

\begin{proposition}[Coreflection of anchored groupoids]
\label{prop:anchored-formal-coreflection}
The canonical-anchor functor
\[
    \operatorname{can}_{X/S}:
    \operatorname{FormLieGpd}(X/S)
    \longrightarrow
    \operatorname{AnchGpd}(X/S)
\]
admits unitwise de Rham formalization as a right adjoint
\[
    P_X^{\mathrm{dR}}:
    \operatorname{AnchGpd}(X/S)
    \longrightarrow
    \operatorname{FormLieGpd}(X/S).
\]
For an anchored groupoid
\[
    \phi_{\mathfrak G}:
    \mathfrak G_\bullet\to \mathcal R^{\mathrm{dR}}_{X/S,\bullet},
\]
the canonical anchor of
$P_X^{\mathrm{dR}}(\mathfrak G_\bullet)$ is equivalently represented by
either composite
\[
    P_X^{\mathrm{dR}}(\mathfrak G_\bullet)
    \xrightarrow{r_{\mathfrak G}}
    \mathfrak G_\bullet
    \xrightarrow{\phi_{\mathfrak G}}
    \mathcal R^{\mathrm{dR}}_{X/S,\bullet}
\]
or
\[
    P_X^{\mathrm{dR}}(\mathfrak G_\bullet)
    \xrightarrow{P_X^{\mathrm{dR}}(\phi_{\mathfrak G})}
    P_X^{\mathrm{dR}}\bigl(\mathcal R^{\mathrm{dR}}_{X/S,\bullet}\bigr)
    \xrightarrow{r_{\mathcal R}}
    \mathcal R^{\mathrm{dR}}_{X/S,\bullet}.
\]
\end{proposition}

\begin{proof}
Naturality of the formalization counit gives
\[
    r_{\mathcal R}
    P_X^{\mathrm{dR}}(\phi_{\mathfrak G})
    \simeq
    \phi_{\mathfrak G}r_{\mathfrak G}.
\]
Because $\mathcal R^{\mathrm{dR}}_{X/S,\bullet}$ is formal, $r_{\mathcal R}$ is an equivalence, so the two
displayed anchors agree canonically.  The formalized groupoid is formal by
Theorem~\ref{thm:unitwise-formalization-idempotent}; its anchor is therefore the
canonical one by Theorem~\ref{thm:maximal-formal-terminal}.

Now let $\mathfrak F_\bullet$ be formal and let
$(\mathfrak G_\bullet,\phi_{\mathfrak G})$ be anchored.  Since
\[
    \Map_{\operatorname{Gpd}_X}
    \left(
        \mathfrak F_\bullet,
        \mathcal R^{\mathrm{dR}}_{X/S,\bullet}
    \right)
    \simeq *,
\]
the anchor-compatibility condition on a morphism
$\mathfrak F_\bullet\to\mathfrak G_\bullet$ has a contractible space of
solutions.  Hence
\[
\begin{aligned}
    \Map_{\operatorname{AnchGpd}}
    \left(
        \operatorname{can}_{X/S}(\mathfrak F_\bullet),
        (\mathfrak G_\bullet,\phi_{\mathfrak G})
    \right)
    &\simeq
    \Map_{\operatorname{Gpd}_X}
    \left(
        \mathfrak F_\bullet,
        \mathfrak G_\bullet
    \right)\\
    &\simeq
    \Map_{\operatorname{Gpd}_X}
    \left(
        \mathfrak F_\bullet,
        P_X^{\mathrm{dR}}(\mathfrak G_\bullet)
    \right)\\
    &\simeq
    \Map_{\operatorname{FormLieGpd}(X/S)}
    \left(
        \mathfrak F_\bullet,
        P_X^{\mathrm{dR}}(\mathfrak G_\bullet)
    \right),
\end{aligned}
\]
where the middle equivalence is
Theorem~\ref{thm:formalization-coreflection}.  This is the claimed adjunction.
\end{proof}

\subsection{Effective formal leaf spaces}

\begin{definition}[Effective formal leaf space]
\label{def:formal-leaf-space}
For
$\mathfrak G_\bullet\in\operatorname{FormLieGpd}(X/S)$ define
\[
    Q_{\mathfrak G}
    :=
    \lvert\mathfrak G_\bullet\rvert.
\]
Its quotient presentation is
\[
    q_{\mathfrak G}:
    X\twoheadrightarrow Q_{\mathfrak G}.
\]
Realizing the canonical anchor gives
\[
    \bar\phi_{\mathfrak G}:
    Q_{\mathfrak G}
    \longrightarrow
    Q_{X/S}
\]
with
\[
    \bar\phi_{\mathfrak G}q_{\mathfrak G}
    \simeq
    e_{X/S}.
\]
We call $Q_{\mathfrak G}$ the \emph{effective formal leaf space} or
\emph{effective formal quotient}.
\end{definition}

\subsection{Anchored quotient classification}

\begin{definition}[Infinitesimal quotient factorizations]
\label{def:infinitesimal-quotient-factorizations}
Let
\[
    \iota_{02}:\Delta^1\longrightarrow\Delta^2,
    \qquad
    0\longmapsto0,\quad 1\longmapsto2,
\]
be the long-edge inclusion.  Define the $\infty$-category of factorizations of
$e_{X/S}$ by the homotopy fibre
\[
    \operatorname{Fact}(e_{X/S})
    :=
    \operatorname{fib}_{e_{X/S}}
    \left(
        \iota_{02}^*:
        \operatorname{Fun}(\Delta^2,(\mathcal H_{\Sp})_{/S})
        \longrightarrow
        \operatorname{Fun}(\Delta^1,(\mathcal H_{\Sp})_{/S})
    \right).
\]
Thus its objects are coherently commutative factorizations
\[
    X\xrightarrow{q}Q\xrightarrow{r}Q_{X/S},
    \qquad
    rq\simeq e_{X/S},
\]
and all mapping-space coherence is inherited from this homotopy fibre.

Let
\[
    \operatorname{Fact}_{\mathrm{eff}}(e_{X/S})
    \subset
    \operatorname{Fact}(e_{X/S})
\]
be the full subcategory on the factorizations for which
$q:X\twoheadrightarrow Q$ is an effective epimorphism.
\end{definition}

\begin{lemma}[Effectivity of the second quotient map]
\label{lem:second-factor-effective}
Let
\[
    A\xrightarrow{f}B\xrightarrow{g}C
\]
be morphisms in an $\infty$-topos.  If $gf$ is an effective epimorphism, then
$g$ is an effective epimorphism.  Consequently every object of
$\operatorname{Fact}_{\mathrm{eff}}(e_{X/S})$ is canonically a chain
\[
    X\twoheadrightarrow Q\twoheadrightarrow Q_{X/S}
\]
of effective epimorphisms.
\end{lemma}

\begin{proof}
Factor $g$ by the effective-epimorphism--monomorphism factorization
\[
    B\twoheadrightarrow I_g\lhook\joinrel\longrightarrow C.
\]
The composite $gf$ factors through the monomorphism $I_g\to C$.  Since the
effective image of the effective epimorphism $gf$ is all of $C$, this
monomorphism is an equivalence.  Hence $g$ is an effective epimorphism.
Apply this to
\[
    X\xrightarrow{q}Q\xrightarrow{r}Q_{X/S},
    \qquad
    rq\simeq e_{X/S}.
\]
\end{proof}

\begin{theorem}[Anchored Quotient Classification Theorem]
\label{thm:anchored-gpd-quotient-classification}
Realization and \v{C}ech nerve induce inverse equivalences
\[
    \operatorname{AnchGpd}(X/S)
    \simeq
    \operatorname{Fact}_{\mathrm{eff}}(e_{X/S}).
\]
Under this equivalence,
\[
    \mathfrak G_\bullet
    \longleftrightarrow
    \left(
        X\twoheadrightarrow Q_{\mathfrak G}
        \twoheadrightarrow Q_{X/S}
    \right),
\]
and
\[
    \mathfrak G_n
    \simeq
    \underbrace{
    X\times_{Q_{\mathfrak G}}\cdots
    \times_{Q_{\mathfrak G}}X
    }_{n+1}.
\]
\end{theorem}

\begin{proof}
Proposition~\ref{prop:groupoids-effective-quotients-equivalence} is an
equivalence of $\infty$-categories
\[
    \operatorname{Gpd}_X((\mathcal H_{\Sp})_{/S})
    \simeq
    \operatorname{EffEpi}(X).
\]
Under it the canonical relation
$\mathcal R^{\mathrm{dR}}_{X/S,\bullet}$ corresponds to the effective epimorphism
\[
    e_{X/S}:X\twoheadrightarrow Q_{X/S}.
\]
Passing to slices therefore gives
\[
\begin{aligned}
    \operatorname{AnchGpd}(X/S)
    &=
    \left(
        \operatorname{Gpd}_X((\mathcal H_{\Sp})_{/S})
    \right)_{/\mathcal R^{\mathrm{dR}}_{X/S,\bullet}}\\
    &\simeq
    \operatorname{EffEpi}(X)_{/e_{X/S}}.
\end{aligned}
\]
An object of the latter slice is exactly an effective epimorphism
$q:X\twoheadrightarrow Q$ together with a morphism under $X$
\[
    r:(X\xrightarrow{q}Q)
    \longrightarrow
    (X\xrightarrow{e_{X/S}}Q_{X/S}),
\]
that is, a coherently commutative factorization
\[
    X\xrightarrow{q}Q\xrightarrow{r}Q_{X/S},
    \qquad
    rq\simeq e_{X/S}.
\]
Thus the slice is canonically
$\operatorname{Fact}_{\mathrm{eff}}(e_{X/S})$ as defined above.  By
Lemma~\ref{lem:second-factor-effective}, the second map $r$ is automatically
an effective epimorphism.  Because the argument is obtained by slicing an
equivalence of $\infty$-categories, all mapping spaces and higher coherence
are already included.  The degreewise formula is the \v{C}ech description
from Proposition~\ref{prop:groupoids-effective-quotients-equivalence}.
\end{proof}

\subsection{Quotient-side formalization and recognition}

\begin{definition}[Centred quotient formalization]
\label{def:quotient-formalization}
Let
\[
    q:X\twoheadrightarrow Q
\]
be an effective epimorphism.  Retain the Chapter~4 de Rham--formally
\'etale factorization notation
\[
    X
    \xrightarrow{\ell_q}
    P^{\mathrm{dR}}(q)
    \xrightarrow{r_q}
    Q,
    \qquad
    P^{\mathrm{dR}}(q)
    =
    Q\times_{Q_{\mathrm{dR}/S}}X_{\mathrm{dR}/S},
\]
where
\[
    \ell_q=(q,\eta_{X/S}^{\mathrm{dR}})
\]
and $r_q$ is the first projection.  Factor $\ell_q$ by the
effective-epimorphism--monomorphism factorization:
\[
    X
    \xrightarrow{e_{\ell_q}}
    Q_{\ell_q}
    \xrightarrow{m_{\ell_q}}
    P^{\mathrm{dR}}(q).
\]
Put
\[
    c_q
    :=
    r_qm_{\ell_q}:
    Q_{\ell_q}\longrightarrow Q.
\]
Then
\[
    c_qe_{\ell_q}\simeq q.
\]
\end{definition}

\begin{lemma}[The quotient comparison is effective]
\label{lem:quotient-formalization-effective}
The morphism
\[
    c_q:Q_{\ell_q}\longrightarrow Q
\]
is an effective epimorphism.
\end{lemma}

\begin{proof}
The composite
\[
    X\xrightarrow{e_{\ell_q}}Q_{\ell_q}
      \xrightarrow{c_q}Q
\]
is the effective epimorphism $q$.  The image of the composite factors through
the image of $c_q$, so the image of $c_q$ is all of $Q$.
\end{proof}

\begin{theorem}[Quotient formula for unitwise formalization]
\label{thm:quotient-formula-unitwise-formalization}
For an effective epimorphism $q:X\twoheadrightarrow Q$ there is a canonical
equivalence
\[
    P_X^{\mathrm{dR}}\bigl(\check C_\bullet(q)\bigr)
    \simeq
    \check C_\bullet(e_{\ell_q}).
\]
Under this equivalence the realization of the formalization counit is
\[
    c_q:Q_{\ell_q}\longrightarrow Q.
\]
\end{theorem}

\begin{proof}
In degree $n$,
\[
    \check C_n(q)
    =
    \underbrace{X\times_Q\cdots\times_QX}_{n+1}
\]
with diagonal centre
\[
    \Delta_n^q:X\longrightarrow\check C_n(q).
\]
Left exactness gives
\[
    \operatorname{dR}_S(\check C_n(q))
    \simeq
    \underbrace{
    \operatorname{dR}_S(X)\times_{\operatorname{dR}_S(Q)}\cdots\times_{\operatorname{dR}_S(Q)}\operatorname{dR}_S(X)
    }_{n+1},
\]
and the localized centre is the full diagonal from $\operatorname{dR}_S(X)$.  Substituting
this into the defining pullback for centred formal \'etalification and regrouping finite
limits gives
\[
\begin{aligned}
    P^{\mathrm{dR}}(\Delta_n^q)
    &\simeq
    \underbrace{
    X\times_{P^{\mathrm{dR}}(q)}\cdots\times_{P^{\mathrm{dR}}(q)}X
    }_{n+1},
\end{aligned}
\]
where
\[
    P^{\mathrm{dR}}(q)
    =
    Q\times_{\operatorname{dR}_S(Q)}\operatorname{dR}_S(X)
\]
and each map $X\to P^{\mathrm{dR}}(q)$ is the de Rham-equivalence centre $\ell_q$.

Because $m_{\ell_q}:Q_{\ell_q}\to P^{\mathrm{dR}}(q)$ is a monomorphism,
\[
    Q_{\ell_q}
    \simeq
    Q_{\ell_q}\times_{P^{\mathrm{dR}}(q)}Q_{\ell_q}.
\]
Since $\ell_q=m_{\ell_q}e_{\ell_q}$, finite-limit pasting therefore gives
\[
    \underbrace{X\times_{P^{\mathrm{dR}}(q)}\cdots\times_{P^{\mathrm{dR}}(q)}X}_{n+1}
    \simeq
    \underbrace{
    X\times_{Q_{\ell_q}}\cdots
    \times_{Q_{\ell_q}}X
    }_{n+1}.
\]
The equivalences are compatible with deletion and repetition of factors, so
they assemble simplicially.  Since $e_{\ell_q}$ is effective, the right-hand
\v{C}ech nerve realizes to $Q_{\ell_q}$.  The counit is induced by
$c_qe_{\ell_q}\simeq q$, hence realizes to $c_q$.
\end{proof}

\begin{theorem}[Recognition of formal quotients]
\label{thm:formal-quotient-monomorphism-recognition}
For an effective epimorphism
\[
    q:X\twoheadrightarrow Q,
\]
the following are equivalent:
\begin{enumerate}
    \item $\check C_\bullet(q)$ is a formal fixed-object groupoid;
    \item
    \[
        \operatorname{dR}_S(q):
        X_{\mathrm{dR}/S}
        \longrightarrow
        Q_{\mathrm{dR}/S}
    \]
    is a monomorphism;
    \item $\operatorname{dR}_S(q)$ is an equivalence;
    \item
    \[
        c_q:Q_{\ell_q}\longrightarrow Q
    \]
    is an equivalence.
\end{enumerate}
\end{theorem}

\begin{proof}
The localized degree-one unit of $\check C_\bullet(q)$ is the diagonal
\[
    \Delta_{\operatorname{dR}_S(q)}:
    X_{\mathrm{dR}/S}
    \longrightarrow
    X_{\mathrm{dR}/S}
    \times_{Q_{\mathrm{dR}/S}}
    X_{\mathrm{dR}/S}.
\]
By Proposition~\ref{prop:degree-one-formality}, the \v{C}ech groupoid is formal
if and only if this diagonal is an equivalence.  A morphism in an
$\infty$-topos is a monomorphism precisely when its diagonal is an
equivalence.  This proves the equivalence of (1) and (2).

Chapter~4 proves that relative de Rham
localization carries effective epimorphisms in the slice to effective
epimorphisms in the ambient slice; this conclusion is not being inferred from
left exactness alone.  Since $q$ is an effective epimorphism,
$\operatorname{dR}_S(q)$ is therefore an effective epimorphism.  In an
$\infty$-topos a morphism which is both an effective epimorphism and a
monomorphism is an equivalence.  Hence (2) and (3) are equivalent.

By Theorem~\ref{thm:quotient-formula-unitwise-formalization}, the formalization
counit realizes to $c_q$.  A morphism of effective groupoid presentations is
an equivalence if and only if its realization is an equivalence, proving the
equivalence of (1) and (4).
\end{proof}

\begin{definition}[Formal effective quotients]
\label{def:formal-effective-quotients}
Let
\[
    \operatorname{EffEpi}^{\mathrm{for}}(X)
    \subset
    \operatorname{EffEpi}(X)
\]
be the full subcategory spanned by effective epimorphisms
$q:X\twoheadrightarrow Q$ for which
$\operatorname{dR}_S(q)$ is an equivalence.  By
Theorem~\ref{thm:formal-quotient-monomorphism-recognition}, this is
equivalently the condition that $\operatorname{dR}_S(q)$ be a monomorphism.
\end{definition}

\begin{theorem}[Quotient formalization coreflection]
\label{thm:quotient-formalization-coreflection}
The inclusion
\[
    \operatorname{EffEpi}^{\mathrm{for}}(X)
    \hookrightarrow
    \operatorname{EffEpi}(X)
\]
admits
\[
    q\longmapsto e_{\ell_q}
\]
as a right adjoint.  Its counit at $q$ is
\[
    c_q:Q_{\ell_q}\to Q.
\]
\end{theorem}

\begin{proof}
Transport the groupoid formalization coreflection
of Theorem~\ref{thm:formalization-coreflection} through the equivalence
\[
    \operatorname{Gpd}_X((\mathcal H_{\Sp})_{/S})
    \simeq
    \operatorname{EffEpi}(X)
\]
of Proposition~\ref{prop:groupoids-effective-quotients-equivalence}.
Theorem~\ref{thm:quotient-formula-unitwise-formalization} identifies the image
of $q$ under the transported right adjoint with $e_{\ell_q}$, and identifies
the transported counit with $c_q$.  The fixed points are
$\operatorname{EffEpi}^{\mathrm{for}}(X)$ by
Theorem~\ref{thm:formal-quotient-monomorphism-recognition}.
\end{proof}

\begin{proposition}[Canonical factorization through the maximal quotient]
\label{prop:formal-quotient-canonical-anchor}
Every formal effective quotient
\[
    q:X\twoheadrightarrow Q
\]
admits a contractibly unique morphism
\[
    \bar\phi_q:Q\longrightarrow Q_{X/S}
\]
such that
\[
    \bar\phi_q q\simeq e_{X/S}.
\]
The morphism $\bar\phi_q$ is an effective epimorphism and a relative de Rham
equivalence:
\[
    \operatorname{dR}_S(\bar\phi_q):
    \operatorname{dR}_S(Q)\xrightarrow{\simeq}\operatorname{dR}_S(Q_{X/S}).
\]
\end{proposition}

\begin{proof}
The \v{C}ech groupoid $\check C_\bullet(q)$ is formal, so it has a
contractibly unique canonical anchor
\[
    \check C_\bullet(q)\longrightarrow \mathcal R^{\mathrm{dR}}_{X/S,\bullet}.
\]
Realization gives $\bar\phi_q$, and the degree-zero identity gives
$\bar\phi_q q\simeq e_{X/S}$.  Conversely, a map $\bar\phi_q$ with this property
determines the corresponding anchored \v{C}ech morphism, so uniqueness of the
anchor gives contractible uniqueness of $\bar\phi_q$.

The composite $\bar\phi_q q=e_{X/S}$ is effective, hence
Lemma~\ref{lem:second-factor-effective} makes $\bar\phi_q$ effective.

By Theorem~\ref{thm:formal-quotient-monomorphism-recognition},
formality of $\check C_\bullet(q)$ gives an equivalence
\[
    \operatorname{dR}_S(q):\operatorname{dR}_S(X)\xrightarrow{\simeq}\operatorname{dR}_S(Q).
\]
Lemma~\ref{lem:effective-quotient-dr-equivalence} gives independently
\[
    \operatorname{dR}_S(e_{X/S}):
    \operatorname{dR}_S(X)\xrightarrow{\simeq}\operatorname{dR}_S(Q_{X/S}).
\]
Applying $\operatorname{dR}_S$ to $\bar\phi_q q\simeq e_{X/S}$ yields
\[
    \operatorname{dR}_S(\bar\phi_q)\operatorname{dR}_S(q)\simeq \operatorname{dR}_S(e_{X/S}).
\]
Thus
\[
    \operatorname{dR}_S(\bar\phi_q)
    \simeq
    \operatorname{dR}_S(e_{X/S})\operatorname{dR}_S(q)^{-1}
\]
is an equivalence.
\end{proof}

\begin{definition}[Formal infinitesimal quotient factorizations]
\label{def:formal-infinitesimal-quotient-factorizations}
Let
\[
    \operatorname{Fact}^{\mathrm{for}}_{\mathrm{eff}}
    (e_{X/S})
    \subset
    \operatorname{Fact}_{\mathrm{eff}}(e_{X/S})
\]
be the full subcategory on factorizations
\[
    X\xrightarrow{q}Q\xrightarrow{r}Q_{X/S}
\]
for which $\operatorname{dR}_S(q)$ is an equivalence.  Equivalently, by
Theorem~\ref{thm:formal-quotient-monomorphism-recognition},
$\operatorname{dR}_S(q)$ is a monomorphism, or
$c_q:Q_{\ell_q}\to Q$ is an equivalence.
For every such factorization both arrows are effective epimorphisms by
Lemma~\ref{lem:second-factor-effective}, and
Proposition~\ref{prop:formal-quotient-canonical-anchor} shows that the second
arrow is also a relative de Rham equivalence.
\end{definition}

\begin{theorem}[Formal Quotient Classification Theorem]
\label{thm:formal-gpd-quotient-classification}
There are canonical equivalences
\[
    \operatorname{FormLieGpd}(X/S)
    \simeq
    \operatorname{EffEpi}^{\mathrm{for}}(X)
    \simeq
    \operatorname{Fact}^{\mathrm{for}}_{\mathrm{eff}}
    (e_{X/S}).
\]
Thus a brave-new formal Lie groupoid is equivalently an effective quotient
\[
    q:X\twoheadrightarrow Q
\]
such that
\[
    \operatorname{dR}_S(q):
    X_{\mathrm{dR}/S}\to Q_{\mathrm{dR}/S}
\]
is monic, equivalently an equivalence.  The quotient anchor
\[
    \bar\phi_q:Q\longrightarrow Q_{X/S}
\]
is then recovered canonically and is both an effective epimorphism and a
relative de Rham equivalence.
\end{theorem}

\begin{proof}
By Definition~\ref{def:unitwise-formal-groupoid},
$\operatorname{FormLieGpd}(X/S)$ is the full subcategory of formal
fixed-object groupoids.  The first equivalence is therefore obtained by restricting
\[
    \operatorname{Gpd}_X((\mathcal H_{\Sp})_{/S})
    \simeq
    \operatorname{EffEpi}(X)
\]
of Proposition~\ref{prop:groupoids-effective-quotients-equivalence} to the
formal objects, using
Theorem~\ref{thm:formal-quotient-monomorphism-recognition}.

By Theorem~\ref{thm:anchored-gpd-quotient-classification}, restricting the
anchored quotient equivalence to formal underlying groupoids identifies
\[
    \operatorname{Fact}^{\mathrm{for}}_{\mathrm{eff}}(e_{X/S})
    \simeq
    \operatorname{EffEpi}^{\mathrm{for}}(X)_{/e_{X/S}}.
\]
Under the formal groupoid/effective quotient equivalence,
$e_{X/S}:X\twoheadrightarrow Q_{X/S}$ corresponds to
$\mathcal R^{\mathrm{dR}}_{X/S,\bullet}$, which is terminal among formal fixed-object
groupoids by Theorem~\ref{thm:maximal-formal-terminal}.  Hence $e_{X/S}$ is
terminal in $\operatorname{EffEpi}^{\mathrm{for}}(X)$.  Slicing an
$\infty$-category over a terminal object gives an equivalence with the
original $\infty$-category, so
\[
    \operatorname{EffEpi}^{\mathrm{for}}(X)_{/e_{X/S}}
    \simeq
    \operatorname{EffEpi}^{\mathrm{for}}(X).
\]
This proves the final equivalence including all mapping-space coherence.
The assertions about the recovered map $Q\to Q_{X/S}$ are
Proposition~\ref{prop:formal-quotient-canonical-anchor}.
\end{proof}

\begin{corollary}[Initial and terminal formal Lie groupoids]
\label{cor:initial-terminal-formal-gpd}
The category $\operatorname{FormLieGpd}(X/S)$ has canonical initial and
terminal objects
\[
    0_{X,\bullet}
    \longrightarrow
    \mathfrak G_\bullet
    \longrightarrow
    \mathcal R^{\mathrm{dR}}_{X/S,\bullet}.
\]
Their quotient factorizations are
\[
    X\xrightarrow{\operatorname{id}}X
      \xrightarrow{e_{X/S}}Q_{X/S}
\]
and
\[
    X\xrightarrow{e_{X/S}}Q_{X/S}
      \xrightarrow{\operatorname{id}}Q_{X/S}.
\]
\end{corollary}

\begin{proof}
Apply Proposition~\ref{prop:extremal-infinitesimal-groupoids} and
Definition~\ref{def:formal-lie-groupoid}.
\end{proof}

\begin{remark}[Nonlinear anchor]
\label{rem:nonlinear-anchor}
Before formalization, an anchor is genuine additional structure.  After
formalization it is the contractibly unique morphism to the maximal formal
relation.  On quotients it becomes
\[
    \bar\phi_{\mathfrak G}:Q_{\mathfrak G}\to Q_{X/S}.
\]
The final handoff section below records the represented first-order
linearization of the canonical nonlinear anchor in the geometric range where
the cotangent theory of Chapter~2 applies.  No bracket,
Lie-algebraic presentation, or recognition theorem is introduced in the
intrinsic formal-groupoid interface.
\end{remark}


\section{Formal orbit disks and formal leaves}
\label{sec:flg-formal-disks}

Fix
\[
    \mathfrak G_\bullet
    \in
    \operatorname{FormLieGpd}(X/S)
\]
and write
\[
    q:=q_{\mathfrak G}:X\twoheadrightarrow Q_{\mathfrak G},
    \qquad
    \bar\phi:=\bar\phi_{\mathfrak G}:Q_{\mathfrak G}\twoheadrightarrow Q_{X/S}.
\]
By Theorem~\ref{thm:formal-quotient-monomorphism-recognition},
\[
    \operatorname{dR}_S(q):
    X_{\mathrm{dR}/S}
    \longrightarrow
    (Q_{\mathfrak G})_{\mathrm{dR}/S}
\]
is an equivalence, and in particular a monomorphism.

\begin{definition}[Formal orbit disk]
\label{def:formal-orbit-disk}
Let $a:T\to X$ be an object of $(\mathcal H_{\Sp})_{/X}$.  Using the
pointwise-disk convention of Chapter~5, define the
$\mathfrak G$-formal orbit disk centred at $a$ by
\[
    \mathsf{Disk}^{\infty}_{\mathfrak G}(a)
    :=
    T\times_{Q_{\mathfrak G}}X.
\]
This notation agrees with the value at \(a\) of the formal-groupoid disk
endofunctor
\(\mathsf{Disk}^{\infty}_{\mathfrak G}=q_{\mathfrak G}^*\Sigma_{q_{\mathfrak G}}\)
defined in Definition~\ref{def:formal-groupoid-disk-jet}.
Write
\[
    \pi^{\mathfrak G}_a:
    \mathsf{Disk}^{\infty}_{\mathfrak G}(a)\to T
\]
for the first projection and
\[
    \operatorname{ev}^{\mathfrak G}_a:
    \mathsf{Disk}^{\infty}_{\mathfrak G}(a)\to X
\]
for the second.
\end{definition}

\begin{proposition}[Arrow-fibre presentation]
\label{prop:formal-disk-arrow-fibre}
There is a canonical equivalence
\[
    \mathsf{Disk}^{\infty}_{\mathfrak G}(a)
    \simeq
    T\times_{a,X,s}\mathfrak G_1,
\]
under which $\pi^{\mathfrak G}_a$ is the pullback of the source and
$\operatorname{ev}^{\mathfrak G}_a$ is the pulled-back target.
\end{proposition}

\begin{proof}
Groupoid effectivity gives
\[
    \mathfrak G_1
    \simeq
    X\times_{Q_{\mathfrak G}}X
\]
with source and target the two projections.  Pulling source back along
$a:T\to X$ yields the claimed equivalence.
\end{proof}


\begin{proposition}[Dependent infinitesimal-neighbourhood presentation]
\label{prop:formal-disk-dependent-presentation}
In the differential-cohesive internal notation, the formal orbit disk over
$t:T$ has fibre
\[
    \left(\mathsf{Disk}^{\infty}_{\mathfrak G}(a)\right)_t
    \simeq
    \sum_{x:X}
    \bigl(a(t)\sim_{\mathfrak G}x\bigr).
\]
Equivalently, the total object may be read as
\[
    \mathsf{Disk}^{\infty}_{\mathfrak G}(a)
    \simeq
    \left\{
        (t,x,\gamma)
        \ \middle|\
        t:T,\ x:X,\ \gamma:a(t)\sim_{\mathfrak G}x
    \right\}.
\]
Under this interpretation
$\pi^{\mathfrak G}_a$ forgets $(x,\gamma)$ and
$\operatorname{ev}^{\mathfrak G}_a$ forgets $(t,\gamma)$.
\end{proposition}

\begin{proof}
By the arrow-fibre presentation,
\[
    \mathsf{Disk}^{\infty}_{\mathfrak G}(a)
    \simeq
    T\times_{a,X,s}\mathfrak G_1.
\]
The fibre over $t:T$ is therefore the source fibre of $\mathfrak G_1$ at
$a(t)$.  Since
\[
    \mathfrak G_1\simeq X\times_{Q_{\mathfrak G}}X,
\]
that source fibre consists precisely of a point $x:X$ together with a
witness that $q_{\mathfrak G}(a(t))$ and $q_{\mathfrak G}(x)$ agree, which is
the notation $a(t)\sim_{\mathfrak G}x$.
\end{proof}

The groupoid unit induces a centre section
\[
    e^{\mathfrak G}_a:
    T\longrightarrow
    \mathsf{Disk}^{\infty}_{\mathfrak G}(a).
\]

\begin{proposition}[Intrinsic formality of formal orbit disks]
\label{prop:formal-disk-centre-formality}
The formal orbit disk is de Rham-formal along its centre:
\[
    \mathsf{Disk}^{\infty}_{\mathfrak G}(a)
    \times_{
        (\mathsf{Disk}^{\infty}_{\mathfrak G}(a))_{\mathrm{dR}/S}
    }
    T_{\mathrm{dR}/S}
    \xrightarrow{\simeq}
    \mathsf{Disk}^{\infty}_{\mathfrak G}(a).
\]
Equivalently,
\[
    \operatorname{dR}_S(e^{\mathfrak G}_a):
    T_{\mathrm{dR}/S}
    \xrightarrow{\simeq}
    (\mathsf{Disk}^{\infty}_{\mathfrak G}(a))_{\mathrm{dR}/S}
\]
is an equivalence.
\end{proposition}

\begin{proof}
Left exactness gives
\[
    (\mathsf{Disk}^{\infty}_{\mathfrak G}(a))_{\mathrm{dR}/S}
    \simeq
    T_{\mathrm{dR}/S}
    \times_{(Q_{\mathfrak G})_{\mathrm{dR}/S}}
    X_{\mathrm{dR}/S}.
\]
By Theorem~\ref{thm:formal-quotient-monomorphism-recognition},
\[
    \operatorname{dR}_S(q):
    X_{\mathrm{dR}/S}
    \xrightarrow{\simeq}
    (Q_{\mathfrak G})_{\mathrm{dR}/S}
\]
is an equivalence.  The projection
\[
    T_{\mathrm{dR}/S}
    \times_{(Q_{\mathfrak G})_{\mathrm{dR}/S}}
    X_{\mathrm{dR}/S}
    \longrightarrow
    T_{\mathrm{dR}/S}
\]
is therefore its base change along
$\operatorname{dR}_S(q\circ a):T_{\mathrm{dR}/S}\to
(Q_{\mathfrak G})_{\mathrm{dR}/S}$ and is an equivalence.  Its inverse is the
localized centre section
\[
    \operatorname{dR}_S(e^{\mathfrak G}_a):
    T_{\mathrm{dR}/S}
    \xrightarrow{\simeq}
    (\mathsf{Disk}^{\infty}_{\mathfrak G}(a))_{\mathrm{dR}/S}.
\]
Apply Proposition~\ref{prop:centred-formalization-recognition}.
\end{proof}

\begin{proposition}[Effectivity and infinitesimal equivalence of formal orbit disks]
\label{prop:formal-disk-effective}
The projection
\[
    \pi^{\mathfrak G}_a:
    \mathsf{Disk}^{\infty}_{\mathfrak G}(a)\twoheadrightarrow T
\]
is both an effective epimorphism and a relative de Rham equivalence.
\end{proposition}

\begin{proof}
It is the base change of
\[
    q:X\twoheadrightarrow Q_{\mathfrak G}
\]
along $qa:T\to Q_{\mathfrak G}$.  The quotient map $q$ is effective by
groupoid effectivity and is a relative de Rham equivalence by
Theorem~\ref{thm:formal-quotient-monomorphism-recognition}.  Effective
epimorphisms are stable under pullback, and Chapter~4 proves that the left
class of the de Rham--formally \'etale factorization system is pullback
stable.  Both properties therefore pass to $\pi^{\mathfrak G}_a$.
\end{proof}

\begin{proposition}[Comparison with the ambient relative de Rham disk]
\label{prop:formal-disk-ambient-comparison}
The quotient anchor
\[
    \bar\phi_{\mathfrak G}:Q_{\mathfrak G}\to Q_{X/S}
\]
induces a canonical morphism over $T$ and $X$
\[
    \mathrm{cmp}^{\mathrm{disk}}_{\mathfrak G,a}:
    \mathsf{Disk}^{\infty}_{\mathfrak G}(a)
    \longrightarrow
    \operatorname{Nbh}^{\mathrm{dR}}_{T/S}(T\times_SX)
    \simeq
    T\times_{Q_{X/S}}X.
\]
Here the relative de Rham neighbourhood is taken along the graph
$\Gamma_a:T\to T\times_SX$, as in Chapter~5.  These maps are natural in $a$
and in $\mathfrak G$.
\end{proposition}

\begin{proof}
Apply $\bar\phi_{\mathfrak G}$ to the base of the defining fibre product
$T\times_{Q_{\mathfrak G}}X$.  Naturality is functoriality of pullback.
\end{proof}

\begin{remark}[Formal leaves]
The fibre of
\[
    q:X\twoheadrightarrow Q_{\mathfrak G}
\]
through $a:T\to X$ is $\mathsf{Disk}^{\infty}_{\mathfrak G}(a)$.  Thus
$Q_{\mathfrak G}$ parametrizes the effective formal leaves.  Each generalized
leaf is an effective relative de Rham equivalence over its parameter object
and is intrinsically de Rham-formal along the chosen centre.  This is geometrically
parallel to the formal-leaf viewpoint in formal integration of partition Lie
algebroids, while the algebraic input and coefficient categories there are
different \cite{BrantnerMagidsonNuiten}.
\end{remark}

\begin{proposition}[Cartesianity of formal disks]
\label{prop:formal-disk-cartesian}
For $b:T'\to T$ and $a:T\to X$ there is a canonical equivalence
\[
    T'\times_T\mathsf{Disk}^{\infty}_{\mathfrak G}(a)
    \simeq
    \mathsf{Disk}^{\infty}_{\mathfrak G}(ab),
\]
compatible with evaluation, centre sections, identities, composition, and
higher pullback pasting.
\end{proposition}

\begin{proof}
Both sides are canonically
\[
    T'\times_{Q_{\mathfrak G}}X.
\]
\end{proof}


\section{Raw and formal isotropy}
\label{sec:flg-isotropy}

A general internal groupoid carries a raw automorphism object over its object
space.  In an $\infty$-topos this raw object also sees ambient higher inertia
of the object space itself.  For the formal theory we therefore distinguish it
from its centred de Rham formalization.  Throughout this section, the phrase
\emph{formal isotropy} means the formalized anchor-kernel isotropy constructed
below; it does not rename the entire raw automorphism object.

\subsection{Raw isotropy and inertia}

In this subsection let
\[
    \mathfrak G_\bullet\in
    \operatorname{Gpd}_X((\mathcal H_{\Sp})_{/S})
\]
be arbitrary; no formality is used in the raw isotropy constructions.

\begin{definition}[Raw isotropy group]
\label{def:raw-isotropy-group}
Let
\[
    (s,t):
    \mathfrak G_1\longrightarrow X\times_SX
\]
be the source-target map.  Define
\[
    I^{\mathrm{raw}}_{\mathfrak G}
    :=
    X\times_{\Delta_X,\;X\times_SX,\;(s,t)}
    \mathfrak G_1.
\]
Write
\[
    p_{\mathrm{raw}}:
    I^{\mathrm{raw}}_{\mathfrak G}\to X
\]
for its structural morphism and
\[
    e_{\mathrm{raw}}:
    X\longrightarrow I^{\mathrm{raw}}_{\mathfrak G}
\]
for the section induced by the groupoid unit.
\end{definition}

\begin{proposition}[Raw isotropy is a group object]
\label{prop:raw-isotropy-group}
The object
$I^{\mathrm{raw}}_{\mathfrak G}\to X$
carries a canonical group-object structure in $(\mathcal H_{\Sp})_{/X}$.
Multiplication, unit, and inverse are the restrictions of groupoid
composition, unit, and inversion.
\end{proposition}

\begin{proof}
The composable-pair object
\[
    I^{\mathrm{raw}}_{\mathfrak G}
    \times_X
    I^{\mathrm{raw}}_{\mathfrak G}
\]
maps canonically to
$\mathfrak G_1\times_{t,X,s}\mathfrak G_1$, and groupoid composition carries
it back to the source-target diagonal fibre.  Thus composition restricts to
multiplication on $I^{\mathrm{raw}}_{\mathfrak G}$.  Inversion preserves the
same fibre, and the groupoid unit factors through it.  The group identities
and all higher coherences are the corresponding restrictions of the
simplicial groupoid identities.
\end{proof}

\begin{definition}[Relative inertia object]
\label{def:relative-inertia-object}
For an object $Q\to S$, regard the diagonal
\[
    \Delta_Q:Q\longrightarrow Q\times_SQ
\]
as a pointed object of the slice $(\mathcal H_{\Sp})_{/Q}$ through the first
projection.  Define its relative inertia object by
\[
    \mathcal I_{Q/S}
    :=
    Q\times_{\Delta_Q,\;Q\times_SQ,\;\Delta_Q}Q.
\]
\end{definition}

\begin{proposition}[Canonical group structure on relative inertia]
\label{prop:relative-inertia-group}
The object
\[
    \mathcal I_{Q/S}\longrightarrow Q
\]
carries a canonical group-object structure in $(\mathcal H_{\Sp})_{/Q}$.  For every
morphism $f:Q\to Q'$ over $S$, functoriality supplies a canonical homomorphism
of group objects over $Q$
\[
    \mathcal I_{Q/S}
    \longrightarrow
    f^*\mathcal I_{Q'/S}.
\]
These morphisms are functorial in $f$ and compatible with arbitrary base
change in $S$.
\end{proposition}

\begin{proof}
Form in the slice $(\mathcal H_{\Sp})_{/Q}$ the \v{C}ech nerve of the section
\[
    \Delta_Q:
    Q\longrightarrow
    (Q\times_SQ\xrightarrow{\operatorname{pr}_1}Q).
\]
Its degree-zero object is the terminal object $Q$ of the slice and its
degree-one object is
\[
    Q\times_{Q\times_SQ}Q
    =
    \mathcal I_{Q/S}.
\]
Every \v{C}ech nerve is a groupoid object.  A groupoid object whose object
space is terminal is equivalently a group object, so the unit, multiplication,
and inverse of $\mathcal I_{Q/S}$ are the degree-zero, degree-two, and inversion
maps of this \v{C}ech groupoid.  Functoriality follows from functoriality of
diagonals and finite limits.  Pullback in an $\infty$-topos preserves the
entire defining \v{C}ech nerve, giving the asserted base-change compatibility.
\end{proof}

\begin{proposition}[Raw isotropy as pulled-back inertia]
\label{prop:isotropy-inertia}
There is a canonical equivalence of group objects over $X$
\[
    I^{\mathrm{raw}}_{\mathfrak G}
    \simeq
    X\times_{Q_{\mathfrak G}}
    \mathcal I_{Q_{\mathfrak G}/S}.
\]
\end{proposition}

\begin{proof}
Using
\[
    \mathfrak G_1
    \simeq
    X\times_{Q_{\mathfrak G}}X,
\]
substitute the effective relation presentation into the defining pullback for
$I^{\mathrm{raw}}_{\mathfrak G}$.  Finite-limit pasting identifies the
result with the pullback to $X$ of
\[
    Q_{\mathfrak G}
    \times_{Q_{\mathfrak G}\times_SQ_{\mathfrak G}}
    Q_{\mathfrak G}.
\]
Composition and inversion agree with loop concatenation and reversal, so the
equivalence respects the group structures.
\end{proof}

\subsection{Formal isotropy as the infinitesimal anchor kernel}

From now through the formal-isotropy constructions, let
\[
    \mathfrak G_\bullet\in\operatorname{FormLieGpd}(X/S).
\]

\begin{definition}[Formal isotropy group]
\label{def:isotropy-formal-group}
Define the \emph{formal isotropy group} by centred formalization along the raw
identity section:
\[
    I^{\mathrm{for}}_{\mathfrak G}
    :=
    P^{\mathrm{dR}}(e_{\mathrm{raw}})
    =
    I^{\mathrm{raw}}_{\mathfrak G}
    \times_{
        (I^{\mathrm{raw}}_{\mathfrak G})_{\mathrm{dR}/S}
    }
    X_{\mathrm{dR}/S}.
\]
Write
\[
    r_{e_{\mathrm{raw}}}:
    I^{\mathrm{for}}_{\mathfrak G}
    \longrightarrow
    I^{\mathrm{raw}}_{\mathfrak G}
\]
for the formalization counit.
\end{definition}

\begin{theorem}[Formal isotropy group theorem]
\label{thm:isotropy-group}
The object
\[
    I^{\mathrm{for}}_{\mathfrak G}\to X
\]
is canonically a group object in $(\mathcal H_{\Sp})_{/X}$, the morphism
$r_{e_{\mathrm{raw}}}$ is a homomorphism, and its identity section is de Rham-formal.  Moreover there is a canonical equivalence of group objects
over $X$
\[
    I^{\mathrm{for}}_{\mathfrak G}
    \xrightarrow{\simeq}
    \mathfrak G_1
    \times_{\mathcal R^{\mathrm{dR}}_{X/S}}
    X,
\]
where $\mathfrak G_1\to \mathcal R^{\mathrm{dR}}_{X/S}$ is the canonical degree-one anchor and
$X\to \mathcal R^{\mathrm{dR}}_{X/S}$ is the unit section.
\end{theorem}

\begin{proof}
The raw multiplication is a morphism of centred objects from
\[
    X\xrightarrow{\,e_{\mathrm{raw}}\times_Xe_{\mathrm{raw}}\,}
    I^{\mathrm{raw}}_{\mathfrak G}\times_XI^{\mathrm{raw}}_{\mathfrak G}
\]
to
\[
    X\xrightarrow{e_{\mathrm{raw}}}I^{\mathrm{raw}}_{\mathfrak G}.
\]
Since $P_X^{\mathrm{dR}}$ preserves pullbacks and fixes
$\operatorname{id}_X$,
\[
    P^{\mathrm{dR}}
    \bigl(e_{\mathrm{raw}}\times_Xe_{\mathrm{raw}}\bigr)
    \simeq
    I^{\mathrm{for}}_{\mathfrak G}
    \times_X
    I^{\mathrm{for}}_{\mathfrak G}.
\]
Applying $P_X^{\mathrm{dR}}$ to multiplication, unit, and inversion therefore gives a
group-object structure on $I^{\mathrm{for}}_{\mathfrak G}$.  Naturality of the fixed-source counit makes $r_{e_{\mathrm{raw}}}$ a homomorphism.  De Rham formality of the identity
section is Theorem~\ref{thm:centred-formalization-idempotent}.

For the kernel formula, the raw isotropy square is a pullback diagram of
centred objects:
\[
    I^{\mathrm{raw}}_{\mathfrak G}
    \simeq
    \mathfrak G_1
    \times_{X\times_SX}
    X,
\]
where the centres are respectively the groupoid unit, the diagonal, and
$\operatorname{id}_X$.  Apply $P_X^{\mathrm{dR}}$.  Since $\mathfrak G$ is formal,
\[
    P^{\mathrm{dR}}(u)\simeq\mathfrak G_1.
\]
By Theorem~\ref{thm:pair-groupoid-formalization},
\[
    P^{\mathrm{dR}}(\Delta_X)\simeq \mathcal R^{\mathrm{dR}}_{X/S},
\]
and
\[
    P^{\mathrm{dR}}(\operatorname{id}_X)\simeq X.
\]
Pullback preservation therefore yields
\[
    I^{\mathrm{for}}_{\mathfrak G}
    \simeq
    \mathfrak G_1\times_{\mathcal R^{\mathrm{dR}}_{X/S}}X.
\]
The formalized map $\mathfrak G_1\to \mathcal R^{\mathrm{dR}}_{X/S}$ is the canonical anchor by
Corollary~\ref{cor:canonical-formal-anchor}.  The multiplication on raw
isotropy is the restriction of groupoid composition, and its domain is the
pullback
\[
    I^{\mathrm{raw}}_{\mathfrak G}
    \times_X
    I^{\mathrm{raw}}_{\mathfrak G}.
\]
Since $P_X^{\mathrm{dR}}$ preserves this pullback and the formalized multiplication is
obtained by applying $P_X^{\mathrm{dR}}$ to the same composition map, the displayed
anchor-kernel equivalence intertwines multiplication, unit, and inversion.
It is therefore an equivalence of group objects.
\end{proof}

\begin{remark}[Infinitesimal automorphisms and the anchor kernel]
In the internal infinitesimal-equality notation, the fibre of raw isotropy at
$x:X$ is
\[
    I^{\mathrm{raw}}_{\mathfrak G,x}
    \simeq
    \bigl(x\sim_{\mathfrak G}x\bigr).
\]
The canonical anchor sends such an automorphism to an ambient infinitesimal
automorphism
\[
    \phi_{\mathfrak G}(\gamma):
    x\sim_{\mathrm{dR}/S}x.
\]
The anchor-kernel theorem therefore reads
\[
    I^{\mathrm{for}}_{\mathfrak G,x}
    \simeq
    \left\{
        \gamma:x\sim_{\mathfrak G}x
        \ \middle|\
        \phi_{\mathfrak G}(\gamma)\simeq 1_x
    \right\}.
\]
This is only the dependent form of the homotopy pullback
\[
    \mathfrak G_1\times_{\mathcal R^{\mathrm{dR}}_{X/S}}X;
\]
it introduces no additional notion of kernel.
\end{remark}

\begin{remark}[Why raw and formal isotropy are distinct]
\label{rem:raw-formal-isotropy-distinct}
The formality of $\mathfrak G_\bullet$ does not imply, from left exactness
alone, that $I^{\mathrm{raw}}_{\mathfrak G}$ is de Rham-formal along its identity.
This is a genuine categorical issue rather than a proof artifact.

The following is an abstract left-exact-localization counterexample; it is not
asserted to arise from the spectral de Rham reflector.  In the
$\infty$-topos of spaces take the accessible left-exact localization
$L=\tau_{\le1}$ and
\[
    X=B\mathbb Z\times K(\mathbb Z,2),
    \qquad
    LX\simeq B\mathbb Z.
\]
The maximal $L$-formal relation
\[
    X\times_{B\mathbb Z}X\rightrightarrows X
\]
is formal.  Since
\[
    \mathcal I_{B\mathbb Z}
    \simeq
    B\mathbb Z\times\mathbb Z,
\]
its raw isotropy is
\[
    I^{\mathrm{raw}}
    \simeq
    X\times\mathbb Z.
\]
The identity section lies in the zero component, and centred formalization gives
\[
    P^{\mathrm{dR}}(e_{\mathrm{raw}})
    \simeq
    (X\times\mathbb Z)
    \times_{B\mathbb Z\times\mathbb Z}
    B\mathbb Z
    \simeq X,
\]
which is not equivalent to $X\times\mathbb Z$.  Thus the separate formal
isotropy construction is essential at the level of general left-exact
localizations.  The de Rham and jet constructions of Chapters~4--5 contain no theorem collapsing this
distinction.
\end{remark}

\subsection{The descended formal inertia kernel}

Put
\[
    Q:=Q_{\mathfrak G},
    \qquad
    \bar\phi:=\bar\phi_{\mathfrak G}:Q\to Q_{X/S}.
\]
Functoriality of relative inertia gives a canonical homomorphism of group objects over
$Q$
\[
    \mathcal I_{Q/S}
    \longrightarrow
    \bar\phi^*\mathcal I_{Q_{X/S}/S}.
\]

\begin{definition}[Formal inertia kernel]
\label{def:formal-inertia-kernel}
Define
\[
    \mathcal K_{\mathfrak G}
    :=
    \mathcal I_{Q/S}
    \times_{\bar\phi^*\mathcal I_{Q_{X/S}/S}}
    Q,
\]
where
\[
    Q\longrightarrow \bar\phi^*\mathcal I_{Q_{X/S}/S}
\]
is the identity-loop section.  This is a group object over $Q$.  No second
``adjoint-group'' symbol is introduced: the descended formal inertia kernel
$\mathcal K_{\mathfrak G}$ is the single group object whose pullback carries
the formal isotropy and its conjugation transport.
\end{definition}

\begin{lemma}[Effective descent for internal group objects]
\label{lem:effective-descent-group-objects}
Let
\[
    e:U\twoheadrightarrow V
\]
be an effective epimorphism in $(\mathcal H_{\Sp})_{/S}$, and let
$U_\bullet=\check C_\bullet(e)$.  Pullback induces a canonical equivalence
\[
    \operatorname{Grp}((\mathcal H_{\Sp})_{/V})
    \xrightarrow{\simeq}
    \operatorname*{Tot}_{[n]\in\Delta}
    \operatorname{Grp}((\mathcal H_{\Sp})_{/U_n}),
\]
where $\operatorname{Grp}$ denotes internal group objects.  The equivalence is
compatible with the forgetful functors to the corresponding slice
$\infty$-categories and with arbitrary base change.
\end{lemma}

\begin{proof}
Effective descent in an $\infty$-topos gives
\[
    (\mathcal H_{\Sp})_{/V}
    \xrightarrow{\simeq}
    \operatorname*{Tot}_{[n]\in\Delta}
    (\mathcal H_{\Sp})_{/U_n},
\]
with every transition functor given by pullback and therefore preserving
finite limits.

An internal group object is a model of the finite-product algebraic theory
of groups, equivalently a reduced Segal groupoid object.  Functor categories
commute with limits in the target, and the condition that a functor preserve
the finite products appearing in that theory is detected componentwise
because every transition functor in the displayed totalization preserves
finite limits.  Passing to the corresponding full subcategories gives
\[
    \operatorname{Grp}((\mathcal H_{\Sp})_{/V})
    \simeq
    \operatorname*{Tot}_{[n]\in\Delta}
    \operatorname{Grp}((\mathcal H_{\Sp})_{/U_n}).
\]
The underlying-object and base-change compatibilities are inherited from the
effective-descent equivalence of slices.
\end{proof}

\begin{theorem}[Descent formula for formal isotropy]
\label{thm:formal-isotropy-descent}
There are canonical equivalences of group objects
\[
    q^*\mathcal I_{Q/S}
    \simeq
    I^{\mathrm{raw}}_{\mathfrak G},
    \qquad
    q^*\mathcal K_{\mathfrak G}
    \simeq
    I^{\mathrm{for}}_{\mathfrak G}.
\]
The second equivalence transports the canonical \v{C}ech descent datum on
$q^*\mathcal K_{\mathfrak G}$ to
$I^{\mathrm{for}}_{\mathfrak G}$.  Under
Lemma~\ref{lem:effective-descent-group-objects},
$\mathcal K_{\mathfrak G}$ is the contractibly unique group object over
$Q$ realizing this descended formal isotropy.  Thus formal isotropy descends canonically, rather than by an
additional descent datum, to the formal leaf space as the kernel of the
inertia morphism induced by
$\bar\phi:Q\to Q_{X/S}$.
\end{theorem}

\begin{proof}
The first equivalence is
Proposition~\ref{prop:isotropy-inertia}.  For the second, pull the defining
finite limit of $\mathcal K_{\mathfrak G}$ back along $q$.  Since
$\bar\phi q\simeq e_{X/S}$,
\[
\begin{aligned}
    q^*\mathcal K_{\mathfrak G}
    &\simeq
    q^*\mathcal I_{Q/S}
    \times_{q^*\bar\phi^*\mathcal I_{Q_{X/S}/S}}
    X\\
    &\simeq
    I^{\mathrm{raw}}_{\mathfrak G}
    \times_{e_{X/S}^*\mathcal I_{Q_{X/S}/S}}
    X.
\end{aligned}
\]
The homomorphism in the second line is the pullback of the inertia morphism
induced by $\bar\phi:Q\to Q_{X/S}$.  Under
\[
    \mathfrak G_\bullet\simeq\check C_\bullet(q),
    \qquad
    \mathcal R^{\mathrm{dR}}_{X/S,\bullet}\simeq\check C_\bullet(e_{X/S}),
\]
it is exactly the homomorphism on raw isotropy induced by the canonical anchor
$\mathfrak G_\bullet\to \mathcal R^{\mathrm{dR}}_{X/S,\bullet}$.

The identity-loop section of
$e_{X/S}^*\mathcal I_{Q_{X/S}/S}$ is the isotropy of the unit section
$X\to \mathcal R^{\mathrm{dR}}_{X/S}$.  Taking its fibre therefore gives the finite-limit identity
\[
    I^{\mathrm{raw}}_{\mathfrak G}
    \times_{e_{X/S}^*\mathcal I_{Q_{X/S}/S}}
    X
    \simeq
    \mathfrak G_1\times_{\mathcal R^{\mathrm{dR}}_{X/S}}X.
\]
Theorem~\ref{thm:isotropy-group} identifies the right-hand side with
$I^{\mathrm{for}}_{\mathfrak G}$ as a group object over $X$.

Since $q:X\twoheadrightarrow Q$ is an effective epimorphism,
Lemma~\ref{lem:effective-descent-group-objects} identifies group objects over
$Q$ with group objects over the complete \v{C}ech nerve
$\check C_\bullet(q)\simeq\mathfrak G_\bullet$ equipped with their coherent
descent data.  The pullback $q^*\mathcal K_{\mathfrak G}$ carries the
tautological descent datum of the group object
$\mathcal K_{\mathfrak G}\to Q$.  Transporting that datum through the
canonical equivalence
\[
    q^*\mathcal K_{\mathfrak G}
    \xrightarrow{\simeq}
    I^{\mathrm{for}}_{\mathfrak G}
\]
constructs coherent descent data on formal isotropy.  Because the descent
functor is an equivalence of $\infty$-categories, the homotopy fibre over this
specified descent object is contractible.  Hence
$\mathcal K_{\mathfrak G}$ is the contractibly unique descended group object
realizing that formal-isotropy descent datum.
\end{proof}

\subsection{The nonlinear formal-isotropy commutator}

\begin{definition}[Nonlinear formal-isotropy commutator]
\label{def:isotropy-commutator}
Define
\[
    \operatorname{comm}_{\mathfrak G}:
    I^{\mathrm{for}}_{\mathfrak G}
    \times_X
    I^{\mathrm{for}}_{\mathfrak G}
    \longrightarrow
    I^{\mathrm{for}}_{\mathfrak G}
\]
internally by
\[
    \operatorname{comm}_{\mathfrak G}(a,b)
    :=
    aba^{-1}b^{-1}.
\]
\end{definition}

\begin{proposition}[Normalization and functoriality of the commutator]
\label{prop:isotropy-commutator-properties}
The commutator is canonically trivial when either input is the identity.
A morphism of formal Lie groupoids carries formal-isotropy commutators to
formal-isotropy commutators.
\end{proposition}

\begin{proof}
The normalization identities are group identities in
$I^{\mathrm{for}}_{\mathfrak G}$.  Functoriality follows because a groupoid
functor preserves composition, unit, inversion, and the canonical anchors,
hence preserves the anchor kernels.
\end{proof}

\begin{remark}[Not yet a bracket]
\label{rem:commutator-not-bracket}
The map $\operatorname{comm}_{\mathfrak G}$ is multiplicative and nonlinear.
No bilinearity, antisymmetry, Jacobi identity, or additive truncation is
asserted.  Later brackets and higher partition-Lie operations arise only
after differentiation and the corresponding formal-moduli recognition
theorems.  No comparison with, or reduction to, an ordinary binary dg Lie
bracket is asserted at this nonlinear stage.
\end{remark}


\section{Infinitesimal differences, arrow fibres, and bitorsors}
\label{sec:flg-bitorsor}

For a one-object group, multiplication identifies every arrow fibre with the
identity fibre.  For a general groupoid two different effective arrow images
must be distinguished: the image of the source-target map and the image of
the infinitesimal anchor.

\begin{theorem}[Kernel bitorsor theorem for fixed-object groupoid morphisms]
\label{thm:kernel-bitorsor}
Let
\[
    \psi_\bullet:
    \mathfrak G_\bullet
    \longrightarrow
    \mathfrak H_\bullet
\]
be a morphism in
$\operatorname{Gpd}_X((\mathcal H_{\Sp})_{/S})$.  Write
$u_{\mathfrak H}:X\to\mathfrak H_1$ for the unit of the target groupoid and
define the arrow-kernel
\[
    K_\psi
    :=
    \mathfrak G_1
    \times_{\psi_1,\;\mathfrak H_1,\;u_{\mathfrak H}}
    X.
\]
Then $K_\psi\to X$ is canonically a group object.  Factor the degree-one map
by its effective image
\[
    \mathfrak G_1
    \xrightarrow{\pi_\psi}
    M_\psi
    \xrightarrow{j_\psi}
    \mathfrak H_1,
\]
with $\pi_\psi$ an effective epimorphism and $j_\psi$ a monomorphism.  The
kernel acts on $\mathfrak G_1$ by target and source composition, and the two
associated shear morphisms are canonical equivalences
\[
\begin{aligned}
    \mathsf{sh}_t:
    t^*K_\psi
    &\xrightarrow{\simeq}
    \mathfrak G_1\times_{M_\psi}\mathfrak G_1,
    &
    (g,h)&\longmapsto(h\circ g,g),\\
    \mathsf{sh}_s:
    s^*K_\psi
    &\xrightarrow{\simeq}
    \mathfrak G_1\times_{M_\psi}\mathfrak G_1,
    &
    (g,k)&\longmapsto(g\circ k,g).
\end{aligned}
\]
The inverse morphisms are represented internally by
\[
    (g_1,g_2)
    \longmapsto
    (g_2,g_1\circ g_2^{-1})
\]
and
\[
    (g_1,g_2)
    \longmapsto
    (g_2,g_2^{-1}\circ g_1),
\]
where the required morphisms into the homotopy pullbacks defining
$t^*K_\psi$ and $s^*K_\psi$ are supplied canonically by the specified
homotopy in
$\mathfrak G_1\times_{\mathfrak H_1}\mathfrak G_1$.
\end{theorem}

\begin{proof}
Because $\psi_\bullet$ is a fixed-object groupoid functor, it induces a
homomorphism of raw isotropy group objects
\[
    I^{\mathrm{raw}}_{\mathfrak G}
    \longrightarrow
    I^{\mathrm{raw}}_{\mathfrak H}.
\]
Finite-limit pasting identifies the arrow-kernel with the group-object fibre
over the identity section:
\[
    K_\psi
    \simeq
    I^{\mathrm{raw}}_{\mathfrak G}
    \times_{I^{\mathrm{raw}}_{\mathfrak H}}
    X.
\]
Limits of internal group objects are created in the underlying slice, so this
pullback gives the canonical group-object structure on $K_\psi\to X$.

Since $j_\psi$ is a monomorphism, finite-limit pasting gives a canonical
equivalence
\[
    P_\psi
    :=
    \mathfrak G_1\times_{M_\psi}\mathfrak G_1
    \xrightarrow{\simeq}
    \mathfrak G_1\times_{\mathfrak H_1}\mathfrak G_1.
\]
Let
\[
    p_1,p_2:P_\psi\rightrightarrows\mathfrak G_1
\]
be the projections.  Through the right-hand homotopy-pullback presentation,
$P_\psi$ carries its universal specified homotopy
\[
    \gamma:
    \psi_1p_1
    \simeq
    \psi_1p_2.
\]
Applying source and target to $\gamma$, and using that $\psi_\bullet$ is the
identity in degree zero, supplies the specified endpoint identifications
needed to form the internal difference morphisms
\[
    d_t:=p_1\circ p_2^{-1},
    \qquad
    d_s:=p_2^{-1}\circ p_1.
\]
Here and below the displayed composition notation denotes the morphisms
obtained from the projections, inversion, the composable-arrow pullback, and
groupoid composition; no strict equality of generalized elements is being
used.

Functoriality of composition and inversion under $\psi_\bullet$ carries the
homotopy $\gamma$ to canonical homotopies
\[
\begin{aligned}
    \psi_1d_t
    &\simeq
    \psi_1p_1\circ(\psi_1p_2)^{-1}
    \simeq
    u_{\mathfrak H}\,t p_2,\\
    \psi_1d_s
    &\simeq
    (\psi_1p_2)^{-1}\circ\psi_1p_1
    \simeq
    u_{\mathfrak H}\,s p_2.
\end{aligned}
\]
These specified homotopies are exactly the data required by the two homotopy
pullbacks.  Hence the universal properties of those pullbacks give canonical
morphisms
\[
    \delta_t:
    P_\psi\longrightarrow t^*K_\psi,
    \qquad
    \delta_s:
    P_\psi\longrightarrow s^*K_\psi
\]
whose first components are $p_2$ and whose kernel components are $d_t$ and
$d_s$, respectively.

Conversely, composition with a kernel arrow has unchanged image under
$\psi_1$.  More precisely, the defining homotopy
$\psi_1(h)\simeq u_{\mathfrak H}$ for a target-kernel arrow gives
\[
    \psi_1(h\circ g)
    \simeq
    \psi_1(h)\circ\psi_1(g)
    \simeq
    \psi_1(g),
\]
and the analogous source-kernel calculation gives
$\psi_1(g\circ k)\simeq\psi_1(g)$.  Thus the two action maps, together with
the unchanged arrow $g$, define the shear morphisms
$\mathsf{sh}_t$ and $\mathsf{sh}_s$ into $P_\psi$.

The simplicial groupoid unit, inverse, and associativity coherences give
\[
    (p_1\circ p_2^{-1})\circ p_2\simeq p_1,
    \qquad
    p_2\circ(p_2^{-1}\circ p_1)\simeq p_1,
\]
in the two relevant composable-arrow diagrams, and similarly
\[
    (h\circ g)\circ g^{-1}\simeq h,
    \qquad
    g^{-1}\circ(g\circ k)\simeq k.
\]
Tracing the specified pullback homotopies shows that these are identities in
the corresponding homotopy pullbacks, not merely identities after forgetting
to $\mathfrak G_1$.  Therefore
\[
    \mathsf{sh}_t\delta_t\simeq\operatorname{id},
    \qquad
    \delta_t\mathsf{sh}_t\simeq\operatorname{id},
    \qquad
    \mathsf{sh}_s\delta_s\simeq\operatorname{id},
    \qquad
    \delta_s\mathsf{sh}_s\simeq\operatorname{id}.
\]
Hence both shear morphisms are equivalences with the stated inverses.  All
higher compatibility is inherited functorially from the groupoid coherences
and the universal properties of the finite homotopy pullbacks.
\end{proof}

\subsection{Raw source-target bitorsors}

Factor the source-target map as
\[
    \mathfrak G_1
    \xrightarrow{\pi_{\mathrm{raw}}}
    M^{\mathrm{raw}}_{\mathfrak G}
    \xrightarrow{j_{\mathrm{raw}}}
    X\times_SX,
\]
where $\pi_{\mathrm{raw}}$ is an effective epimorphism and
$j_{\mathrm{raw}}$ is a monomorphism.

\begin{proposition}[Raw isotropy actions]
\label{prop:isotropy-actions}
Groupoid composition induces commuting left and right actions
\[
    a_L:
    t^*I^{\mathrm{raw}}_{\mathfrak G}
    \longrightarrow
    \mathfrak G_1,
    \qquad
    a_L(g,h)=h\circ g,
\]
and
\[
    a_R:
    s^*I^{\mathrm{raw}}_{\mathfrak G}
    \longrightarrow
    \mathfrak G_1,
    \qquad
    a_R(g,k)=g\circ k.
\]
\end{proposition}

\begin{proof}
The action maps are restrictions of groupoid composition.  Associativity gives
the action laws and
\[
    (h\circ g)\circ k
    \simeq
    h\circ(g\circ k).
\]
All higher compatibility is inherited from the simplicial associativity
coherence.
\end{proof}

\begin{theorem}[Raw Isotropy Bitorsor Theorem]
\label{thm:raw-bitorsor}
There are canonical equivalences
\[
\begin{aligned}
    t^*I^{\mathrm{raw}}_{\mathfrak G}
    &\xrightarrow{\simeq}
    \mathfrak G_1
    \times_{M^{\mathrm{raw}}_{\mathfrak G}}
    \mathfrak G_1,
    &
    (g,h)&\longmapsto(h\circ g,g),\\
    s^*I^{\mathrm{raw}}_{\mathfrak G}
    &\xrightarrow{\simeq}
    \mathfrak G_1
    \times_{M^{\mathrm{raw}}_{\mathfrak G}}
    \mathfrak G_1,
    &
    (g,k)&\longmapsto(g\circ k,g).
\end{aligned}
\]
Their inverses are
\[
    (g_1,g_2)
    \longmapsto
    (g_2,g_1\circ g_2^{-1})
\]
and
\[
    (g_1,g_2)
    \longmapsto
    (g_2,g_2^{-1}\circ g_1),
\]
respectively.
\end{theorem}

\begin{proof}
Apply Theorem~\ref{thm:kernel-bitorsor} to the canonical fixed-object vertex
morphism
\[
    \psi_\bullet:
    \mathfrak G_\bullet
    \longrightarrow
    \operatorname{Pair}(X/S)_\bullet.
\]
Its degree-one component is
\[
    \psi_1=(s,t):
    \mathfrak G_1\longrightarrow X\times_SX.
\]
The arrow-kernel of $\psi_1$ over the diagonal unit is exactly
$I^{\mathrm{raw}}_{\mathfrak G}$, and the effective image
$M_\psi$ is $M^{\mathrm{raw}}_{\mathfrak G}$.  The two shear equivalences and
their internally constructed inverse morphisms are therefore precisely the
displayed maps.
\end{proof}

\begin{remark}
Fibrewise, this says that every \emph{inhabited} source-target arrow fibre is
a left torsor for target raw isotropy and a right torsor for source raw
isotropy.  Passing through
$M^{\mathrm{raw}}_{\mathfrak G}$ rather than all of $X\times_SX$ records the
possible absence of arrows between two objects.
\end{remark}

\subsection{Formal anchor bitorsors}

Factor the canonical degree-one anchor as
\[
    \mathfrak G_1
    \xrightarrow{\pi_{\mathrm{for}}}
    M^{\mathrm{for}}_{\mathfrak G}
    \xrightarrow{j_{\mathrm{for}}}
    \mathcal R^{\mathrm{dR}}_{X/S},
\]
with $\pi_{\mathrm{for}}$ effective and $j_{\mathrm{for}}$ monic.

\begin{proposition}[Formal isotropy actions preserve anchor fibres]
\label{prop:formal-isotropy-actions}
The left and right raw isotropy actions, precomposed with the formalization
homomorphism
\[
    r_{e_{\mathrm{raw}}}:
    I^{\mathrm{for}}_{\mathfrak G}
    \longrightarrow
    I^{\mathrm{raw}}_{\mathfrak G},
\]
induce canonically commuting actions
\[
    t^*I^{\mathrm{for}}_{\mathfrak G}
    \longrightarrow
    \mathfrak G_1,
    \qquad
    s^*I^{\mathrm{for}}_{\mathfrak G}
    \longrightarrow
    \mathfrak G_1
\]
which preserve the fibres of
$\pi_{\mathrm{for}}:\mathfrak G_1\to
M^{\mathrm{for}}_{\mathfrak G}$.
\end{proposition}

\begin{proof}
By Theorem~\ref{thm:isotropy-group}, formal isotropy is canonically the
homotopy fibre
\[
    I^{\mathrm{for}}_{\mathfrak G}
    \simeq
    \mathfrak G_1\times_{\mathcal R^{\mathrm{dR}}_{X/S}}X
\]
of the canonical anchor over the unit section.  Under this identification,
the composite with
$r_{e_{\mathrm{raw}}}:I^{\mathrm{for}}_{\mathfrak G}\to
I^{\mathrm{raw}}_{\mathfrak G}$
is the canonical map from this homotopy fibre to raw isotropy.  Since
$\phi_{\mathfrak G}$ is a groupoid morphism, composing an arrow $g$ with a
formal-isotropy arrow $h$ gives
\[
    \phi(h\circ g)
    \simeq
    \phi(h)\circ\phi(g)
    \simeq
    1\circ\phi(g)
    \simeq
    \phi(g),
\]
and similarly on the right.  Thus the actions obtained by precomposition with
$r_{e_{\mathrm{raw}}}$ preserve anchor fibres.  The action coherences are inherited from
the raw actions.
\end{proof}

\begin{theorem}[Formal Bitorsor Theorem]
\label{thm:formal-bitorsor}
There are canonical equivalences
\[
\begin{aligned}
    t^*I^{\mathrm{for}}_{\mathfrak G}
    &\xrightarrow{\simeq}
    \mathfrak G_1
    \times_{M^{\mathrm{for}}_{\mathfrak G}}
    \mathfrak G_1,
    &
    (g,h)&\longmapsto(h\circ g,g),\\
    s^*I^{\mathrm{for}}_{\mathfrak G}
    &\xrightarrow{\simeq}
    \mathfrak G_1
    \times_{M^{\mathrm{for}}_{\mathfrak G}}
    \mathfrak G_1,
    &
    (g,k)&\longmapsto(g\circ k,g).
\end{aligned}
\]
Equivalently, because $j_{\mathrm{for}}$ is monic,
\[
    t^*I^{\mathrm{for}}_{\mathfrak G}
    \simeq
    \mathfrak G_1\times_{\mathcal R^{\mathrm{dR}}_{X/S}}\mathfrak G_1
    \simeq
    s^*I^{\mathrm{for}}_{\mathfrak G}.
\]
Thus every inhabited fibre of the infinitesimal anchor is canonically a
left torsor for target formal isotropy and a right torsor for source formal
isotropy, with the two torsor structures commuting coherently.
\end{theorem}

\begin{proof}
Apply Theorem~\ref{thm:kernel-bitorsor} to the canonical infinitesimal anchor
\[
    \phi_{\mathfrak G}:
    \mathfrak G_\bullet
    \longrightarrow
    \mathcal R^{\mathrm{dR}}_{X/S,\bullet}.
\]
Its arrow-kernel is
\[
    K_{\phi_{\mathfrak G}}
    =
    \mathfrak G_1\times_{\mathcal R^{\mathrm{dR}}_{X/S}}X
    \simeq
    I^{\mathrm{for}}_{\mathfrak G}
\]
by Theorem~\ref{thm:isotropy-group}, and its effective arrow image is
$M^{\mathrm{for}}_{\mathfrak G}$.  The general kernel-bitorsor theorem
therefore gives both displayed equivalences, including the specified
homotopy-pullback witnesses in the inverse maps.
\end{proof}

\begin{remark}[Infinitesimal difference calculus]
The formal bitorsor theorem permits the following internal manipulation.  If
\[
    g_1,g_2:x\sim_{\mathfrak G}y
\]
have the same image under the canonical infinitesimal anchor, then their
target-side difference
\[
    g_1g_2^{-1}
\]
is a formal isotropy element at $y$, while their source-side difference
\[
    g_2^{-1}g_1
\]
is a formal isotropy element at $x$.  Conversely, composing $g_2$ on either
side by the corresponding formal isotropy element recovers $g_1$.  The two
shear equivalences of Theorem~\ref{thm:formal-bitorsor} are the external
homotopy-pullback semantics of these two difference calculations.
\end{remark}

\begin{remark}[Why the one-object case is exceptional]
\label{rem:group-case-exceptional}
For a one-object groupoid all arrows have the same source and target and the
formal isotropy group is the arrow group itself.  The bitorsor therefore
admits a global translation trivialization by the single identity fibre.
For a general groupoid formal isotropy varies over $X$, and the intrinsic
replacement for a global translation coordinate is the formal anchor
bitorsor above.
\end{remark}


\section{Conjugation and the formal adjoint local system}
\label{sec:flg-adjoint}

The descended inertia kernel already carries all transport coherence.  The only
remaining point is to identify its canonical \v{C}ech transport with internal
groupoid conjugation.  We make that identification once at the level of raw
inertia and then pass to the formal kernel by finite limits.

\begin{definition}[Internal conjugation on raw isotropy]
\label{def:raw-conjugation}
For an arbitrary fixed-object groupoid
$\mathfrak G_\bullet\rightrightarrows X$, composition and inversion define a
canonical homomorphism of group objects over $\mathfrak G_1$
\[
    \operatorname{conj}^{\mathrm{raw}}_{\mathfrak G}:
    s^*I^{\mathrm{raw}}_{\mathfrak G}
    \longrightarrow
    t^*I^{\mathrm{raw}}_{\mathfrak G}.
\]
It is the morphism obtained by sending an arrow $g:x\to y$ and an isotropy
arrow $h:x\to x$ to the composite
\[
    g\circ h\circ g^{-1}:y\longrightarrow y.
\]
The displayed generalized-element formula abbreviates the morphism constructed
from the composable-arrow pullbacks, inversion, and groupoid composition.
\end{definition}

\begin{theorem}[\v{C}ech inertia transport is conjugation]
\label{thm:inertia-descent-conjugation}
Let
\[
    q:X\twoheadrightarrow Q
\]
be an effective epimorphism in $(\mathcal H_{\Sp})_{/S}$, and put
$\mathfrak G_\bullet:=\check C_\bullet(q)$.  Under the canonical equivalence
of group objects
\[
    q^*\mathcal I_{Q/S}
    \xrightarrow{\simeq}
    I^{\mathrm{raw}}_{\mathfrak G},
\]
the degree-one transport of the canonical effective-descent datum on
$q^*\mathcal I_{Q/S}$ is
\[
    \operatorname{conj}^{\mathrm{raw}}_{\mathfrak G}:
    s^*I^{\mathrm{raw}}_{\mathfrak G}
    \xrightarrow{\simeq}
    t^*I^{\mathrm{raw}}_{\mathfrak G}.
\]
The degree-two cocycle and all higher coherence are the corresponding
simplicial coherences of the \v{C}ech nerve.
\end{theorem}

\begin{proof}
The equivalence
\[
    q^*\mathcal I_{Q/S}
    \simeq
    I^{\mathrm{raw}}_{\mathfrak G}
\]
is Proposition~\ref{prop:isotropy-inertia}.  Over
\[
    \mathfrak G_1=X\times_QX
\]
the defining homotopy of the homotopy pullback gives
\[
    q s\simeq q t.
\]
Pulling the single group object $\mathcal I_{Q/S}\to Q$ along this homotopy
gives the canonical degree-one descent equivalence
\[
    \theta:
    s^*I^{\mathrm{raw}}_{\mathfrak G}
    \xrightarrow{\simeq}
    t^*I^{\mathrm{raw}}_{\mathfrak G}.
\]

It remains to identify this transport.  Let
\[
    a:T\longrightarrow\mathfrak G_1
\]
be an arbitrary generalized point over $S$, and put
\[
    x:=sa,\qquad y:=ta.
\]
The homotopy-pullback datum of
$\mathfrak G_1=X\times_QX$ supplies a path
\[
    g:q(x)\simeq q(y)
\]
in $\Map_S(T,Q)$.  Pull the two candidate transport maps back
along $a$.  Through the standard identification of the fibre of the diagonal
self-intersection
\[
    \mathcal I_{Q/S}
    =
    Q\times_{Q\times_SQ}Q
\]
with the corresponding based-loop object, a generalized element of
$x^*q^*\mathcal I_{Q/S}$ is a loop $h$ at $q(x)$.

Transport in this based-loop family along the path $g$ is, by path induction
in the space $\Map_S(T,Q)$, the change-of-base loop represented
by
\[
    g\circ h\circ g^{-1}.
\]
Under the \v{C}ech-groupoid identification this is exactly the pullback along
$a$ of the morphism constructed from groupoid composition and inversion in
Definition~\ref{def:raw-conjugation}.  Hence
\[
    a^*\theta
    \simeq
    a^*\operatorname{conj}^{\mathrm{raw}}_{\mathfrak G}
\]
naturally in every generalized point $a:T\to\mathfrak G_1$.  Yoneda in the
slice over $\mathfrak G_1$ gives
\[
    \theta
    \simeq
    \operatorname{conj}^{\mathrm{raw}}_{\mathfrak G}.
\]

The degree-two cocycle is functoriality of path transport under concatenation,
which is the composition law of the \v{C}ech groupoid.  The higher simplices
supply the corresponding higher-pasting coherences.  Thus the complete
effective-descent datum on raw inertia is precisely coherent conjugation.
\end{proof}

Now return to
$\mathfrak G_\bullet\in\operatorname{FormLieGpd}(X/S)$ and write
\[
    q:=q_{\mathfrak G}:X\twoheadrightarrow Q_{\mathfrak G},
    \qquad
    \bar\phi:=\bar\phi_{\mathfrak G}:Q_{\mathfrak G}\to Q_{X/S}.
\]

\begin{definition}[Formal adjoint transport]
\label{def:formal-adjoint-transport}
Under the anchor-kernel equivalence
\[
    I^{\mathrm{for}}_{\mathfrak G}
    \simeq
    \mathfrak G_1\times_{\mathcal R^{\mathrm{dR}}_{X/S}}X,
\]
internal conjugation preserves the fibre over the unit section and defines
\[
    \operatorname{Ad}_{\mathfrak G}:
    s^*I^{\mathrm{for}}_{\mathfrak G}
    \longrightarrow
    t^*I^{\mathrm{for}}_{\mathfrak G}.
\]
\end{definition}

\begin{theorem}[Coherent formal adjoint descent]
\label{thm:coherent-adjoint-transport}
The formal adjoint transport is an equivalence and is exactly the degree-one
\v{C}ech transport on
\[
    q^*\mathcal K_{\mathfrak G}
    \simeq
    I^{\mathrm{for}}_{\mathfrak G}.
\]
For composable arrows
\[
    x\xrightarrow{g}y\xrightarrow{h}z
\]
there is a canonical coherent equivalence
\[
    \operatorname{Ad}_{h\circ g}
    \simeq
    \operatorname{Ad}_{h}\operatorname{Ad}_{g},
\]
and all higher cocycle coherence is inherited from effective descent of the
single group object
$\mathcal K_{\mathfrak G}\to Q_{\mathfrak G}$.
\end{theorem}

\begin{proof}
The inertia morphism
\[
    \mathcal I_{Q_{\mathfrak G}/S}
    \longrightarrow
    \bar\phi^*\mathcal I_{Q_{X/S}/S}
\]
and the identity-loop section are morphisms over $Q_{\mathfrak G}$.
Consequently their pullbacks along the \v{C}ech nerve preserve the canonical
descent transport.  The finite-limit fibre defining
$\mathcal K_{\mathfrak G}$ therefore inherits the restriction of the raw
inertia transport.  By
Theorem~\ref{thm:inertia-descent-conjugation} that raw transport is
conjugation.  Under
\[
    q^*\mathcal K_{\mathfrak G}
    \simeq
    I^{\mathrm{for}}_{\mathfrak G},
\]
its restriction is precisely
$\operatorname{Ad}_{\mathfrak G}$.

Since it is descent transport, it is an equivalence.  Its unit, composition,
and all higher coherences are the degree-zero, degree-two, and higher
\v{C}ech coherences of the descended group object.  No independent adjoint
local-system datum is introduced.
\end{proof}

\begin{corollary}[Formal inertia kernel and conjugation descent]
\label{cor:adjoint-inertia-object}
For the formal inertia kernel
\[
    \mathcal K_{\mathfrak G}\longrightarrow Q_{\mathfrak G}
\]
there is a canonical equivalence
\[
    q^*\mathcal K_{\mathfrak G}
    \simeq
    I^{\mathrm{for}}_{\mathfrak G},
\]
and its canonical \v{C}ech transport is the formal adjoint transport.  Raw
isotropy separately satisfies
\[
    I^{\mathrm{raw}}_{\mathfrak G}
    \simeq
    q^*\mathcal I_{Q_{\mathfrak G}/S}.
\]
\end{corollary}

\begin{proof}
The first equivalence is
Theorem~\ref{thm:formal-isotropy-descent}.  The transport statement is
Theorem~\ref{thm:coherent-adjoint-transport}, and the raw statement is
Proposition~\ref{prop:isotropy-inertia}.
\end{proof}

\begin{proposition}[Adjoint covariance of the nonlinear commutator]
\label{prop:adjoint-commutator}
Over $\mathfrak G_1$ there is a canonical equivalence
\[
    \operatorname{Ad}_{\mathfrak G}
    \left(\operatorname{comm}_{\mathfrak G}(a,b)\right)
    \simeq
    \operatorname{comm}_{\mathfrak G}
    \left(
        \operatorname{Ad}_{\mathfrak G}(a),
        \operatorname{Ad}_{\mathfrak G}(b)
    \right).
\]
\end{proposition}

\begin{proof}
Formal adjoint transport is a homomorphism of the corresponding formal
isotropy group objects and therefore preserves multiplication and inversion.
Apply it to $aba^{-1}b^{-1}$.
\end{proof}

\section{Dependent infinitesimal transport: disks, jets, and effective relations}
\label{sec:flg-jet}

Chapter~5 developed the disk--jet calculus for the maximal de Rham quotient
$e_{X/S}:X\twoheadrightarrow Q_{X/S}$.  The same dependent-adjoint mechanism
applies to any relation after replacing its target by its effective image.  We
record the general statement once and then specialize it to formal Lie
groupoids.

\subsection{Effective-image relation calculus}
\label{sec:flg-effective-relation-calculus}

The only extra categorical input needed to pass from an arbitrary relation map
to its effective image is cancellation along a monomorphism.

\begin{lemma}[Dependent-adjoint cancellation along a monomorphism]
\label{lem:monomorphism-dependent-adjoint-cancellation}
Let
\[
    m:I\lhook\joinrel\longrightarrow B
\]
be a monomorphism in $(\mathcal H_{\Sp})_{/S}$, with dependent adjoint triple
\[
    \Sigma_m\dashv m^*\dashv\Pi_m.
\]
Then there are canonical equivalences
\[
    m^*\Sigma_m\simeq\operatorname{id}_{(\mathcal H_{\Sp})_{/I}},
    \qquad
    m^*\Pi_m\simeq\operatorname{id}_{(\mathcal H_{\Sp})_{/I}}.
\]
They are natural in $m$, compatible with Cartesian base change, and with the
units, counits, and mate transformations of the dependent adjunctions.
\end{lemma}

\begin{proof}
For $Z\to I$,
\[
    m^*\Sigma_m Z
    \simeq
    Z\times_B I
    \simeq
    Z\times_I(I\times_BI).
\]
Since $m$ is a monomorphism, its diagonal is an equivalence
\[
    I\xrightarrow{\simeq}I\times_BI,
\]
so
\[
    m^*\Sigma_mZ\simeq Z
\]
naturally in $Z$.

For $Z,E\in(\mathcal H_{\Sp})_{/I}$, the two dependent adjunctions give
\[
\begin{aligned}
    \Map_I(Z,m^*\Pi_mE)
    &\simeq
    \Map_B(\Sigma_mZ,\Pi_mE)\\
    &\simeq
    \Map_I(m^*\Sigma_mZ,E)\\
    &\simeq
    \Map_I(Z,E).
\end{aligned}
\]
Yoneda therefore gives the second equivalence.  Both comparisons are built
from the diagonal equivalence of the monomorphism and the adjunction
equivalences, so their base-change and mate compatibilities are the canonical
ones supplied by Cartesian pasting and functoriality of adjunction.
\end{proof}

\begin{theorem}[Effective-image relation disk--jet calculus]
\label{thm:effective-relation-calculus}
Let
\[
    q:X\longrightarrow B
\]
be an arbitrary morphism in $(\mathcal H_{\Sp})_{/S}$, with effective-image factorization
\[
    X\xrightarrow{e_q}Q_q\xrightarrow{m_q}B,
\]
where $e_q$ is an effective epimorphism and $m_q$ is a monomorphism.  Write
\[
    s,t:X\times_BX\rightrightarrows X
\]
for the two projections, and define
\[
    \mathsf{Disk}^{\infty}_q:=q^*\Sigma_q,
    \qquad
    J^\infty_q:=q^*\Pi_q
\]
as endofunctors of $(\mathcal H_{\Sp})_{/X}$.
Then:
\begin{enumerate}
    \item the kernel-pair nerve depends only on the effective image:
    \[
        \check C_\bullet(q)\simeq\check C_\bullet(e_q),
        \qquad
        \left|\check C_\bullet(q)\right|\simeq Q_q;
    \]
    \item there are canonical equivalences of monads and comonads
    \[
        \mathsf{Disk}^{\infty}_q\simeq e_q^*\Sigma_{e_q},
        \qquad
        J^\infty_q\simeq e_q^*\Pi_{e_q};
    \]
    \item the relation-correspondence formulas are
    \[
        \mathsf{Disk}^{\infty}_q\simeq\Sigma_s t^*
        \simeq\Sigma_t s^*,
        \qquad
        J^\infty_q\simeq\Pi_s t^*;
    \]
    \item there is a canonical adjunction
    \[
        \mathsf{Disk}^{\infty}_q\dashv J^\infty_q;
    \]
    \item a $\mathsf{Disk}^{\infty}_q$-algebra structure, a $J^\infty_q$-coalgebra
    structure, and coherent transport along $\check C_\bullet(q)$ are canonically
    equivalent presentations of the same descent structure;
    \item pullback along $e_q$ induces canonical equivalences
    \[
        (\mathcal H_{\Sp})_{/Q_q}
        \simeq
        \operatorname{Alg}_{\mathsf{Disk}^{\infty}_q}
        \bigl((\mathcal H_{\Sp})_{/X}\bigr)
        \simeq
        \operatorname{Coalg}_{J^\infty_q}
        \bigl((\mathcal H_{\Sp})_{/X}\bigr).
    \]
\end{enumerate}
For every commutative square
\[
\begin{tikzcd}
X \arrow[r,"f"] \arrow[d,"q"'] & X' \arrow[d,"q'"]\\
B \arrow[r,"g"'] & B'
\end{tikzcd}
\]
there are canonical mate transformations
\[
    \mathsf{Disk}^{\infty}_q f^*\longrightarrow f^*\mathsf{Disk}^{\infty}_{q'},
    \qquad
    f^*J^\infty_{q'}\longrightarrow J^\infty_qf^*,
\]
compatible with the monad and comonad structures and with composition.  They
are equivalences when the square is Cartesian.
\end{theorem}

\begin{proof}
Because $m_q$ is a monomorphism,
\[
    Q_q\xrightarrow{\simeq}Q_q\times_BQ_q.
\]
Finite-limit pasting therefore identifies every iterated fibre product over
$B$ with the corresponding iterated fibre product over $Q_q$.  Hence
\[
    \check C_\bullet(q)\simeq\check C_\bullet(e_q),
\]
and effectivity of $e_q$ gives
\[
    \left|\check C_\bullet(q)\right|\simeq Q_q.
\]

Lemma~\ref{lem:monomorphism-dependent-adjoint-cancellation}, applied to
$m_q:Q_q\hookrightarrow B$, gives
\[
    m_q^*\Sigma_{m_q}\simeq\operatorname{id},
    \qquad
    m_q^*\Pi_{m_q}\simeq\operatorname{id}.
\]
Since $q=m_qe_q$, composition of dependent adjunctions gives coherent
identifications
\[
    q^*\Sigma_q\simeq e_q^*\Sigma_{e_q},
    \qquad
    q^*\Pi_q\simeq e_q^*\Pi_{e_q},
\]
including their induced monad and comonad structures.

The relation square
\[
\begin{tikzcd}
X\times_BX \arrow[r,"t"] \arrow[d,"s"'] & X \arrow[d,"q"]\\
X \arrow[r,"q"'] & B
\end{tikzcd}
\]
is Cartesian.  Left and right Beck--Chevalley therefore give
\[
    \mathsf{Disk}^{\infty}_q\simeq\Sigma_s t^*,
    \qquad
    J^\infty_q\simeq\Pi_s t^*,
\]
and relation inversion gives
$\mathsf{Disk}^{\infty}_q\simeq\Sigma_t s^*$.
Composition of
$\Sigma_q\dashv q^*\dashv\Pi_q$ yields the adjunction
\[
    \mathsf{Disk}^{\infty}_q\dashv J^\infty_q.
\]

Now apply the homotopy-coherent adjoint--monad duality proved in Chapter~5:
the complete monad structure on $\mathsf{Disk}^{\infty}_q$ and its iterates is
carried by $(\infty,2)$-categorical mate formation to the complete comonad
structure on $J^\infty_q$.  Consequently there is a canonical equivalence over
$(\mathcal H_{\Sp})_{/X}$
\[
    \operatorname{Alg}_{\mathsf{Disk}^{\infty}_q}
    \bigl((\mathcal H_{\Sp})_{/X}\bigr)
    \simeq
    \operatorname{Coalg}_{J^\infty_q}
    \bigl((\mathcal H_{\Sp})_{/X}\bigr),
\]
and the objectwise action/coaction correspondence carries all unit,
associativity, counit, coassociativity, structured-morphism, and higher-pasting
coherences simultaneously.

Finally $e_q:X\twoheadrightarrow Q_q$ is an effective epimorphism.  Effective
\v{C}ech descent gives
\[
    (\mathcal H_{\Sp})_{/Q_q}
    \simeq
    \operatorname*{Tot}_{[n]\in\Delta}
    (\mathcal H_{\Sp})_{/\check C_n(q)}.
\]
Equivalently, the Barr--Beck argument of Chapter~5 applies verbatim to the
adjunction
$\Sigma_{e_q}\dashv e_q^*$: the pullback $e_q^*$ is conservative by
effective descent and preserves the required realizations because it also has
the right adjoint $\Pi_{e_q}$.  Thus
\[
    (\mathcal H_{\Sp})_{/Q_q}
    \simeq
    \operatorname{Alg}_{\mathsf{Disk}^{\infty}_q}
    \simeq
    \operatorname{Coalg}_{J^\infty_q}.
\]
Taking the homotopy fibre over a fixed degree-zero object identifies these
structures with coherent transport along $\check C_\bullet(q)$.

For a commutative square as in the statement, the canonical left
Beck--Chevalley mate
\[
    \Sigma_q f^*\longrightarrow g^*\Sigma_{q'}
\]
and right Beck--Chevalley mate
\[
    g^*\Pi_{q'}\longrightarrow\Pi_q f^*
\]
become, after pullback, the displayed transformations of disk monads and jet
comonads.  Mate-pasting coherence makes them compatible with composition and
all monad/comonad structure; ordinary Beck--Chevalley makes them equivalences
for Cartesian squares.
\end{proof}


\subsection{The formal-groupoid specialization}

Let
\[
    \mathfrak G_\bullet\in\operatorname{FormLieGpd}(X/S),
    \qquad
    q:=q_{\mathfrak G}:X\twoheadrightarrow Q_{\mathfrak G}.
\]
Since $q$ is already effective, its relation groupoid is
\[
    \mathfrak G_\bullet\simeq\check C_\bullet(q).
\]

\begin{definition}[Formal-groupoid disk and jet endofunctors]
\label{def:formal-groupoid-disk-jet}
Define
\[
    \mathsf{Disk}^{\infty}_{\mathfrak G}:=q^*\Sigma_q,
    \qquad
    J^\infty_{\mathfrak G}:=q^*\Pi_q.
\]
The first carries its canonical monad structure, the second its canonical
comonad structure, and Theorem~\ref{thm:effective-relation-calculus} gives
\[
    \mathsf{Disk}^{\infty}_{\mathfrak G}\dashv J^\infty_{\mathfrak G}.
\]
For a generalized point \(a:T\to X\), its value
\(\mathsf{Disk}^{\infty}_{\mathfrak G}(a)\) is the formal orbit disk of
Definition~\ref{def:formal-orbit-disk}.
\end{definition}

Writing $s,t:\mathfrak G_1\rightrightarrows X$, there are canonical
relation-correspondence equivalences
\[
    \mathsf{Disk}^{\infty}_{\mathfrak G}\simeq\Sigma_s t^*
    \simeq\Sigma_t s^*,
    \qquad
    J^\infty_{\mathfrak G}\simeq\Pi_s t^*.
\]
The unit and counit are induced by the groupoid unit, and multiplication and
comultiplication are induced by groupoid composition.


\begin{proposition}[Dependent infinitesimal-equality formulas]
\label{prop:formal-groupoid-dependent-sum-product}
Let $E\to X$ be read as a dependent family
\[
    x:X\vdash E_x.
\]
In the internal infinitesimal-equality notation, the disk and jet functors have
the fibrewise descriptions
\[
\boxed{
    \bigl(\mathsf{Disk}^{\infty}_{\mathfrak G}E\bigr)_x
    \simeq
    \sum_{y:X}
    \sum_{\gamma:x\sim_{\mathfrak G}y}
    E_y
},
\]
and
\[
\boxed{
    \bigl(J^\infty_{\mathfrak G}E\bigr)_x
    \simeq
    \prod_{y:X}
    \prod_{\gamma:x\sim_{\mathfrak G}y}
    E_y
}.
\]
Accordingly,
\[
    \mathsf{Disk}^{\infty}_{\mathfrak G}\dashv J^\infty_{\mathfrak G}
\]
is the dependent-sum/dependent-product adjunction along the chosen effective
infinitesimal equality relation.
\end{proposition}

\begin{proof}
The source fibre of
\[
    \mathfrak G_1\simeq X\times_{Q_{\mathfrak G}}X
\]
over $x$ consists of pairs $(y,\gamma)$ with
$\gamma:x\sim_{\mathfrak G}y$.  Evaluating
\[
    \mathsf{Disk}^{\infty}_{\mathfrak G}
    \simeq
    \Sigma_s t^*,
    \qquad
    J^\infty_{\mathfrak G}
    \simeq
    \Pi_s t^*
\]
on that fibre gives the two displayed dependent formulas.  Their adjunction
is the already constructed adjunction
$\mathsf{Disk}^{\infty}_{\mathfrak G}\dashv J^\infty_{\mathfrak G}$; the notation only
exposes its dependent logical content.
\end{proof}

\begin{proposition}[Fibrewise formal-disk formula]
\label{prop:formal-groupoid-fibrewise-jets}
Let $a:T\to X$.  For every $E\to X$ there is a canonical equivalence
\[
    a^*J^\infty_{\mathfrak G}E
    \simeq
    \Pi_{\pi^{\mathfrak G}_a}
    \left(\operatorname{ev}^{\mathfrak G}_a\right)^*E.
\]
Thus the generalized fibre of $J^\infty_{\mathfrak G}E$ at $a$ is the object of
sections of $E$ over the formal orbit disk
$\mathsf{Disk}^{\infty}_{\mathfrak G}(a)=T\times_{Q_{\mathfrak G}}X$.
\end{proposition}

\begin{proof}
Pull back $J^\infty_{\mathfrak G}\simeq\Pi_s t^*$ along $a:T\to X$ and apply right
Beck--Chevalley to the Cartesian arrow-fibre square of
Proposition~\ref{prop:formal-disk-arrow-fibre}.
\end{proof}

\subsection{Comparison with the maximal crystalline calculus}

The canonical quotient map
\[
    \bar\phi_{\mathfrak G}:Q_{\mathfrak G}\longrightarrow Q_{X/S},
    \qquad
    e_{X/S}\simeq \bar\phi_{\mathfrak G}q_{\mathfrak G},
\]
compares the chosen formal relation with the maximal de Rham relation.

\begin{proposition}[Ambient disk and jet comparisons]
\label{prop:formal-groupoid-ambient-monad-comonad}
There are canonical natural transformations
\[
    \mathsf{Disk}^{\infty}_{\mathfrak G}
    \longrightarrow
    \mathsf{Disk}^{\mathrm{dR},\infty}_{X/S},
    \qquad
    J^\infty_{X/S}\longrightarrow J^\infty_{\mathfrak G}.
\]
The first is induced by the unit
$\operatorname{id}\to \bar\phi_{\mathfrak G}^*\Sigma_{\bar\phi_{\mathfrak G}}$, and the
second by the counit
$\bar\phi_{\mathfrak G}^*\Pi_{\bar\phi_{\mathfrak G}}\to\operatorname{id}$.  Fibrewise,
they are extension and restriction along the canonical map
\[
    \mathsf{Disk}^{\infty}_{\mathfrak G}(a)
    \longrightarrow
    \operatorname{Nbh}^{\mathrm{dR}}_{T/S}(T\times_SX).
\]
\end{proposition}

\begin{proof}
Use $e_{X/S}=\bar\phi_{\mathfrak G}q_{\mathfrak G}$ and compose the dependent
adjunctions.  The two transformations are the corresponding unit and counit
mates.  Their fibrewise description follows from the formal-disk formulas of
Chapter~5 and Proposition~\ref{prop:formal-groupoid-fibrewise-jets}.
\end{proof}

\section{External infinitesimal comparisons and effective images}
\label{sec:flg-full-effective}

The intrinsic theory requires no target beyond the effective formal quotient.
Nevertheless a later construction may produce a comparison
\[
    X\xrightarrow{a}\mathfrak X\xrightarrow{b}X_{\mathrm{dR}/S}
\]
together with a specified homotopy
$ba\simeq\eta_{X/S}^{\mathrm{dR}}$.  Such a diagram is external input to the
present chapter; no representability, effectivity, or locality of
$\mathfrak X$ is inferred from it.  The only intrinsic operation we perform
on this input is its effective image.

Factor
\[
    X\xrightarrow{e_a}Q_a\xrightarrow{m_a}\mathfrak X
\]
by the effective-epimorphism--monomorphism factorization.

\begin{proposition}[Effective-image invariance of an external comparison]
\label{prop:kernel-pair-image-invariance}
For every such comparison there are canonical simplicial equivalences
\[
    \underbrace{
    X\times_{\mathfrak X}\cdots\times_{\mathfrak X}X
    }_{n+1}
    \simeq
    \underbrace{
    X\times_{Q_a}\cdots\times_{Q_a}X
    }_{n+1}
\]
for all $n\geq0$.  Hence
\[
    \check C_\bullet(X\to\mathfrak X)
    \simeq
    \check C_\bullet(e_a).
\]
The homotopy $ba\simeq\eta_{X/S}^{\mathrm{dR}}$ induces canonically a map
\[
    Q_a\longrightarrow Q_{X/S}
\]
under $X$.  Thus the kernel-pair groupoid is canonically anchored; its
unitwise de Rham formalization is a formal Lie groupoid; and the relation disk--jet
calculus attached to $a$ depends only on the effective image $e_a$.
\end{proposition}

\begin{proof}
Since $m_a$ is a monomorphism,
\[
    Q_a\xrightarrow{\simeq}Q_a\times_{\mathfrak X}Q_a.
\]
Finite-limit pasting gives
\[
    X\times_{\mathfrak X}X
    \simeq
    X\times_{Q_a}X,
\]
and the same argument in every degree gives the simplicial equivalence.

The square of arrows from $a$ to the de Rham unit, determined by $b$ and the
specified homotopy, is carried by functorial effective-image factorization to
a morphism from the effective image of $a$ to the effective image of the de
Rham unit.  This is the displayed map $Q_a\to Q_{X/S}$.  The quotient
classification produces the corresponding anchor, and unitwise formalization
produces the formal fixed point.  Finally
Theorem~\ref{thm:effective-relation-calculus} identifies the relation monad,
comonad, and transport calculus of $a$ with those of $e_a$.
\end{proof}

\begin{remark}[Effective image versus an external target]
The kernel-pair geometry sees $Q_a$, not arbitrary excess information in
$\mathfrak X$.  Formalization may further replace $Q_a$ by the coreflective
formal quotient $Q_{\ell_{e_a}}$.  Thus
\[
    \mathfrak X,
    \qquad
    Q_a,
    \qquad
    Q_{\ell_{e_a}}
\]
may be different objects.  Any later comparison with a represented or
formal-moduli leaf space must prove which one is modeled; no identification is
built into the definition of a formal Lie groupoid.
\end{remark}

\section{Crystalline descent and dependent infinitesimal transport}
\label{sec:flg-descent}

\begin{theorem}[Formal-groupoid transport and mate compatibility]
\label{thm:formal-groupoid-transport-mates}
For $E\to X$, the following data are canonically equivalent:
\begin{enumerate}
    \item a $\mathsf{Disk}^{\infty}_{\mathfrak G}$-algebra structure on $E$;
    \item a $J^\infty_{\mathfrak G}$-coalgebra structure on $E$;
    \item coherent descent transport along $\mathfrak G_\bullet$, whose
    degree-one term is an equivalence
    \[
        \theta:s^*E\xrightarrow{\simeq}t^*E
    \]
    satisfying
    \[
        u^*\theta\simeq\operatorname{id}_E
    \]
    and the \v{C}ech cocycle identity on $\mathfrak G_2$.
\end{enumerate}
Under these equivalences the algebra action and coalgebra coaction are mates
under
$\mathsf{Disk}^{\infty}_{\mathfrak G}\dashv J^\infty_{\mathfrak G}$.
\end{theorem}

\begin{proof}
By groupoid effectivity,
\[
    \mathfrak G_\bullet
    \simeq
    \check C_\bullet(q).
\]
Apply Theorem~\ref{thm:effective-relation-calculus} to
$q:X\twoheadrightarrow Q_{\mathfrak G}$.  Its relation groupoid is
$\mathfrak G_\bullet$, and its monad and comonad are precisely
\[
    \mathsf{Disk}^{\infty}_{\mathfrak G}=q^*\Sigma_q,
    \qquad
    J^\infty_{\mathfrak G}=q^*\Pi_q.
\]
The theorem therefore gives the stated equivalence between algebra structure,
coalgebra structure, and coherent transport, including the unit, cocycle, and
all higher \v{C}ech coherences.  It also identifies the algebra action and
coalgebra coaction as mates under
$\mathsf{Disk}^{\infty}_{\mathfrak G}\dashv J^\infty_{\mathfrak G}$.
\end{proof}

\begin{theorem}[Effective formal-groupoid descent]
\label{thm:formal-groupoid-effective-descent}
Pullback along $q:X\twoheadrightarrow Q_{\mathfrak G}$ induces canonical
equivalences
\[
    (\mathcal H_{\Sp})_{/Q_{\mathfrak G}}
    \simeq
    \operatorname{Alg}_{\mathsf{Disk}^{\infty}_{\mathfrak G}}((\mathcal H_{\Sp})_{/X})
    \simeq
    \operatorname{Coalg}_{J^\infty_{\mathfrak G}}((\mathcal H_{\Sp})_{/X}).
\]
Both equivalences commute with the forgetful functors to
$(\mathcal H_{\Sp})_{/X}$.
\end{theorem}

\begin{proof}
Apply Theorem~\ref{thm:effective-relation-calculus} to the effective quotient
\[
    q:X\twoheadrightarrow Q_{\mathfrak G}.
\]
Groupoid effectivity identifies its \v{C}ech nerve with
$\mathfrak G_\bullet$, so the theorem gives simultaneously
\[
    (\mathcal H_{\Sp})_{/Q_{\mathfrak G}}
    \simeq
    \operatorname*{Tot}_{[n]\in\Delta}
    (\mathcal H_{\Sp})_{/\mathfrak G_n}
    \simeq
    \operatorname{Alg}_{\mathsf{Disk}^{\infty}_{\mathfrak G}}((\mathcal H_{\Sp})_{/X})
    \simeq
    \operatorname{Coalg}_{J^\infty_{\mathfrak G}}((\mathcal H_{\Sp})_{/X}),
\]
with the common underlying functor given by pullback along $q$.
\end{proof}

\begin{definition}[Descended formal-groupoid objects]
\label{def:formal-cartan-objects}
A \emph{descended $\mathfrak G$-object} is an object of the quotient slice
\[
    (\mathcal H_{\Sp})_{/Q_{\mathfrak G}}.
\]
No additional category symbol is introduced.  Through
Theorem~\ref{thm:formal-groupoid-effective-descent}, the same objects may
equivalently be presented as
\[
    \operatorname{Alg}_{\mathsf{Disk}^{\infty}_{\mathfrak G}}((\mathcal H_{\Sp})_{/X})
    \quad\text{or}\quad
    \operatorname{Coalg}_{J^\infty_{\mathfrak G}}((\mathcal H_{\Sp})_{/X}).
\]
\end{definition}

\begin{remark}[Crystalline transport is derived]
A descended formal-groupoid object is intrinsically an object over the formal leaf space.
Its disk action and jet coaction are derived by effective descent and are not
independent structures.  This is the exact analogue, for a chosen formal Lie
groupoid, of the intrinsic crystalline-descent philosophy of the maximal jet chapter.
\end{remark}


\begin{proposition}[Crystalline transport as dependent substitution]
\label{prop:cartan-dependent-substitution}
Let
\[
    \bar E\longrightarrow Q_{\mathfrak G}
\]
be a descended formal-groupoid object and put $E:=q^*\bar E$.  In the internal
infinitesimal-equality notation, every witness
\[
    \gamma:x\sim_{\mathfrak G}y
\]
induces a canonical equivalence
\[
    \gamma_*:E_x\xrightarrow{\simeq}E_y.
\]
These equivalences satisfy reflexivity, composition, inversion, and all higher
coherence supplied by substitution in the single family $\bar E$ over
$Q_{\mathfrak G}$.  Under
Theorem~\ref{thm:formal-groupoid-effective-descent}, this substitution is
exactly the coherent \v{C}ech transport
\[
    s^*E\xrightarrow{\simeq}t^*E
\]
and hence exactly the
$\mathsf{Disk}^{\infty}_{\mathfrak G}$-algebra and
$J^\infty_{\mathfrak G}$-coalgebra transport.
\end{proposition}

\begin{proof}
The relation
\[
    \mathfrak G_1
    \simeq
    X\times_{Q_{\mathfrak G}}X
\]
carries its canonical homotopy
$q s\simeq q t$.  Pulling the single family
$\bar E\to Q_{\mathfrak G}$ along this homotopy gives
\[
    s^*q^*\bar E
    \xrightarrow{\simeq}
    t^*q^*\bar E.
\]
In internal notation this is substitution along
$\gamma:q(x)=q(y)$, written
$\gamma:x\sim_{\mathfrak G}y$.  The higher substitution laws are exactly the
higher \v{C}ech descent coherences.  The mate comparison follows from
Theorem~\ref{thm:formal-groupoid-transport-mates}.
\end{proof}

\subsection{Invariant and descended sections}

Let $\bar E\to Q_{\mathfrak G}$ and set $E:=q^*\bar E$.

\begin{definition}[Coherent $\mathfrak G$-invariant sections]
\label{def:formal-groupoid-invariant-sections}
The space of coherent $\mathfrak G$-invariant sections of $E$ is the mapping
space
\[
    \Map_{(\mathcal H_{\Sp})_{/Q_{\mathfrak G}}}
    \left(
        \operatorname{id}_{Q_{\mathfrak G}},
        \bar E
    \right).
\]
No separate section-space notation is introduced.
\end{definition}

\begin{proposition}[Cech presentation of invariant sections]
\label{prop:invariant-sections-totalization}
There is a canonical equivalence
\[
    \Map_{(\mathcal H_{\Sp})_{/Q_{\mathfrak G}}}
    \left(
        \operatorname{id}_{Q_{\mathfrak G}},
        \bar E
    \right)
    \simeq
    \operatorname*{Tot}_{[n]\in\Delta}
    \Map_{(\mathcal H_{\Sp})_{/\mathfrak G_n}}
    \left(
        \operatorname{id}_{\mathfrak G_n},
        E|_{\mathfrak G_n}
    \right),
\]
where, for the zeroth vertex map
$p_0:\mathfrak G_n\to X$, we write
\[
    E|_{\mathfrak G_n}:=p_0^*E.
\]
The coherent descent transport canonically identifies this object with the
pullback along every other vertex and induces the displayed cosimplicial
structure.
\end{proposition}

\begin{proof}
Apply effective descent to mapping spaces in the slice over
$Q_{\mathfrak G}$.  Mapping spaces in a limit of $\infty$-categories are the
corresponding limits of mapping spaces.
\end{proof}

\subsection{Interface with the stable differential-cohomological reconstruction}

The crystalline descent theory above is nonlinear: its objects lie in the slice
\[
    (\mathcal H_{\Sp})_{/Q_{\mathfrak G}}.
\]
Chapter~2 separately stabilized the differential-cohesive endpoint semantics
and proved the BNV-type exact-adjoint fracture for stable infinitesimal
coefficients.  These two levels must not be conflated.  In particular, a
nonlinear descended object is not automatically equipped with a cycle sector,
curvature morphism, or differential-cohomology class.

Nevertheless the present descent construction provides the nonlinear input
to which that stable reconstruction may later be applied.  Once a later
coefficient construction produces a stable infinitesimal coefficient
\[
    E\in\mathcal D_{\inf}
\]
from a pointed or linearized descended object, Chapter~2 supplies canonically
\[
    \mathcal A(E):=\operatorname{fib}(E\to\Im^{\mathrm{st}}E),
    \qquad
    \mathcal Z(E):=\operatorname{cofib}(\&^{\mathrm{st}}E\to E),
\]
the adjacent-mates characteristic morphism, and the reconstruction
equivalence
\[
    E
    \simeq
    \mathcal Z(E)
    \times_{\Im^{\mathrm{st}}\mathcal Z(E)}
    \Im^{\mathrm{st}}E.
\]
Thus the nonlinear sequence of the present chapter and the stable
differential-cohomological sequence of Chapter~2 fit as
\[
\boxed{
\begin{aligned}
\text{formal infinitesimal quotient}
&\longrightarrow
\text{crystalline descent and infinitesimal transport}
\\
&\longrightarrow
\text{pointed or linearized stable coefficient}
\\
&\longrightarrow
\text{BNV deformation, cycle, and characteristic sectors}.
\end{aligned}}
\]

\begin{remark}[The relative stable step is a theorem domain]
The global stabilized endpoint adjunctions of Chapter~2 do not by themselves
identify the stabilization of every slice
$(\mathcal H_{\Sp})_{/Q_{\mathfrak G}}$ with a quasi-coherent
coefficient category, nor do they automatically supply a sliced BNV fracture
over an arbitrary formal leaf object.  Such a relative identification must be
constructed on its stated theorem domain.  The present chapter therefore uses
the BNV reconstruction as a precise coefficient-level handoff and does not
insert unproved ``differential-crystalline'' data into
the quotient slice $(\mathcal H_{\Sp})_{/Q_{\mathfrak G}}$.
\end{remark}


\section{Functoriality and base change}
\label{sec:flg-functoriality}

\subsection{Morphisms over a fixed object}

Let
\[
    F:\mathfrak G_\bullet\longrightarrow\mathfrak H_\bullet
\]
be a morphism of formal fixed-object groupoids on $X$.  Since
$\mathcal R^{\mathrm{dR}}_{X/S,\bullet}$ is terminal among formal fixed-object groupoids,
\[
    \phi_{\mathfrak H}F
    \simeq
    \phi_{\mathfrak G}
\]
through a contractible space of compatible homotopies.  Realization therefore gives a map
\[
    a:=Q_F:
    Q_{\mathfrak G}\longrightarrow Q_{\mathfrak H}
\]
under $X$ and over $Q_{X/S}$, with
\[
    q_{\mathfrak H}\simeq aq_{\mathfrak G}.
\]

\begin{proposition}[Functoriality over fixed $X$]
\label{prop:formal-gpd-functoriality-fixed-X}
A morphism
$F:\mathfrak G\to\mathfrak H$
induces canonically:
\begin{enumerate}
    \item for every $b:T\to X$, a morphism of formal orbit disks
    \[
        \mathsf{Disk}^{\infty}_{\mathfrak G}(b)
        \longrightarrow
        \mathsf{Disk}^{\infty}_{\mathfrak H}(b);
    \]
    \item homomorphisms
    \[
        I^{\mathrm{raw}}_{\mathfrak G}
        \longrightarrow
        I^{\mathrm{raw}}_{\mathfrak H},
        \qquad
        I^{\mathrm{for}}_{\mathfrak G}
        \longrightarrow
        I^{\mathrm{for}}_{\mathfrak H};
    \]
    \item over $Q_{\mathfrak G}$, homomorphisms
    \[
        \mathcal I_{Q_{\mathfrak G}/S}
        \longrightarrow
        a^*\mathcal I_{Q_{\mathfrak H}/S},
        \qquad
        \mathcal K_{\mathfrak G}
        \longrightarrow
        a^*\mathcal K_{\mathfrak H};
    \]
    \item a natural transformation of monads
    \[
        \mathsf{Disk}^{\infty}_{\mathfrak G}
        \longrightarrow
        \mathsf{Disk}^{\infty}_{\mathfrak H};
    \]
    \item a natural transformation of comonads
    \[
        J^\infty_{\mathfrak H}
        \longrightarrow
        J^\infty_{\mathfrak G}.
    \]
\end{enumerate}
These constructions preserve identities, composition, commutators,
adjoint transport, bitorsor structures, and all higher functorial
coherences.
\end{proposition}

\begin{proof}
The quotient map $a$ induces
\[
    T\times_{Q_{\mathfrak G}}X
    \longrightarrow
    T\times_{Q_{\mathfrak H}}X,
\]
giving the disk map.  Raw isotropy is a finite-limit construction on
$F_1$, and formal isotropy is its centred formalization, so
Proposition~\ref{prop:centred-formalization-varying-centre} gives the two
isotropy homomorphisms.  Equivalently, formal isotropy functoriality follows
from the anchor-kernel formula.  Relative inertia is functorial in the
quotient, and the square over $Q_{X/S}$ carries the identity-loop fibre
defining $\mathcal K_{\mathfrak G}$ to that defining
$\mathcal K_{\mathfrak H}$.

For the disk monads,
\[
\begin{aligned}
    \mathsf{Disk}^{\infty}_{\mathfrak H}
    &=
    q_{\mathfrak H}^*\Sigma_{q_{\mathfrak H}}\\
    &\simeq
    q_{\mathfrak G}^*a^*\Sigma_a\Sigma_{q_{\mathfrak G}}.
\end{aligned}
\]
Apply
$q_{\mathfrak G}^*(-)\Sigma_{q_{\mathfrak G}}$
to the unit
\[
    \operatorname{id}\longrightarrow a^*\Sigma_a
\]
of $\Sigma_a\dashv a^*$.

For jets,
\[
\begin{aligned}
    J^\infty_{\mathfrak H}
    &=
    q_{\mathfrak H}^*\Pi_{q_{\mathfrak H}}\\
    &\simeq
    q_{\mathfrak G}^*a^*\Pi_a\Pi_{q_{\mathfrak G}}.
\end{aligned}
\]
Apply
$q_{\mathfrak G}^*(-)\Pi_{q_{\mathfrak G}}$
to the counit
\[
    a^*\Pi_a\longrightarrow\operatorname{id}
\]
of $a^*\dashv\Pi_a$.  Compatibility with the monad and comonad structures
is naturality of units, counits, and composition of dependent adjunctions.
The remaining structures are formed functorially from groupoid
multiplication, inversion, finite limits, and the canonical anchor.
\end{proof}

\subsection{Functoriality in the geometric object}

\begin{definition}[Functor of formal Lie groupoids]
\label{def:functor-formal-lie-groupoids}
Let $f:X\to Y$ be a morphism over $S$.  A functor
\[
    (\mathfrak G,X)\longrightarrow(\mathfrak H,Y)
\]
of formal Lie groupoids over $S$ is a simplicial functor
\[
    F_\bullet:
    \mathfrak G_\bullet\longrightarrow\mathfrak H_\bullet
\]
whose degree-zero component is $f$.
\end{definition}

\begin{proposition}[Automatic anchor compatibility]
\label{prop:automatic-anchor-functoriality}
Every functor of formal Lie groupoids as above fits canonically into a
homotopy-commutative square
\[
\begin{tikzcd}
\mathfrak G_\bullet \arrow[r,"F_\bullet"] \arrow[d,"\phi_{\mathfrak G}"'] &
\mathfrak H_\bullet \arrow[d,"\phi_{\mathfrak H}"]\\
\mathcal R^{\mathrm{dR}}_{X/S,\bullet} \arrow[r] &
\mathcal R^{\mathrm{dR}}_{Y/S,\bullet}.
\end{tikzcd}
\]
The construction supplies a canonical compatible homotopy, functorial in
$(F,f)$ and coherently compatible with identities and composition.
\end{proposition}

\begin{proof}
For every fixed-object groupoid there is a canonical vertex morphism to the
relative pair groupoid.  Hence the square
\[
\begin{tikzcd}
\mathfrak G_\bullet \arrow[r,"F_\bullet"] \arrow[d] &
\mathfrak H_\bullet \arrow[d]\\
\operatorname{Pair}(X/S)_\bullet
\arrow[r,"\operatorname{Pair}(f)"'] &
\operatorname{Pair}(Y/S)_\bullet
\end{tikzcd}
\]
commutes by functoriality of the vertex maps.  Apply the varying-centre
formalization of Proposition~\ref{prop:centred-formalization-varying-centre}.
Using
Theorem~\ref{thm:pair-groupoid-formalization} identifies the bottom formalized
map with the canonical simplicial map induced by $f$ and the vertical formalized maps with the canonical anchors.
\end{proof}

\begin{proposition}[Quotient functoriality]
\label{prop:quotient-functoriality}
A functor
$(\mathfrak G,X)\to(\mathfrak H,Y)$
induces a canonical morphism
\[
    Q_{\mathfrak G}\longrightarrow Q_{\mathfrak H}
\]
compatible with
\[
    Q_{X/S}\longrightarrow Q_{Y/S}.
\]
The assignment preserves identities, composition, and all higher
functorial coherences.
\end{proposition}

\begin{proof}
Apply geometric realization to the simplicial functor and to the canonical
anchor square of
Proposition~\ref{prop:automatic-anchor-functoriality}.
\end{proof}

\subsection{Arbitrary change of base}

Let $c:S'\to S$, put
\[
    X':=X\times_SS',
    \qquad
    \mathfrak G'_\bullet
    :=
    \mathfrak G_\bullet\times_SS'.
\]

\begin{theorem}[Base change of formal Lie groupoids]
\label{thm:formal-gpd-base-change}
There are canonical equivalences
\[
    Q_{\mathfrak G'}
    \simeq
    Q_{\mathfrak G}\times_SS',
    \qquad
    \mathcal R^{\mathrm{dR}}_{X'/S',\bullet}
    \simeq
    \mathcal R^{\mathrm{dR}}_{X/S,\bullet}\times_SS',
\]
under which $\mathfrak G'_\bullet$ is a formal Lie groupoid with the
base-changed canonical anchor.  Moreover,
\[
    P_{X'}^{\mathrm{dR}}(\mathfrak G'_\bullet)
    \simeq
    P_X^{\mathrm{dR}}(\mathfrak G_\bullet)\times_SS',
\]
\[
    I^{\mathrm{raw}}_{\mathfrak G'}
    \simeq
    I^{\mathrm{raw}}_{\mathfrak G}\times_SS',
    \qquad
    I^{\mathrm{for}}_{\mathfrak G'}
    \simeq
    I^{\mathrm{for}}_{\mathfrak G}\times_SS',
\]
and
\[
    \mathcal K_{\mathfrak G'}
    \simeq
    \mathcal K_{\mathfrak G}\times_SS'.
\]
For every generalized point $a:T\to X$ and its base change
$a':T':=T\times_SS'\to X'$,
\[
    \mathsf{Disk}^{\infty}_{\mathfrak G'}(a')
    \simeq
    \mathsf{Disk}^{\infty}_{\mathfrak G}(a)\times_SS'.
\]
The quotient formalization, disk monads, jet comonads, adjunction, transport
data, bitorsor structures, adjoint transport, and effective descent
equivalences commute with base change through their canonical
Beck--Chevalley maps.
\end{theorem}

\begin{proof}
Pullback between slices of an $\infty$-topos preserves finite limits and all
small colimits.  It therefore preserves groupoid realizations, effective
epimorphisms, monomorphisms, and the iterated fibre products of \v{C}ech nerves.

The relative de Rham base-change theorem and
Proposition~\ref{prop:centred-formalization-base-change} give
\[
    P_{X'}^{\mathrm{dR}}(\mathfrak G'_\bullet)
    \simeq
    P_X^{\mathrm{dR}}(\mathfrak G_\bullet)\times_SS'.
\]
Thus formality is preserved.  Base change for $\mathcal R^{\mathrm{dR}}_{X/S,\bullet}$ follows
either from the base-change theorem of Chapter~4 or from
Theorem~\ref{thm:pair-groupoid-formalization}.  Geometric realization gives
the quotient formula.

Raw isotropy, relative inertia, formal disks, and the two arrow images are
finite-limit or effective-image constructions, hence commute with pullback.
Formal isotropy commutes with base change by centred-formalization base change,
and the formula for $\mathcal K$ follows from finite limits and the
base-changed quotient map.  Effective-image factorization also gives
\[
    Q_{\ell_{q'}}
    \simeq
    Q_{\ell_q}\times_SS'
\]
for the quotient formalization.  The dependent-adjoint and descent
statements are the canonical Beck--Chevalley and effective-descent
comparisons.
\end{proof}


\section{Fixed-object quotient invariance and atlas dependence}
\label{sec:flg-presentation-invariance}

The quotient classification gives complete presentation invariance while the
object space $X$ is fixed.  It does not give arbitrary-atlas Morita
invariance, because the formality condition itself depends on the chosen
centre.

\begin{definition}[Fixed-object quotient equivalence]
\label{def:quotient-equivalence}
A morphism
\[
    F:\mathfrak G\to\mathfrak H
\]
in $\operatorname{FormLieGpd}(X/S)$ is a
\emph{fixed-object quotient equivalence} if
\[
    Q_F:
    Q_{\mathfrak G}\longrightarrow Q_{\mathfrak H}
\]
is an equivalence.
\end{definition}

\begin{proposition}[Fixed-object presentation invariance]
\label{prop:presentation-invariance}
A fixed-object quotient equivalence is an equivalence of formal Lie
groupoids.  It induces equivalences of the formal orbit disks, raw and formal
isotropy groups, formal inertia kernels, raw and formal arrow images, adjoint
transport, bitorsors, disk monads, jet comonads, and descent categories.
\end{proposition}

\begin{proof}
By groupoid effectivity,
\[
    \mathfrak G_\bullet
    \simeq
    \check C_\bullet(X\to Q_{\mathfrak G}),
    \qquad
    \mathfrak H_\bullet
    \simeq
    \check C_\bullet(X\to Q_{\mathfrak H}).
\]
An equivalence of quotient objects under $X$ gives degreewise equivalences of
the \v{C}ech nerves.  Every remaining construction is obtained functorially from
the quotient, its relation groupoid, finite limits, effective images, centred
formalization, or dependent adjunction.
\end{proof}

\begin{proposition}[Unitwise formality and formalization are not arbitrary-atlas invariant]
\label{prop:formality-not-morita-invariant}
Unitwise de Rham formality is not preserved by arbitrary
replacement of the object space by another effective atlas of the same
quotient.  Moreover the canonical unitwise formalization constructed in this
chapter does not restore arbitrary-atlas Morita invariance.
\end{proposition}

\begin{proof}
Let $X$ be noninitial and start with the zero formal groupoid
\[
    \check C_\bullet(\operatorname{id}_X).
\]
Change atlas by the split twofold effective cover
\[
    u:
    U:=X\amalg X
    \twoheadrightarrow
    X.
\]
The resulting groupoid is $\check C_\bullet(u)$.

By Theorem~\ref{thm:formal-quotient-monomorphism-recognition}, this groupoid is
formal if and only if
\[
    \operatorname{dR}_S(u):
    \operatorname{dR}_S(U)\longrightarrow \operatorname{dR}_S(X)
\]
is a monomorphism.  Chapter~4 proves finite-coproduct
compatibility, so
\[
    \operatorname{dR}_S(U)
    \simeq
    \operatorname{dR}_S(X)\amalg \operatorname{dR}_S(X),
\]
and $\operatorname{dR}_S(u)$ is the fold map.  We claim that $\operatorname{dR}_S(X)$ is noninitial.  If
$\operatorname{dR}_S(X)\simeq\varnothing$, then the relative de Rham unit would supply a
morphism
\[
    X\longrightarrow \operatorname{dR}_S(X)\simeq\varnothing.
\]
The initial object of an $\infty$-topos is strict, so this would force
$X\simeq\varnothing$, contrary to the choice of $X$.  Thus $\operatorname{dR}_S(X)$ is
noninitial.

Coproducts in an $\infty$-topos are disjoint.  The self-pullback of the fold
map
\[
    \operatorname{dR}_S(X)\amalg \operatorname{dR}_S(X)
    \longrightarrow
    \operatorname{dR}_S(X)
\]
contains the two off-diagonal coproduct summands, and these are noninitial.
Its diagonal is therefore not an equivalence.  Hence the fold map is not
monic and $\check C_\bullet(u)$ is not formal.

We now compute its canonical formalization.  The centred quotient formalization
of $u:U\twoheadrightarrow X$ is
\[
\begin{aligned}
    P^{\mathrm{dR}}(u)
    &=
    X\times_{\operatorname{dR}_S(X)}\operatorname{dR}_S(U)\\
    &\simeq
    X\times_{\operatorname{dR}_S(X)}
    \bigl(\operatorname{dR}_S(X)\amalg \operatorname{dR}_S(X)\bigr)\\
    &\simeq
    \bigl(X\times_{\operatorname{dR}_S(X)}\operatorname{dR}_S(X)\bigr)
    \amalg
    \bigl(X\times_{\operatorname{dR}_S(X)}\operatorname{dR}_S(X)\bigr)\\
    &\simeq
    X\amalg X
    =
    U.
\end{aligned}
\]
Under this equivalence, the formalized centre
\[
    \ell_u=(u,\eta_{U/S}^{\mathrm{dR}}):U\longrightarrow P^{\mathrm{dR}}(u)
\]
restricts to the identity on each coproduct component and is therefore an
equivalence.  Its effective image factorization is consequently
\[
    U
    \xrightarrow{\operatorname{id}_U}
    U
    \xrightarrow{\simeq}
    P^{\mathrm{dR}}(u).
\]
Thus
\[
    Q_{\ell_u}\simeq U,
    \qquad
    e_{\ell_u}\simeq\operatorname{id}_U,
    \qquad
    c_u\simeq u.
\]
The quotient formula for unitwise formalization now gives
\[
    P_U^{\mathrm{dR}}\bigl(\check C_\bullet(u)\bigr)
    \simeq
    \check C_\bullet(\operatorname{id}_U)
    =
    0_{U,\bullet}.
\]
Its quotient is $U$, not $X$.  Hence the chapter's canonical formalization,
although a coreflection at fixed object space, does not turn this atlas
replacement back into the original zero groupoid on $X$.
\end{proof}

\begin{remark}[Scope of presentation invariance]
The constructions are presentation-independent under equivalences of quotient
presentations with fixed object space $X$.  No stronger arbitrary-atlas
invariance is asserted.  Unitwise formalization is itself a fixed-centre
coreflection and, as the preceding calculation shows, is not a Morita
replacement procedure.  If a later recognition theorem uses another atlas,
its comparison with the present formal groupoid must therefore construct and
prove the relevant comparison rather than identify the two presentations
solely from a common quotient.
\end{remark}


\section{Basic geometric examples}
\label{sec:flg-examples}

\subsection{The zero formal Lie groupoid}

\begin{example}[Zero groupoid]
The identity groupoid
\[
    0_{X,\bullet}
    =
    \check C_\bullet(\operatorname{id}_X)
\]
is a formal Lie groupoid with quotient $X$.  For every $a:T\to X$,
\[
    \mathsf{Disk}^{\infty}_{0_X}(a)\simeq T,
    \qquad
    \mathsf{Disk}^{\infty}_{0_X}\simeq\operatorname{id},
    \qquad
    J^\infty_{0_X}\simeq\operatorname{id}.
\]
Its raw isotropy is the relative inertia of $X$,
\[
    I^{\mathrm{raw}}_{0_X}
    \simeq
    \mathcal I_{X/S},
\]
while its formal isotropy is
\[
    I^{\mathrm{for}}_{0_X}
    \simeq
    \mathcal I_{X/S}
    \times_{e_{X/S}^*\mathcal I_{Q_{X/S}/S}}
    X.
\]
\end{example}

\begin{proof}
The quotient is $\operatorname{id}_X$, so the orbit disk and the two
dependent-adjoint composites are identities.  The raw isotropy formula is
Proposition~\ref{prop:isotropy-inertia}.  Here the canonical quotient anchor is
$e_{X/S}:X\to Q_{X/S}$, so
Theorem~\ref{thm:formal-isotropy-descent} gives the displayed formal
isotropy kernel.
\end{proof}

\begin{remark}
The raw isotropy of the zero groupoid need not be trivial in an
$\infty$-topos: it records the ambient relative inertia of the object $X$.
The term ``zero groupoid'' refers to the simplicial relation structure, not
to the vanishing of ambient higher loops.
\end{remark}

\subsection{The maximal relative de Rham groupoid}

\begin{example}[Maximal groupoid]
The terminal formal Lie groupoid is
\[
    \mathcal R^{\mathrm{dR}}_{X/S,\bullet}.
\]
Its formal leaf space is $Q_{X/S}$, its formal orbit disks are
\[
    \operatorname{Nbh}^{\mathrm{dR}}_{T/S}(T\times_SX),
\]
and
\[
    \mathsf{Disk}^{\infty}_{\mathcal R^{\mathrm{dR}}_{X/S}}=\mathsf{Disk}^{\mathrm{dR},\infty}_{X/S},
    \qquad
    J^\infty_{\mathcal R^{\mathrm{dR}}_{X/S}}=J^\infty_{X/S}.
\]
Its raw isotropy is
\[
    I^{\mathrm{raw}}_{\mathcal R^{\mathrm{dR}}_{X/S}}
    \simeq
    e_{X/S}^*\mathcal I_{Q_{X/S}/S},
\]
whereas its formal isotropy is the trivial group object
\[
    I^{\mathrm{for}}_{\mathcal R^{\mathrm{dR}}_{X/S}}
    \simeq X.
\]
\end{example}

\begin{proof}
The quotient and jet assertions are immediate from the definitions.  The raw
isotropy formula is
Proposition~\ref{prop:isotropy-inertia}.  The canonical anchor of the maximal
relation is the identity, so the anchor-kernel formula gives
\[
    I^{\mathrm{for}}_{\mathcal R^{\mathrm{dR}}_{X/S}}
    \simeq
    \mathcal R^{\mathrm{dR}}_{X/S}\times_{\mathcal R^{\mathrm{dR}}_{X/S}}X
    \simeq X.
\]
\end{proof}

\subsection{Formalization of the pair groupoid}

\begin{example}[Pair groupoid]
For the relative pair groupoid
\[
    \operatorname{Pair}(X/S)_n
    =
    X^{\times_S(n+1)},
\]
Theorem~\ref{thm:pair-groupoid-formalization} gives
\[
    P_X^{\mathrm{dR}}\bigl(\operatorname{Pair}(X/S)_\bullet\bigr)
    \simeq
    \mathcal R^{\mathrm{dR}}_{X/S,\bullet}.
\]
Thus the maximal infinitesimal relation is precisely the unitwise
de Rham formalization of the codiscrete relative groupoid.
\end{example}

\subsection{Formal action groupoids}

Let $G\to S$ be a group object and define its formal identity fibre by
\[
    \operatorname{Nbh}^{\mathrm{dR}}_{S/S}(G)
    =
    G\times_{G_{\mathrm{dR}/S}}S,
\]
where $S\to G_{\mathrm{dR}/S}$ is the relative de Rham image of the identity
section.  For commutative group objects this is the translation object already
appearing in Chapter~5; here no commutativity is assumed.

\begin{proposition}[Constructed formal identity group]
\label{prop:formal-identity-group}
The object $\operatorname{Nbh}^{\mathrm{dR}}_{S/S}(G)$ carries a canonical group-object
structure over $S$.  The projection
\[
    \operatorname{Nbh}^{\mathrm{dR}}_{S/S}(G)\longrightarrow G
\]
is a group homomorphism, $\operatorname{Nbh}^{\mathrm{dR}}_{S/S}(G)$ is de Rham-formal along its
identity section, and the construction is compatible with arbitrary base
change in $S$.
\end{proposition}

\begin{proof}
Relative de Rham localization is left exact, hence carries the finite-limit
diagram defining the group object $G$ to a group object
$G_{\mathrm{dR}/S}$, and naturality makes
\[
    \eta_{G/S}^{\mathrm{dR}}:G\longrightarrow G_{\mathrm{dR}/S}
\]
a group homomorphism.  The terminal object $S$ of $(\mathcal H_{\Sp})_{/S}$ carries
its unique group structure, and the map
$S\to G_{\mathrm{dR}/S}$ is the identity section of the localized group.
Therefore the pullback
\[
    \operatorname{Nbh}^{\mathrm{dR}}_{S/S}(G)
    =
    G\times_{G_{\mathrm{dR}/S}}S
\]
is a group object, with projection to $G$ a homomorphism.

Moreover
\[
    \operatorname{Nbh}^{\mathrm{dR}}_{S/S}(G)
    =
    P^{\mathrm{dR}}(e).
\]
Theorem~\ref{thm:centred-formalization-idempotent} therefore makes it de Rham-formal
along the identity section.  Base-change compatibility is
Proposition~\ref{prop:centred-formalization-base-change}, together with
base change of the group-object finite-limit diagrams.
\end{proof}

Now let
\[
    \alpha:G\times_SX\longrightarrow X
\]
be a left action.  Write $G\ltimes X$ for the action groupoid, with
\[
    s(g,x)=x,
    \qquad
    t(g,x)=\alpha(g,x).
\]

\begin{theorem}[Formalization of an action groupoid]
\label{thm:action-groupoid-formalization}
The $G$-action precomposes canonically along the homomorphism
$\operatorname{Nbh}^{\mathrm{dR}}_{S/S}(G)\to G$ to an action
\[
    \operatorname{Nbh}^{\mathrm{dR}}_{S/S}(G)\times_SX
    \longrightarrow X,
\]
and there is a canonical equivalence
\[
    P_X^{\mathrm{dR}}(G\ltimes X)
    \simeq
    \operatorname{Nbh}^{\mathrm{dR}}_{S/S}(G)\ltimes X.
\]
\end{theorem}

\begin{proof}
Proposition~\ref{prop:formal-identity-group} gives a canonical group
homomorphism
\[
    \operatorname{Nbh}^{\mathrm{dR}}_{S/S}(G)\longrightarrow G.
\]
Precomposing the $G$-action with this homomorphism constructs
\[
    \operatorname{Nbh}^{\mathrm{dR}}_{S/S}(G)\times_SX\longrightarrow X,
\]
and its unit, associativity, and higher coherences are the pullbacks of those
of the original action.

The degree-one centre is
\[
    u_1:X\longrightarrow G\times_SX,
    \qquad
    x\longmapsto(e,x).
\]
By left exactness,
\[
    (G\times_SX)_{\mathrm{dR}/S}
    \simeq
    G_{\mathrm{dR}/S}
    \times_S
    X_{\mathrm{dR}/S}.
\]
Hence
\[
\begin{aligned}
    P^{\mathrm{dR}}(u_1)
    &\simeq
    (G\times_SX)
    \times_{
        G_{\mathrm{dR}/S}\times_SX_{\mathrm{dR}/S}
    }
    X_{\mathrm{dR}/S}\\
    &\simeq
    \left(
        G\times_{G_{\mathrm{dR}/S}}S
    \right)
    \times_SX\\
    &=
    \operatorname{Nbh}^{\mathrm{dR}}_{S/S}(G)\times_SX.
\end{aligned}
\]
The formalized target map is the action obtained by precomposition along
$\operatorname{Nbh}^{\mathrm{dR}}_{S/S}(G)\to G$.  Pullback preservation of centred
formalization carries the Segal pullbacks of the action nerve to the Segal
pullbacks of this precomposed action, and hence identifies the formalized
simplicial groupoid with $\operatorname{Nbh}^{\mathrm{dR}}_{S/S}(G)\ltimes X$.
\end{proof}

\subsection{Principal group objects and formal gauge groupoids}

\begin{definition}[Classifying object of a group object]
\label{def:classifying-object-group}
Let $B_\bullet G$ be the one-object groupoid
\[
    (B_\bullet G)_0=S,
    \qquad
    (B_\bullet G)_n=G^{\times_Sn}\quad(n\geq1).
\]
We use the convention compatible with right action nerves: for a right
$G$-object $P$, the action nerve has
\[
    (P/\!/G)_n=P\times_SG^{\times_Sn},
\]
the first face applies the action by the first group element, interior faces
multiply adjacent group elements in that order, the last face forgets the
last group element, and degeneracies insert the identity.  The reduced bar
object $B_\bullet G$ is the corresponding one-object specialization.  Define
\[
    BG:=|B_\bullet G|.
\]
Groupoid effectivity gives an effective epimorphism
\[
    b_G:S\twoheadrightarrow BG,
    \qquad
    \check C_\bullet(b_G)\simeq B_\bullet G.
\]
We call $b_G$ the universal classifying basepoint.  If $S'\to S$ and
$G':=G\times_SS'$, pullback preserves the finite products defining
$B_\bullet G$ and the realization, so there are canonical coherent
equivalences
\[
    B_\bullet G'\simeq B_\bullet G\times_SS',
    \qquad
    BG'\simeq BG\times_SS',
\]
under which $b_{G'}$ is the pullback of $b_G$.
\end{definition}

\begin{definition}[Principal right group objects by classification]
\label{def:principal-object-space}
For $X\to S$ define the space of principal right $G$-objects by
\[
    \operatorname{Prin}_G(X)
    :=
    \Map_{(\mathcal H_{\Sp})_{/S}}(X,BG).
\]
For a point $\chi:X\to BG$, define its associated object
\[
    P_\chi:=X\times_{BG}S
    \xrightarrow{p_\chi}X.
\]
The right $G$-action on $P_\chi$ is not additional structure: it is the
degree-one structure obtained by pulling back the effective groupoid
$B_\bullet G\simeq\check C_\bullet(b_G)$ along $\chi$.
\end{definition}

\begin{proposition}[Constructed principal action and shear]
\label{prop:principal-classifying-map}
For every $\chi:X\to BG$, the projection
\[
    p_\chi:P_\chi\twoheadrightarrow X
\]
is an effective epimorphism and the pulled-back \v{C}ech equivalence gives a
canonical right $G$-action
\[
    \alpha_\chi:P_\chi\times_SG\longrightarrow P_\chi
\]
whose shear morphism is an equivalence
\[
    (\operatorname{pr}_1,\alpha_\chi):
    P_\chi\times_SG
    \xrightarrow{\simeq}
    P_\chi\times_XP_\chi.
\]
All action coherence is inherited from the pulled-back one-object groupoid.
The construction is natural in $\chi$ and under arbitrary base change in
$S$.
\end{proposition}

\begin{proof}
The projection $p_\chi$ is the base change of
$b_G:S\twoheadrightarrow BG$, hence is effective.  Pulling back the
identity
\[
    S\times_{BG}S\simeq G
\]
along $P_\chi\to S$ gives
\[
    P_\chi\times_XP_\chi
    \simeq
    P_\chi\times_SG.
\]
Under this equivalence the second projection is the constructed right action,
so the equivalence itself is the shear map.  Higher iterates are the higher
terms of the same pulled-back \v{C}ech nerve.
\end{proof}

For comparison with action-theoretic formulations, let
$\operatorname{Prin}^{\mathrm{rec}}_G(X)$ denote the maximal
$\infty$-groupoid of coherent right $G$-objects $P\to X$ whose projection is
an effective epimorphism and whose shear map is an equivalence.  This notation
is used only for the following recognition theorem; it is not the primary
definition of principality.

\begin{theorem}[Recognition of principal right group objects]
\label{thm:principal-objects-classified-by-bg}
The construction above induces a canonical equivalence
\[
    \operatorname{Prin}_G(X)
    =
    \Map_{(\mathcal H_{\Sp})_{/S}}(X,BG)
    \xrightarrow{\simeq}
    \operatorname{Prin}^{\mathrm{rec}}_G(X).
\]
Conversely, a coherent right $G$-object
$p:P\to X$ with effective projection and invertible shear canonically
constructs a classifying morphism
\[
    \chi_P:X\longrightarrow BG
\]
and a canonical equivalence of right $G$-objects
\[
    P\xrightarrow{\simeq}X\times_{BG}S.
\]
Thus effective epimorphism and invertible shear are recognition properties of
the objects constructed from classifying maps, not additional principal-object
axioms.
\end{theorem}

\begin{proof}
Let $P\to X$ satisfy the recognition conditions and let $P/\!/G$ be its action
groupoid.  Iterating the shear equivalence gives a simplicial equivalence
\[
    P/\!/G\xrightarrow{\simeq}\check C_\bullet(p).
\]
Since $p$ is effective,
\[
    |P/\!/G|\simeq X.
\]
Forgetting the $P$-coordinate gives a simplicial morphism
\[
    P/\!/G\longrightarrow B_\bullet G.
\]
Realization constructs
\[
    \chi_P:X\longrightarrow BG.
\]
The degree-zero square gives a canonical comparison
\[
\begin{tikzcd}
P \arrow[r] \arrow[d,"p"'] & S \arrow[d,"b_G"]\\
X \arrow[r,"\chi_P"'] & BG,
\end{tikzcd}
\]
and hence
\[
    \mathrm{cmp}^{\mathrm{prin}}_P:P\longrightarrow X\times_{BG}S.
\]
Pull $\mathrm{cmp}^{\mathrm{prin}}_P$ back along the effective epimorphism $p$.  On the source the shear
identifies
\[
    P\times_XP\simeq P\times_SG.
\]
On the target, using the displayed square and
$S\times_{BG}S\simeq G$, one obtains
\[
    P\times_X(X\times_{BG}S)
    \simeq
    P\times_SG.
\]
Tracing the action nerve identifies the pulled-back comparison with the shear
equivalence.  Effective pullback is conservative, so $\mathrm{cmp}^{\mathrm{prin}}_P$ is an equivalence.

Starting instead from $\chi$, the action groupoid of
$P_\chi=X\times_{BG}S$ is the pullback of
$B_\bullet G\simeq\check C_\bullet(b_G)$; realization recovers $\chi$ by
universality of colimits in the $\infty$-topos.  The two constructions are
therefore inverse on objects, morphisms, and all higher homotopies.
\end{proof}

\begin{definition}[Gauge groupoid]
\label{def:gauge-groupoid}
For $\chi:X\to BG$ define
\[
    \operatorname{Gauge}(\chi)
    :=
    \check C_\bullet(\chi)
    \rightrightarrows X.
\]
When $P_\chi=X\times_{BG}S$ is the associated principal object, we also write
$\operatorname{Gauge}(P_\chi)$.  Its arrow object is
\[
    \operatorname{Gauge}(P_\chi)_1=X\times_{BG}X.
\]
\end{definition}

\begin{proposition}[Gauge arrows as an effective diagonal quotient]
\label{prop:gauge-arrow-quotient}
Let $G$ act diagonally on $P_\chi\times_SP_\chi$.  The realization of this
action groupoid is canonically
\[
    (P_\chi\times_SP_\chi)/G
    \xrightarrow{\simeq}
    X\times_{BG}X.
\]
Under this equivalence the descended structure maps are the source, target,
unit, inversion, and composition maps of $\check C_\bullet(\chi)$.
\end{proposition}

\begin{proof}
Effective descent along $b_G:S\twoheadrightarrow BG$ is an equivalence of
$\infty$-categories and preserves finite limits.  The pullback of
$X\to BG$ is $P_\chi\to S$, so the pullback of
$X\times_{BG}X$ is $P_\chi\times_SP_\chi$, with descent transport the
diagonal $G$-action.  Applying the inverse descent equivalence gives the
asserted quotient formula and its simplicial structure.
\end{proof}

\begin{definition}[Formal gauge groupoid]
\label{def:formal-gauge-groupoid}
Define
\[
    \operatorname{Gauge}^{\mathrm{for}}(P_\chi)
    :=
    P_X^{\mathrm{dR}}\bigl(\operatorname{Gauge}(P_\chi)\bigr).
\]
\end{definition}

\begin{proposition}[Gauge groupoids give formal Lie groupoids]
\label{prop:gauge-formal-lie-groupoid}
The formal gauge groupoid is canonically an object of
$\operatorname{FormLieGpd}(X/S)$, natural in the classifying map and under
arbitrary base change.
\end{proposition}

\begin{proof}
Apply unitwise formalization to the explicit \v{C}ech groupoid
$\check C_\bullet(\chi)$.  Naturality is functoriality of pullback,
effective-image factorization, and centred formalization.
\end{proof}

\begin{remark}[Preview of the Atiyah comparison]
The formal gauge groupoid is a nonlinear object to which the following
differentiation stage may be applied once the required represented cotangent or
formal-moduli hypotheses have been verified.  No tangent fibre sequence or
Atiyah-type Lie algebroid is asserted in this chapter.
\end{remark}

\section{Group objects as the one-object special case}
\label{sec:flg-one-object-groups}

Let $G\to S$ be a group object and regard it as the one-object groupoid
\[
    B_\bullet G,
    \qquad
    (B_\bullet G)_0=S,
    \qquad
    (B_\bullet G)_1=G.
\]

\begin{theorem}[One-object de Rham formalization]
\label{thm:one-object-formalization}
The unitwise de Rham formalization of $B_\bullet G$ is the reduced bar groupoid of
\[
    \operatorname{Nbh}^{\mathrm{dR}}_{S/S}(G)
    :=
    G\times_{G_{\mathrm{dR}/S}}S.
\]
In particular, its degree-one object is canonically
\(\operatorname{Nbh}^{\mathrm{dR}}_{S/S}(G)\), and
\[
    P_S^{\mathrm{dR}}(B_\bullet G)
    \simeq
    B_\bullet\operatorname{Nbh}^{\mathrm{dR}}_{S/S}(G).
\]
\end{theorem}

\begin{proof}
The centre in degree one is the identity section
$e:S\to G$.  Since $S$ is terminal in $(\mathcal H_{\Sp})_{/S}$,
\[
    S_{\mathrm{dR}/S}\simeq S.
\]
Therefore
\[
    P^{\mathrm{dR}}(e)
    =
    G\times_{G_{\mathrm{dR}/S}}S
    =
    \operatorname{Nbh}^{\mathrm{dR}}_{S/S}(G).
\]
Pullback preservation of centred formalization carries the higher Segal
pullbacks over $S$ to the reduced bar construction of this formal identity
group.
\end{proof}

\begin{corollary}[Recognition of one-object formality]
\label{cor:one-object-formality-recognition}
The one-object groupoid $B_\bullet G$ is formal if and only if
\[
    \operatorname{Nbh}^{\mathrm{dR}}_{S/S}(G)
    \longrightarrow
    G
\]
is an equivalence.  Equivalently, the localized identity section
\[
    S
    \longrightarrow
    G_{\mathrm{dR}/S}
\]
is an equivalence.
\end{corollary}

\begin{proof}
The first criterion is
Proposition~\ref{prop:degree-one-formality} together with
Theorem~\ref{thm:one-object-formalization}.  The second is
Proposition~\ref{prop:centred-formalization-recognition} applied to
$e:S\to G$.
\end{proof}

\begin{corollary}[Classifying quotient of the formal identity group]
\label{cor:formal-identity-classifying-quotient}
The quotient of the formalized one-object groupoid is
\[
    B\operatorname{Nbh}^{\mathrm{dR}}_{S/S}(G),
\]
and the canonical classifying basepoint
\[
    S\twoheadrightarrow B\operatorname{Nbh}^{\mathrm{dR}}_{S/S}(G)
\]
is an effective epimorphism whose \v{C}ech nerve is
$B_\bullet\operatorname{Nbh}^{\mathrm{dR}}_{S/S}(G)$.  Pullback along this basepoint identifies
objects over $B\operatorname{Nbh}^{\mathrm{dR}}_{S/S}(G)$ with objects over $S$ equipped with
coherent $\operatorname{Nbh}^{\mathrm{dR}}_{S/S}(G)$-transport.
\end{corollary}

\begin{proof}
Apply groupoid effectivity to
$B_\bullet\operatorname{Nbh}^{\mathrm{dR}}_{S/S}(G)$ for the effective epimorphism and \v{C}ech
nerve statement.  Then apply
Theorem~\ref{thm:effective-relation-calculus} to that effective epimorphism;
its relation groupoid is precisely the one-object groupoid
$B_\bullet\operatorname{Nbh}^{\mathrm{dR}}_{S/S}(G)$, so its effective-descent equivalence is
the asserted classification by coherent transport.
\end{proof}

\begin{proposition}[One-object isotropy and adjoint structure]
\label{prop:one-object-isotropy}
For the formalized one-object groupoid associated to $G/S$,
\[
    I^{\mathrm{raw}}
    \simeq
    I^{\mathrm{for}}
    \simeq
    \operatorname{Nbh}^{\mathrm{dR}}_{S/S}(G).
\]
The nonlinear formal-isotropy commutator is the ordinary internal group
commutator of $\operatorname{Nbh}^{\mathrm{dR}}_{S/S}(G)$ and formal adjoint transport is
conjugation.
\end{proposition}

\begin{proof}
The object space is the terminal object $S$ of the slice, so the source-target
diagonal fibre is the entire arrow group
$\operatorname{Nbh}^{\mathrm{dR}}_{S/S}(G)$.  This arrow group is already the output of centred
formalization and is therefore formal.  The commutator and adjoint statements are
the general definitions specialized to a one-object groupoid.
\end{proof}

\begin{remark}
An arbitrary group object $G/S$ therefore determines a canonical one-object
formal Lie groupoid by unitwise formalization.  Only when
$\operatorname{Nbh}^{\mathrm{dR}}_{S/S}(G)\to G$ is an equivalence is the original
one-object groupoid itself formal.  Any later Lie-algebraic object is to be
obtained by differentiating the formalized groupoid in a range where a
differentiation theorem has been established.
\end{remark}

\section{Infinitesimal equality on a group and formal translation}
\label{sec:flg-group-translation}

We now recover the nonlinear translation theory of a group object as a theorem
inside the general groupoid framework.  Chapter~5 treated the translation
description for commutative group objects, where infinitesimal displacement
could be written additively.  The construction below is the noncommutative
extension: no commutativity is assumed, and additive differences are replaced
by the left and right group quotients $g^{-1}h$ and $hg^{-1}$.

Apply the canonical relative infinitesimal-equality construction to the object
$G\to S$:
\[
    \mathcal R^{\mathrm{dR}}_{G/S}
    =
    G\times_{G_{\mathrm{dR}/S}}G
    \rightrightarrows G.
\]
The established object
\[
    \operatorname{Nbh}^{\mathrm{dR}}_{S/S}(G)
    =
    P^{\mathrm{dR}}(e)
    =
    G\times_{G_{\mathrm{dR}/S}}S
\]
is the de Rham formal neighbourhood of the identity.
In the internal infinitesimal-equality notation,
\[
    g\sim_{\mathrm{dR}/S}h
\]
means precisely that $(g,h)$ is a point of
$\mathcal R^{\mathrm{dR}}_{G/S}$.  Translation sends such an infinitesimal equality
to an infinitesimal element at the identity.

\subsection{Translation trivializations}

\begin{theorem}[Group Translation Theorem]
\label{thm:group-translation}
Left and right multiplication define canonical equivalences over $G$
\[
\begin{aligned}
    \Phi^L_{G,1}:
    \mathcal R^{\mathrm{dR}}_{G/S}
    &\xrightarrow{\simeq}
    G\times_S\operatorname{Nbh}^{\mathrm{dR}}_{S/S}(G),
    &
    \Phi^L_{G,1}(g,h)&=(g,g^{-1}h),
\\
    \Phi^R_{G,1}:
    \mathcal R^{\mathrm{dR}}_{G/S}
    &\xrightarrow{\simeq}
    G\times_S\operatorname{Nbh}^{\mathrm{dR}}_{S/S}(G),
    &
    \Phi^R_{G,1}(g,h)&=(g,hg^{-1}).
\end{aligned}
\]
Their inverses are
\[
    \bigl(\Phi^L_{G,1}\bigr)^{-1}(g,u)=(g,gu),
    \qquad
    \bigl(\Phi^R_{G,1}\bigr)^{-1}(g,u)=(g,ug).
\]
\end{theorem}

\begin{proof}
Write
\[
    \eta_{G/S}^{\mathrm{dR}}:G\longrightarrow G_{\mathrm{dR}/S}
\]
for the relative de Rham unit.  Left exactness makes
$G_{\mathrm{dR}/S}$ a group object and naturality makes $\eta_{G/S}^{\mathrm{dR}}$ a group
homomorphism.

Let
\[
    d_L:
    G\times_SG\longrightarrow G,
    \qquad
    d_L=m_G(\iota_G\times\operatorname{id}_G)
\]
be the left-difference morphism.  On
\[
    \mathcal R^{\mathrm{dR}}_{G/S}
    =
    G\times_{G_{\mathrm{dR}/S}}G
\]
the defining pullback supplies a specified homotopy
\[
    \eta_{G/S}^{\mathrm{dR}}\operatorname{pr}_1
    \simeq
    \eta_{G/S}^{\mathrm{dR}}\operatorname{pr}_2.
\]
Since $\eta_{G/S}^{\mathrm{dR}}$ is a group homomorphism, this induces a canonical homotopy
\[
    \eta_{G/S}^{\mathrm{dR}}d_L
    \simeq
    \eta_{G/S}^{\mathrm{dR}}\circ e,
\]
so the universal property of
\[
    \operatorname{Nbh}^{\mathrm{dR}}_{S/S}(G)
    =
    G\times_{G_{\mathrm{dR}/S}}S
\]
constructs a canonical factorization
\[
    \overline d_L:
    \mathcal R^{\mathrm{dR}}_{G/S}\longrightarrow\operatorname{Nbh}^{\mathrm{dR}}_{S/S}(G).
\]
Together with the first projection it gives
\[
    \Phi^L_{G,1}
    :=
    (\operatorname{pr}_1,\overline d_L):
    \mathcal R^{\mathrm{dR}}_{G/S}\longrightarrow G\times_S\operatorname{Nbh}^{\mathrm{dR}}_{S/S}(G).
\]

Conversely, let $i_{\mathrm{Nbh}}:\operatorname{Nbh}^{\mathrm{dR}}_{S/S}(G)\to G$ be the pullback projection and define
\[
    \rho_L:
    G\times_S\operatorname{Nbh}^{\mathrm{dR}}_{S/S}(G)
    \longrightarrow
    G\times_SG
\]
by
\[
    \rho_L
    :=
    \left(
        \operatorname{pr}_G,\,
        m_G(\operatorname{pr}_G,i_{\mathrm{Nbh}}\operatorname{pr}_{\mathrm{Nbh}})
    \right).
\]
The defining homotopy $\eta_{G/S}^{\mathrm{dR}}i_{\mathrm{Nbh}}\simeq \eta_{G/S}^{\mathrm{dR}}\circ e$ implies that the
two components of $\rho_L$ have canonically equivalent de Rham images.
Therefore $\rho_L$ factors uniquely through
$\mathcal R^{\mathrm{dR}}_{G/S}$.  The group-object unit, inverse, and associativity
coherences give
\[
    g^{-1}(gu)\simeq u,
    \qquad
    g(g^{-1}h)\simeq h,
\]
as identities of the corresponding pullback morphisms.  Hence the two
factorizations are mutually inverse and $\Phi^L_{G,1}$ is an equivalence.

The right-difference morphism
\[
    d_R=m_G(\operatorname{id}_G\times\iota_G)
\]
gives, by the identical pullback argument, the equivalence $\Phi^R_{G,1}$ with
inverse induced by $(g,u)\mapsto(g,ug)$.  The generalized-element formulas in
the statement are shorthand for these internally constructed morphisms.
\end{proof}

\begin{theorem}[Higher translation coordinates]
\label{thm:higher-translation-coordinates}
For every $n\geq0$ there is a canonical equivalence
\[
    \Phi^L_{G,n}:
    \mathcal R^{\mathrm{dR}}_{G/S,n}
    \xrightarrow{\simeq}
    G\times_S\operatorname{Nbh}^{\mathrm{dR}}_{S/S}(G)^{\times_S n}
\]
given by
\[
    (g_0,\ldots,g_n)
    \longmapsto
    \left(
        g_0,
        g_0^{-1}g_1,
        g_1^{-1}g_2,
        \ldots,
        g_{n-1}^{-1}g_n
    \right).
\]
Its inverse is
\[
    (g,u_1,\ldots,u_n)
    \longmapsto
    \left(
        g,
        gu_1,
        gu_1u_2,
        \ldots,
        gu_1\cdots u_n
    \right).
\]
Under these equivalences,
\[
    \mathcal R^{\mathrm{dR}}_{G/S,\bullet}
    \simeq
    (G/\!/\operatorname{Nbh}^{\mathrm{dR}}_{S/S}(G))_\bullet
\]
for the right action $(g,u)\mapsto gu$.
\end{theorem}

\begin{proof}
Iterate Theorem~\ref{thm:group-translation} along the consecutive edges of
the relation simplex.  The resulting morphism
\[
    \mathcal R^{\mathrm{dR}}_{G/S,n}
    \longrightarrow
    G\times_S\operatorname{Nbh}^{\mathrm{dR}}_{S/S}(G)^{\times_Sn}
\]
has components the first vertex and the $n$ internally constructed left
differences.  Its inverse is obtained by iterated multiplication
\[
    (g,u_1,\ldots,u_n)
    \longmapsto
    (g,gu_1,gu_1u_2,\ldots,gu_1\cdots u_n).
\]
The group-object coherences identify the two composites with the identity.

It remains to identify the simplicial operators.  Under these coordinates,
the initial face applies the right action by $u_1$, an interior face replaces
$(u_i,u_{i+1})$ by $u_iu_{i+1}$, the terminal face forgets $u_n$, and a
degeneracy inserts the identity section.  By the right-action convention fixed
in Definition~\ref{def:classifying-object-group}, these are exactly the face
and degeneracy maps of $(G/\!/\operatorname{Nbh}^{\mathrm{dR}}_{S/S}(G))_\bullet$.  Hence the degreewise
equivalences assemble to the asserted simplicial equivalence.
\end{proof}

\begin{corollary}[Principal formal identity torsor]
\label{cor:principal-formal-identity-torsor}
There is a canonical quotient equivalence
\[
    \left|G/\!/\operatorname{Nbh}^{\mathrm{dR}}_{S/S}(G)\right|
    \simeq
    Q_{G/S}.
\]
Thus
\[
    G\twoheadrightarrow Q_{G/S}
\]
is a principal right $\operatorname{Nbh}^{\mathrm{dR}}_{S/S}(G)$-object in the ambient $\infty$-topos.
\end{corollary}

\begin{proof}
The map $G\twoheadrightarrow Q_{G/S}$ is an effective epimorphism by the
definition of the effective infinitesimal quotient.  By
Theorem~\ref{thm:higher-translation-coordinates}, the degree-one translation
equivalence identifies the shear map for the right $\operatorname{Nbh}^{\mathrm{dR}}_{S/S}(G)$-action,
\[
    G\times_S\operatorname{Nbh}^{\mathrm{dR}}_{S/S}(G)
    \longrightarrow
    G\times_{Q_{G/S}}G,
    \qquad
    (g,u)\longmapsto(g,gu),
\]
with an equivalence.  The recognition direction of
Theorem~\ref{thm:principal-objects-classified-by-bg} therefore identifies
$G\to Q_{G/S}$ with the principal right $\operatorname{Nbh}^{\mathrm{dR}}_{S/S}(G)$-object classified by the
realization of the action nerve.  The same theorem identifies
the entire action groupoid with the \v{C}ech groupoid of this effective
epimorphism, and realization gives
\[
    |G/\!/\operatorname{Nbh}^{\mathrm{dR}}_{S/S}(G)|\simeq Q_{G/S}.
\]
\end{proof}

\subsection{The integrated Maurer--Cartan cocycle}

\begin{definition}[Formal Maurer--Cartan projections]
\label{def:formal-mc-projections}
Define
\[
    \mathrm{MC}_{L,1}:
    \mathcal R^{\mathrm{dR}}_{G/S}\to\operatorname{Nbh}^{\mathrm{dR}}_{S/S}(G),
    \qquad
    \mathrm{MC}_{L,1}(g,h):=g^{-1}h,
\]
and
\[
    \mathrm{MC}_{R,1}:
    \mathcal R^{\mathrm{dR}}_{G/S}\to\operatorname{Nbh}^{\mathrm{dR}}_{S/S}(G),
    \qquad
    \mathrm{MC}_{R,1}(g,h):=hg^{-1}.
\]
\end{definition}

\begin{theorem}[Integrated group-valued Maurer--Cartan cocycle]
\label{thm:integrated-mc-cocycle}
For every $n\geq0$, the successive left increments define
\[
    \mathrm{MC}_{L,n}:
    \mathcal R^{\mathrm{dR}}_{G/S,n}\longrightarrow B_n\operatorname{Nbh}^{\mathrm{dR}}_{S/S}(G)=\operatorname{Nbh}^{\mathrm{dR}}_{S/S}(G)^{\times_Sn}
\]
by
\[
    \mathrm{MC}_{L,n}(g_0,\ldots,g_n)
    =
    \left(
        g_0^{-1}g_1,
        g_1^{-1}g_2,
        \ldots,
        g_{n-1}^{-1}g_n
    \right).
\]
These maps assemble into a simplicial morphism
\[
    \mathrm{MC}_{L,\bullet}:
    \mathcal R^{\mathrm{dR}}_{G/S,\bullet}\longrightarrow B_\bullet\operatorname{Nbh}^{\mathrm{dR}}_{S/S}(G).
\]
Its degree-one component is $\mathrm{MC}_{L,1}$, and its realization is the
classifying morphism
\[
    Q_{G/S}\longrightarrow B\operatorname{Nbh}^{\mathrm{dR}}_{S/S}(G)
\]
of the principal right $\operatorname{Nbh}^{\mathrm{dR}}_{S/S}(G)$-object $G\to Q_{G/S}$.

The successive right increments similarly define a simplicial morphism
\[
    \mathrm{MC}_{R,\bullet}:
    \mathcal R^{\mathrm{dR}}_{G/S,\bullet}
    \longrightarrow
    B_\bullet(\operatorname{Nbh}^{\mathrm{dR}}_{S/S}(G)^{\mathrm{op}}).
\]
\end{theorem}

\begin{proof}
For left increments $u_i=g_{i-1}^{-1}g_i$, deleting an interior vertex $g_i$
replaces $(u_i,u_{i+1})$ by
\[
    g_{i-1}^{-1}g_{i+1}
    =
    (g_{i-1}^{-1}g_i)(g_i^{-1}g_{i+1})
    =
    u_iu_{i+1}.
\]
Repeating a vertex inserts the identity.  These are the reduced bar face and
degeneracy maps, so the maps are simplicial.  By
Theorem~\ref{thm:higher-translation-coordinates}, the source simplicial object
is the action groupoid of the principal right $\operatorname{Nbh}^{\mathrm{dR}}_{S/S}(G)$-object
$G\to Q_{G/S}$.  The morphism $\mathrm{MC}_{L,\bullet}$ is exactly the simplicial
forgetful map from this action groupoid to $B_\bullet\operatorname{Nbh}^{\mathrm{dR}}_{S/S}(G)$; hence its
realization is the classifying morphism constructed in
Theorem~\ref{thm:principal-objects-classified-by-bg}.  For right increments the
multiplication order reverses, giving the opposite group.
\end{proof}

\begin{proposition}[Integrated Maurer--Cartan defect and Bianchi coherence]
\label{prop:integrated-curvature}
Define
\[
\begin{aligned}
    F_L(g_0,g_1,g_2)
    &:=
    \mathrm{MC}_{L,1}(g_0,g_1)
    \mathrm{MC}_{L,1}(g_1,g_2)
    \mathrm{MC}_{L,1}(g_0,g_2)^{-1},\\
    F_R(g_0,g_1,g_2)
    &:=
    \mathrm{MC}_{R,1}(g_1,g_2)
    \mathrm{MC}_{R,1}(g_0,g_1)
    \mathrm{MC}_{R,1}(g_0,g_2)^{-1}.
\end{aligned}
\]
Both admit canonical coherent trivializations at the identity of $\operatorname{Nbh}^{\mathrm{dR}}_{S/S}(G)$.  The
higher simplices of $\mathrm{MC}_{L,\bullet}$ and $\mathrm{MC}_{R,\bullet}$ provide the corresponding higher
Bianchi coherences.
\end{proposition}

\begin{proof}
For the left expression,
\[
    (g_0^{-1}g_1)(g_1^{-1}g_2)
    =
    g_0^{-1}g_2,
\]
and for the right expression,
\[
    (g_2g_1^{-1})(g_1g_0^{-1})
    =
    g_2g_0^{-1}.
\]
Thus both Maurer--Cartan defects are canonically the identity.  On a $3$-simplex the two
ways of composing three successive increments are related by associativity,
and the complete hierarchy is exactly the hierarchy of simplicial identities
of the bar cocycle.
\end{proof}

\begin{remark}[Why there is no general fixed-group Maurer--Cartan map]
\label{rem:no-general-fixed-group-mc}
The maps $\mathrm{MC}_{L,\bullet}$ and $\mathrm{MC}_{R,\bullet}$ exist because translation trivializes every
formal source fibre of the group object $G$ by the single identity fibre
$\operatorname{Nbh}^{\mathrm{dR}}_{S/S}(G)$.  For a general formal Lie groupoid the anchor fibres form only the varying
formal-isotropy bitorsors of Theorem~\ref{thm:formal-bitorsor}, and no canonical map to the bar construction of one fixed group
exists.  The general replacement is the groupoid nerve itself together with
its crystalline and adjoint transport.  Additive and operator-valued
Maurer--Cartan objects arise after differentiation and Spencer normalization.
\end{remark}

\section{The formal Lie groupoid interface}
\label{sec:flg-synthesis}

The constructions above form a single continuation of the differentiated
spectral doctrine.  The de Rham--formally \'etale factorization of Chapter~4
provides the fixed-centre operation; groupoid effectivity turns formal
relations into effective infinitesimal quotients; and the dependent adjoint
calculus of Chapter~5 turns those quotients into disk, jet, and crystalline
transport.  The following theorem collects the resulting nonlinear interface.

\begin{theorem}[Brave-New Formal Lie Groupoid Interface]
\label{thm:formal-lie-groupoid-package}
Let $X\to S$ be a morphism in $\mathcal H_{\Sp}$.  The preceding
constructions canonically provide the following.
\begin{enumerate}
    \item \textbf{Centred formal \'etalification.}
    For every $u:X\to Y$,
    \[
        P^{\mathrm{dR}}(u)
        =
        Y\times_{Y_{\mathrm{dR}/S}}X_{\mathrm{dR}/S}
    \]
    fits into
    \[
        X
        \xrightarrow{\ell_u}
        P^{\mathrm{dR}}(u)
        \xrightarrow{r_u}
        Y.
    \]
    When regarded on the fixed-source undercategory this is the endofunctor
    $P_X^{\mathrm{dR}}$.  Fixed-centre orthogonality makes
    $P_X^{\mathrm{dR}}$ the right adjoint to the inclusion of centred objects whose centre
    is a relative de Rham equivalence.  The resulting comonad is idempotent
    and preserves the pullbacks required by the Segal conditions.

    \item \textbf{Unitwise de Rham formalization.}
    Applying $P_X^{\mathrm{dR}}$ along the fully degenerate unit simplices gives an
    idempotent comonad
    \[
        P_X^{\mathrm{dR}}:
        \operatorname{Gpd}_X((\mathcal H_{\Sp})_{/S})
        \longrightarrow
        \operatorname{Gpd}_X((\mathcal H_{\Sp})_{/S}).
    \]
    A fixed-object groupoid is formal precisely when its unit
    \[
        u:X\to\mathfrak G_1
    \]
    is a relative de Rham equivalence, equivalently when the formalization
    counit is an equivalence.

    \item \textbf{Maximal infinitesimal equality and canonical anchor.}
    The canonical relative infinitesimal groupoid is
    \[
        \mathcal R^{\mathrm{dR}}_{X/S,\bullet}
        \simeq
        \check C_\bullet
        \left(
            X\twoheadrightarrow Q_{X/S}
        \right)
        \simeq
        P_X^{\mathrm{dR}}\bigl(\operatorname{Pair}(X/S)_\bullet\bigr).
    \]
    It is terminal among formal fixed-object groupoids.  Every
    $\mathfrak G\in\operatorname{FormLieGpd}(X/S)$ therefore has a
    contractibly unique canonical infinitesimal anchor
    \[
        \phi_{\mathfrak G}:
        \mathfrak G_\bullet
        \longrightarrow
        \mathcal R^{\mathrm{dR}}_{X/S,\bullet}.
    \]
    In internal notation this sends
    \[
        \gamma:x\sim_{\mathfrak G}y
        \longmapsto
        \phi_{\mathfrak G}(\gamma):
        x\sim_{\mathrm{dR}/S}y.
    \]

    \item \textbf{Effective infinitesimal quotient classification.}
    Realization and \v{C}ech nerve give canonical equivalences
    \[
        \operatorname{FormLieGpd}(X/S)
        \simeq
        \operatorname{EffEpi}^{\mathrm{for}}(X)
        \simeq
        \operatorname{Fact}^{\mathrm{for}}_{\mathrm{eff}}(e_{X/S}).
    \]
    Equivalently, a formal Lie groupoid is an effective epimorphism
    \[
        q:X\twoheadrightarrow Q
    \]
    which is a relative de Rham equivalence.  Its canonical anchor is the
    uniquely recovered factorization
    \[
        X
        \xrightarrow{q}
        Q
        \xrightarrow{\bar\phi_q}
        Q_{X/S},
        \qquad
        \bar\phi_q q\simeq e_{X/S},
    \]
    and both $q$ and $\bar\phi_q$ are effective relative de Rham equivalences.

    \item \textbf{Formal leaves and dependent infinitesimal neighbourhoods.}
    For
    $a:T\to X$,
    \[
        \mathsf{Disk}^{\infty}_{\mathfrak G}(a)
        =
        T\times_{Q_{\mathfrak G}}X
        \simeq
        T\times_{a,X,s}\mathfrak G_1.
    \]
    Fibrewise,
    \[
        \left(\mathsf{Disk}^{\infty}_{\mathfrak G}(a)\right)_t
        \simeq
        \sum_{x:X}
        \bigl(a(t)\sim_{\mathfrak G}x\bigr).
    \]
    The projection
    $\mathsf{Disk}^{\infty}_{\mathfrak G}(a)\to T$ is an effective relative de Rham
    equivalence, the centre is de Rham-formal, and the canonical anchor gives
    the comparison to the maximal relative de Rham disk of Chapter~5.

    \item \textbf{Raw inertia, formal isotropy, and adjoint transport.}
    Raw isotropy is the pulled-back relative inertia
    \[
        I^{\mathrm{raw}}_{\mathfrak G}
        \simeq
        q_{\mathfrak G}^*
        \mathcal I_{Q_{\mathfrak G}/S}.
    \]
    Formal isotropy is its centred formalization and equivalently the anchor
    kernel
    \[
        I^{\mathrm{for}}_{\mathfrak G}
        \simeq
        \mathfrak G_1
        \times_{\mathcal R^{\mathrm{dR}}_{X/S}}
        X
        \simeq
        q_{\mathfrak G}^*\mathcal K_{\mathfrak G}.
    \]
    Fibrewise this is
    \[
        \left\{
            \gamma:x\sim_{\mathfrak G}x
            \ \middle|\
            \phi_{\mathfrak G}(\gamma)\simeq1_x
        \right\}.
    \]
    The raw source--target and formal anchor fibres are the corresponding
    isotropy bitorsors, and the effective descent transport on
    $\mathcal K_{\mathfrak G}$ is exactly conjugation.  The nonlinear
    commutator is covariant under this adjoint transport.

    \item \textbf{Dependent disk--jet calculus.}
    For
    $q=q_{\mathfrak G}:X\twoheadrightarrow Q_{\mathfrak G}$,
    \[
        \mathsf{Disk}^{\infty}_{\mathfrak G}
        :=
        q^*\Sigma_q,
        \qquad
        J^\infty_{\mathfrak G}
        :=
        q^*\Pi_q,
        \qquad
        \mathsf{Disk}^{\infty}_{\mathfrak G}\dashv J^\infty_{\mathfrak G}.
    \]
    With $s,t:\mathfrak G_1\rightrightarrows X$,
    \[
        \mathsf{Disk}^{\infty}_{\mathfrak G}
        \simeq
        \Sigma_s t^*
        \simeq
        \Sigma_t s^*,
        \qquad
        J^\infty_{\mathfrak G}
        \simeq
        \Pi_s t^*.
    \]
    For a family $E\to X$,
    \[
        \bigl(\mathsf{Disk}^{\infty}_{\mathfrak G}E\bigr)_x
        \simeq
        \sum_{y:X}
        \sum_{\gamma:x\sim_{\mathfrak G}y}
        E_y,
    \]
    while
    \[
        \bigl(J^\infty_{\mathfrak G}E\bigr)_x
        \simeq
        \prod_{y:X}
        \prod_{\gamma:x\sim_{\mathfrak G}y}
        E_y.
    \]
    The maximal crystalline disk--jet calculus is recovered when
    $\mathfrak G=\mathcal R^{\mathrm{dR}}_{X/S}$.

    \item \textbf{Crystalline descent as substitution.}
    Effective descent gives
    \[
        (\mathcal H_{\Sp})_{/Q_{\mathfrak G}}
        \simeq
        \operatorname{Alg}_{\mathsf{Disk}^{\infty}_{\mathfrak G}}
        ((\mathcal H_{\Sp})_{/X})
        \simeq
        \operatorname{Coalg}_{J^\infty_{\mathfrak G}}
        ((\mathcal H_{\Sp})_{/X}).
    \]
    If $\bar E\to Q_{\mathfrak G}$ and $E=q^*\bar E$, then
    $\gamma:x\sim_{\mathfrak G}y$ induces
    \[
        \gamma_*:E_x\xrightarrow{\simeq}E_y
    \]
    by ordinary dependent substitution in the descended family.  This
    substitution is precisely the \v{C}ech transport and hence the disk
    action and jet coaction.

    \item \textbf{Functoriality, base change, and presentation scope.}
    Formalization, quotients, formal disks, raw and formal isotropy, inertia
    kernels, bitorsors, adjoint transport, disk monads, jet comonads, and
    crystalline descent are functorial in formal Lie groupoids and commute with
    arbitrary base change through the canonical Beck--Chevalley maps.
    Fixed-object quotient equivalences give equivalences of the entire
    package.  Unitwise de Rham formality is not invariant under arbitrary
    replacement of the object space by a different effective atlas, as shown
    by the split twofold-atlas calculation.

    \item \textbf{Groups, actions, principal objects, and gauge groupoids.}
    For every group object $G/S$,
    \[
        \operatorname{Nbh}^{\mathrm{dR}}_{S/S}(G)
        =
        G\times_{G_{\mathrm{dR}/S}}S
        =
        P^{\mathrm{dR}}(e)
    \]
    is the de Rham formal neighbourhood of the identity and is a group
    object.  Every $G$-action satisfies
    \[
        P_X^{\mathrm{dR}}(G\ltimes X)
        \simeq
        \operatorname{Nbh}^{\mathrm{dR}}_{S/S}(G)\ltimes X.
    \]
    Principal right $G$-objects are classified by maps $X\to BG$ and
    recognized by effective epimorphism plus invertible shear; their gauge
    groupoids are the corresponding \v{C}ech groupoids and admit canonical
    unitwise formalization.

    \item \textbf{Noncommutative infinitesimal translation.}
    The commutative translation result of Chapter~5 extends to arbitrary
    group objects by
    \[
        \mathcal R^{\mathrm{dR}}_{G/S,\bullet}
        \simeq
        (G/\!/\operatorname{Nbh}^{\mathrm{dR}}_{S/S}(G))_\bullet,
    \]
    using the left and right quotients $g^{-1}h$ and $hg^{-1}$.  Successive
    quotients give the integrated group-valued Maurer--Cartan cocycles and
    their canonically trivial Maurer--Cartan defects and higher Bianchi coherences.

    \item \textbf{Stable differential-cohomological handoff.}
    The formal-groupoid and crystalline-descent constructions above are nonlinear.  When a
    later pointed or linearized coefficient construction produces a stable
    infinitesimal coefficient from this data, the stable endpoint semantics of
    Chapter~2 supplies its BNV deformation and cycle sectors and its
    characteristic reconstruction square.  No pure-cycle, curvature, or
    differential-cohomology datum is added to the nonlinear definition
    before that stable coefficient theorem domain has been established.
\end{enumerate}
\end{theorem}

\begin{proof}
Item~(1) is
Proposition~\ref{prop:centred-formalization-recognition},
Theorem~\ref{thm:centred-formalization-coreflection}, and
Proposition~\ref{prop:centred-formalization-pullbacks}.
Item~(2) is
Proposition~\ref{prop:unitwise-formalization-groupoid},
Theorem~\ref{thm:unitwise-formalization-idempotent},
Proposition~\ref{prop:degree-one-formality}, and
Theorem~\ref{thm:formalization-coreflection}.
Item~(3) is
Theorems~\ref{thm:pair-groupoid-formalization} and~\ref{thm:maximal-formal-terminal} together with
Corollary~\ref{cor:canonical-formal-anchor}.
Item~(4) is
Theorems~\ref{thm:formal-quotient-monomorphism-recognition},~\ref{thm:quotient-formalization-coreflection}, and~\ref{thm:formal-gpd-quotient-classification}, together with
Proposition~\ref{prop:formal-quotient-canonical-anchor}.
Item~(5) is
Propositions~\ref{prop:formal-disk-arrow-fibre},~\ref{prop:formal-disk-dependent-presentation},~\ref{prop:formal-disk-centre-formality},~\ref{prop:formal-disk-effective}, and~\ref{prop:formal-disk-ambient-comparison}.
Item~(6) is
Proposition~\ref{prop:isotropy-inertia},
Theorems~\ref{thm:isotropy-group},~\ref{thm:formal-isotropy-descent},~\ref{thm:raw-bitorsor},~\ref{thm:formal-bitorsor},~\ref{thm:inertia-descent-conjugation}, and~\ref{thm:coherent-adjoint-transport}.
Item~(7) is
Theorem~\ref{thm:effective-relation-calculus},
Definition~\ref{def:formal-groupoid-disk-jet},
Proposition~\ref{prop:formal-groupoid-dependent-sum-product}, and
Proposition~\ref{prop:formal-groupoid-fibrewise-jets}.
Item~(8) is
Theorems~\ref{thm:formal-groupoid-transport-mates} and~\ref{thm:formal-groupoid-effective-descent}, together with
Definition~\ref{def:formal-cartan-objects} and
Proposition~\ref{prop:cartan-dependent-substitution}.
Item~(9) is
Propositions~\ref{prop:formal-gpd-functoriality-fixed-X},~\ref{prop:automatic-anchor-functoriality},~\ref{prop:presentation-invariance}, and~\ref{prop:formality-not-morita-invariant}, together with
Theorem~\ref{thm:formal-gpd-base-change}.
Item~(10) is
Proposition~\ref{prop:formal-identity-group},
Theorem~\ref{thm:action-groupoid-formalization},
Definitions~\ref{def:principal-object-space},~\ref{def:gauge-groupoid}, and~\ref{def:formal-gauge-groupoid}, together with
Theorem~\ref{thm:principal-objects-classified-by-bg}.
Item~(11) is
Theorems~\ref{thm:group-translation},~\ref{thm:higher-translation-coordinates}, and~\ref{thm:integrated-mc-cocycle}.
Item~(12) is the stable differential-cohomological interface recorded in
Section~\ref{sec:flg-descent}, using the exact-adjoint fracture theorem and
stable BNV reconstruction of Chapter~2.
\end{proof}


\section{First-order handoff to brave-new differentiation}
\label{sec:flg-handoff}
\label{sec:flg-represented-interface}

The intrinsic formal Lie groupoid theory does not require a represented
all-orders completion of the groupoid unit.  Chapter~3 already supplies a
canonical adic formal-order filtration for represented affine augmentations,
while Chapter~5 explains why a global finite-order jet tower requires a
separate globalization theorem.  Neither construction is needed for the
nonlinear formal-groupoid interface above.

The spectral-scheme interface used here is only first order.  Chapter~2 constructs
the canonical structured square-zero neighbourhood
\[
    X\xrightarrow{i_1}D^{(1)}_{X/S}\xrightarrow{\nu_1}X\times_SX
\]
with coefficient $L_{X/S}$, and Chapter~4 constructs its canonical map
\[
    \phi_1:D^{(1)}_{X/S}\longrightarrow
    \mathcal R^{\mathrm{dR}}_{X/S,1}
\]
over $X\times_SX$.  The final section compares the first structured
neighbourhood of a spectral-scheme groupoid unit with this diagonal neighbourhood
and proves that its coefficient is the first-order linearization of the
canonical nonlinear anchor.

\begin{remark}[Higher represented refinements]
Completed principal parts, distribution objects, refined Koszul duals,
partition-Lie controllers, PBW comparisons, and finite-order operator
filtrations belong to the later differentiation and Atiyah--Spencer stages.
Whenever such a spectral-geometric model is introduced, its comparison with the
intrinsic formal groupoid is a theorem rather than part of the definition.
\end{remark}




The nonlinear input for differentiation is now intrinsic.  A formal Lie
groupoid carries its canonical anchor
\[
    \phi_{\mathfrak G}:
    \mathfrak G_\bullet\longrightarrow \mathcal R^{\mathrm{dR}}_{X/S,\bullet},
\]
its effective formal quotient
\[
\begin{tikzcd}
X \arrow[r,two heads,"q_{\mathfrak G}"]
  \arrow[dr,two heads,"e_{X/S}"'] &
Q_{\mathfrak G}
  \arrow[d,two heads,"\bar\phi_{\mathfrak G}"]\\
&Q_{X/S},
\end{tikzcd}
\]
and its formal isotropy
\[
    I^{\mathrm{for}}_{\mathfrak G}
    \simeq
    \mathfrak G_1\times_{\mathcal R^{\mathrm{dR}}_{X/S}}X
    \simeq
    q_{\mathfrak G}^*\mathcal K_{\mathfrak G},
\]
with nonlinear commutator, bitorsor action, and coherent adjoint transport.
The formal orbit disks map canonically to the maximal de Rham formal disks, and
$\mathsf{Disk}^{\infty}_{\mathfrak G}\dashv J^\infty_{\mathfrak G}$ is a quotient of
the crystalline transport calculus of Chapter~5.

The stable slice doctrine of Chapters~3--5 exists over arbitrary objects of
$\mathcal H_{\Sp}$, but no identification of the stable shadow of an
arbitrary stack with a quasi-coherent tangent complex is built into that
doctrine.  The quasi-coherent first-order comparison is therefore made on the
spectral-scheme range, where Chapter~2 supplies cotangent
complexes and structured square-zero neighbourhoods and Chapter~4 embeds the
first diagonal neighbourhood into the de Rham relation.

\subsection{Representability from degrees zero and one}

\begin{proposition}[Spectral-scheme closure from degrees zero and one]
\label{prop:represented-groupoid-segal}
Let
\[
    \mathfrak G_\bullet\in
    \operatorname{Gpd}_X\bigl((\mathcal H_{\Sp})_{/S}\bigr)
\]
be a fixed-object groupoid.  Suppose
\[
    X=\mathfrak G_0,\qquad \mathfrak G_1
\]
lie in the full spectral-scheme subcategory
\[
    \operatorname{Sch}^{\mathrm{sp}}
    \hookrightarrow
    \mathcal H_{\Sp}
\]
of Chapter~1.  Then every $\mathfrak G_n$ lies in
$\operatorname{Sch}^{\mathrm{sp}}$, and the simplicial operators are morphisms in
that full subcategory.  Thus $\mathfrak G_\bullet$ is intrinsically a
simplicial spectral scheme; no representing presentation or atlas is chosen.
\end{proposition}

\begin{proof}
For $n\geq2$, the Segal equivalence gives
\[
    \mathfrak G_n
    \simeq
    \underbrace{
        \mathfrak G_1\times_X\cdots\times_X\mathfrak G_1
    }_{n\text{ factors}}.
\]
Chapter~1 proves that spectral schemes admit finite fibre products computed in
the ambient spectral Zariski $\infty$-topos.  Hence every degree belongs to
$\operatorname{Sch}^{\mathrm{sp}}$.  Since the inclusion
$\operatorname{Sch}^{\mathrm{sp}}\hookrightarrow\mathcal H_{\Sp}$ is fully
faithful, each simplicial operator of the ambient groupoid is uniquely the
image of a morphism of spectral schemes, and the simplicial identities are
unchanged.
\end{proof}

\subsection{The first structured neighbourhood of the groupoid unit}

Assume from now on that $S$, $X$, and $\mathfrak G_1$ lie in the full
spectral-scheme subcategory $\operatorname{Sch}^{\mathrm{sp}}\subset
\mathcal H_{\Sp}$.  Write
\[
    \mathfrak G_1
    \mathop{\rightrightarrows}^{s}_{t}
    X,
    \qquad
    u:X\longrightarrow\mathfrak G_1.
\]

\begin{definition}[First structured neighbourhood of the groupoid unit]
\label{def:first-unit-neighbourhood}
Define the first-order source coefficient
\[
    \mathfrak a_{\mathfrak G}^{*}
    :=
    u^*L_{\mathfrak G_1/X}.
\]
Since $su=\operatorname{id}_X$, cotangent transitivity for
\[
    X\xrightarrow{u}\mathfrak G_1\xrightarrow{s}X
\]
produces a canonical connecting equivalence
\[
    \partial_u:
    L_{X/\mathfrak G_1}
    \xrightarrow{\simeq}
    \mathfrak a_{\mathfrak G}^{*}[1].
\]
We use the same rotation convention as for the diagonal in Chapter~2.
Applying the structured square-zero construction of Chapter~2 gives a
canonical thickening
\[
    X\xrightarrow{i_u}D^{(1)}_{\mathfrak G}
    \xrightarrow{\nu_u}\mathfrak G_1
\]
with coefficient $\mathfrak a_{\mathfrak G}^{*}$.
\end{definition}

No smoothness or dualizability is used here.  The relative cotangent complex is
connective for spectral schemes by Chapter~2, so the coefficient lies in the domain of
the structured square-zero construction.

\subsection{First-order comparison with the diagonal}

Put
\[
    Y:=X\times_SX,
    \qquad
    f:=(s,t):\mathfrak G_1\longrightarrow Y,
    \qquad
    \Delta:=fu:X\longrightarrow Y.
\]
Chapter~2 constructed
\[
    X\xrightarrow{i_1}D^{(1)}_{X/S}
    \xrightarrow{\nu_1}X\times_SX
\]
with coefficient $L_{X/S}$, and Chapter~4 constructed the canonical map
\[
    \phi_1:
    D^{(1)}_{X/S}\longrightarrow \mathcal R^{\mathrm{dR}}_{X/S,1}
\]
over $X\times_SX$.

\begin{proposition}[Functoriality of structured first neighbourhoods]
\label{prop:first-neighbourhood-functoriality}
Let
\[
\begin{tikzcd}
X \arrow[r,"u"] \arrow[d,equal] &
Y \arrow[d,"f"] \arrow[r,"p"] &
X \arrow[d,equal]\\
X \arrow[r,"u'"'] &
Y' \arrow[r,"p'"'] &
X
\end{tikzcd}
\]
be a diagram of spectral schemes with
\[
    pu\simeq\operatorname{id}_X,
    \qquad
    p'u'\simeq\operatorname{id}_X,
    \qquad
    fu\simeq u',
    \qquad
    p'f\simeq p.
\]
Put
\[
    M_u:=u^*L_{Y/X},
    \qquad
    M_{u'}:=u'^*L_{Y'/X}.
\]
Let
\[
    X\xrightarrow{i_u}D^{(1)}_u\xrightarrow{\nu_u}Y,
    \qquad
    X\xrightarrow{i_{u'}}D^{(1)}_{u'}\xrightarrow{\nu_{u'}}Y'
\]
denote the structured first neighbourhoods obtained from the connecting
equivalences \(\partial_u\) and \(\partial_{u'}\) by the Chapter~2
construction.  The symbols \(D^{(1)}_u\) and \(D^{(1)}_{u'}\) are local to
this proposition and its proof.
Cotangent functoriality constructs a canonical morphism
\[
    \alpha_f:
    M_{u'}
    \longrightarrow
    M_u,
\]
and naturality of transitivity gives a canonical commutative square
\[
\begin{tikzcd}
L_{X/Y'} \arrow[r,"\partial_{u^{\prime}}"] \arrow[d] &
M_{u'}[1] \arrow[d,"{\alpha_f[1]}"]\\
L_{X/Y} \arrow[r,"\partial_u"'] &
M_u[1].
\end{tikzcd}
\]
Here the horizontal maps are the connecting equivalences determined by the
two retractions.

The square determines a canonical morphism of first structured
neighbourhoods
\[
    d^{(1)}f:
    D^{(1)}_u
    \longrightarrow
    D^{(1)}_{u'}
\]
over $f$, restricting to $\operatorname{id}_X$ on the common centre, whose
coefficient morphism is $\alpha_f$.  The construction is functorial under
identities and composition of diagrams of retractions and is compatible with
arbitrary Cartesian base change of spectral schemes.
\end{proposition}

\begin{proof}
The morphism
\[
    f:(Y,p)\longrightarrow(Y',p')
\]
over $X$ induces the canonical cotangent morphism
\[
    f^*L_{Y'/X}
    \longrightarrow
    L_{Y/X}.
\]
Pulling back along $u$ and using the specified homotopy
$fu\simeq u'$ gives
\[
\begin{aligned}
    M_{u'}
    =
    u'^*L_{Y'/X}
    &\simeq
    u^*f^*L_{Y'/X}\\
    &\longrightarrow
    u^*L_{Y/X}
    =
    M_u.
\end{aligned}
\]
This is $\alpha_f$.

Apply naturality of the global cotangent transitivity sequence to the
morphism between the two composable pairs
\[
    X\xrightarrow{u}Y\xrightarrow{p}X,
    \qquad
    X\xrightarrow{u'}Y'\xrightarrow{p'}X.
\]
Since both composites are identified with
$\operatorname{id}_X$, their middle cotangent terms vanish.  Rotating the two
transitivity cofibre sequences therefore gives the connecting equivalences
\[
    \partial_u:
    L_{X/Y}
    \xrightarrow{\simeq}
    M_u[1],
    \qquad
    \partial_{u'}:
    L_{X/Y'}
    \xrightarrow{\simeq}
    M_{u'}[1],
\]
and transitivity naturality gives the displayed commutative square.

We first construct the comparison on a compatible affine retractive chart.
Write
\[
    Y'=\Spec(A'),
    \qquad
    Y=\Spec(A),
    \qquad
    X=\Spec(B),
\]
so that the displayed diagram gives compatible algebra maps
\[
    A'\longrightarrow A\longrightarrow B
\]
together with the algebra maps corresponding to the two retractions.
The two connecting equivalences classify the relative derivation sections
\[
    d_{u'}:
    B\longrightarrow B\oplus M_{u'}[1],
    \qquad
    d_u:
    B\longrightarrow B\oplus M_u[1].
\]
The coefficient morphism
\[
    \alpha_f:M_{u'}\longrightarrow M_u
\]
induces a morphism of split square-zero algebras
\[
    1_B\oplus\alpha_f[1]:
    B\oplus M_{u'}[1]
    \longrightarrow
    B\oplus M_u[1].
\]
The transitivity-naturality square says precisely that
\[
    (1_B\oplus\alpha_f[1])\,d_{u'}
    \simeq
    d_u
\]
through the canonical homotopy; the same map carries the zero derivation to
the zero derivation.

The structured extension algebras are
\[
    B_{u'}
    :=
    B\times_{B\oplus M_{u'}[1]}B,
    \qquad
    B_u
    :=
    B\times_{B\oplus M_u[1]}B,
\]
where in each pullback the two maps from $B$ are the corresponding
derivation section and the zero section.  Functoriality of finite limits
therefore gives a canonical morphism
\[
    B_{u'}
    \longrightarrow
    B_u.
\]
Passing contravariantly to affine spectral schemes gives
\[
    D^{(1)}_u
    \longrightarrow
    D^{(1)}_{u'}.
\]
The defining diagrams show simultaneously that this morphism restricts to
$\operatorname{id}_X$ on the centre, lies over
$f:Y\to Y'$, and has coefficient morphism $\alpha_f$.

We now globalize this affine construction.  Choose an affine Zariski cover
\[
    \coprod_i U_i\longrightarrow X.
\]
After refining this cover, choose for every $i$ affine open spectral
subschemes
\[
    V_i\subset Y\times_XU_i,
    \qquad
    V'_i\subset Y'\times_XU_i
\]
such that the restricted sections factor through $V_i$ and $V'_i$ and the
restricted morphism $f$ carries $V_i$ into $V'_i$.  Such refinements exist
because affine opens form a Zariski basis, and the inverse image of an affine
open may be refined by affine opens.

On each resulting affine retractive diagram the preceding construction gives
a morphism of first structured neighbourhoods.  On an overlap, choose a
common compatible affine refinement.  Cotangent functoriality, transitivity
naturality, split square-zero formation, and finite-limit functoriality
identify the two restrictions with the same affine morphism on that common
refinement.  The Zariski-restriction theorem for structured square-zero
extensions from Chapter~2 identifies these affine restrictions with the
actual restrictions of
$D^{(1)}_u$ and $D^{(1)}_{u'}$.  Full faithfulness of spectral Zariski descent
for mapping spaces therefore glues the affine comparison maps through a
contractible space to a unique global morphism
\[
    d^{(1)}f:
    D^{(1)}_u
    \longrightarrow
    D^{(1)}_{u'}.
\]

For an identity diagram, every cotangent and finite-limit comparison above is
the identity, so
$d^{(1)}(\operatorname{id})\simeq\operatorname{id}$.
For composable diagrams
\[
    (Y,u,p)\xrightarrow{f}(Y',u',p')
    \xrightarrow{g}(Y'',u'',p''),
\]
cotangent functoriality gives
\[
    \alpha_{gf}
    \simeq
    \alpha_f\alpha_g
\]
with its canonical higher coherence.  The corresponding morphisms of split
square-zero diagrams therefore compose, by finite-limit functoriality, to
the morphism attached to $gf$.  Zariski descent transports this affine
composition comparison to the global one.  This proves functoriality and all
higher composition coherence.

Finally let the entire retractive diagram be pulled back along a Cartesian
base change of spectral schemes.  Global cotangent base change identifies
the pulled-back coefficient morphism with the coefficient morphism of the
pulled-back diagram, and Chapter~2 identifies the pullback of each structured
first neighbourhood with the first neighbourhood of the pulled-back
retraction.  On compatible affine charts the morphism constructed above is
therefore carried to the corresponding base-changed finite-limit morphism.
Full faithfulness of Zariski descent gives the canonical global base-change
comparison and its identity, composition, and higher-pasting coherences.
\end{proof}

\begin{proposition}[First-order source--target comparison]
\label{prop:first-order-source-target-comparison}
There is a canonical cotangent morphism
\[
    a_{\mathfrak G}^{*}:
    L_{X/S}\longrightarrow\mathfrak a_{\mathfrak G}^{*}
\]
and a canonical morphism of structured first neighbourhoods
\[
    d^{(1)}(s,t):
    D^{(1)}_{\mathfrak G}\longrightarrow D^{(1)}_{X/S}
\]
fitting into
\[
\begin{tikzcd}
X \arrow[r,"i_u"] \arrow[d,equal] &
D^{(1)}_{\mathfrak G}
    \arrow[r,"\nu_u"]
    \arrow[d,"{d^{(1)}(s,t)}"'] &
\mathfrak G_1
    \arrow[d,"{(s,t)}"]
\\
X \arrow[r,"i_1"'] &
D^{(1)}_{X/S}
    \arrow[r,"\nu_1"'] &
X\times_SX.
\end{tikzcd}
\]
The coefficient morphism of this structured square-zero comparison is
$a_{\mathfrak G}^{*}$, and the construction is compatible with Cartesian base change of spectral
schemes.
\end{proposition}

\begin{proof}
Cotangent base change for the first projection gives
\[
    L_{(X\times_SX)/X}
    \simeq
    \operatorname{pr}_2^*L_{X/S}.
\]
The morphism
\[
    (s,t):(\mathfrak G_1,s)
    \longrightarrow
    (X\times_SX,\operatorname{pr}_1)
\]
therefore induces
\[
    t^*L_{X/S}\longrightarrow L_{\mathfrak G_1/X}.
\]
Pulling back along $u$ and using $tu\simeq\operatorname{id}_X$ gives
\[
    a_{\mathfrak G}^{*}:
    L_{X/S}\longrightarrow u^*L_{\mathfrak G_1/X}.
\]

Apply Proposition~\ref{prop:first-neighbourhood-functoriality} to the
morphism of retractive contexts
\[
\begin{tikzcd}
X \arrow[r,"u"] \arrow[d,equal] &
\mathfrak G_1 \arrow[d,"{(s,t)}"] \arrow[r,"s"] & X \arrow[d,equal]\\
X \arrow[r,"\Delta"'] &
X\times_SX \arrow[r,"\operatorname{pr}_1"'] & X.
\end{tikzcd}
\]
Its coefficient morphism is exactly $a_{\mathfrak G}^{*}$, and the
resulting morphism of structured first neighbourhoods is the asserted
$d^{(1)}(s,t)$.  Base-change compatibility is part of the same functorial
construction.
\end{proof}

Equivalently, the cotangent anchor is the composite
\[
    L_{X/S}
    \simeq
    u^*t^*L_{X/S}
    \longrightarrow
    u^*L_{\mathfrak G_1/S}
    \longrightarrow
    u^*L_{\mathfrak G_1/X}.
\]

\subsection{Linearization of the canonical nonlinear anchor}

\begin{theorem}[First-order linearization of the canonical anchor]
\label{thm:first-order-linearization-canonical-anchor}
Let
$\mathfrak G_\bullet\in\operatorname{FormLieGpd}(X/S)$ and assume that
$S$, $X$, and $\mathfrak G_1$ lie in
$\operatorname{Sch}^{\mathrm{sp}}$.  Then there is a canonical homotopy
\[
    \phi_{\mathfrak G,1}\nu_u
    \simeq
    \phi_1d^{(1)}(s,t)
\]
through maps
\[
    D^{(1)}_{\mathfrak G}
    \longrightarrow
    \mathcal R^{\mathrm{dR}}_{X/S,1}
\]
over the fixed map $(s,t)\nu_u:D^{(1)}_{\mathfrak G}\to X\times_SX$, and the
space of such compatible homotopies is contractible.  Consequently
\[
    a_{\mathfrak G}^{*}:
    L_{X/S}\longrightarrow\mathfrak a_{\mathfrak G}^{*}
\]
is the first structured coefficient of the canonical nonlinear infinitesimal
anchor.
\end{theorem}

\begin{proof}
The centre
\[
    i_u:X\longrightarrow D^{(1)}_{\mathfrak G}
\]
is a structured square-zero infinitesimal extension and hence a relative de
Rham equivalence by Chapter~4.  Thus $D^{(1)}_{\mathfrak G}$ is de Rham-formal along
its centre.  The centred formalization coreflection therefore makes
composition with
\[
    \mathcal R^{\mathrm{dR}}_{X/S,1}
    =P^{\mathrm{dR}}(\Delta)
    \longrightarrow X\times_SX
\]
an equivalence on mapping spaces out of $D^{(1)}_{\mathfrak G}$.  The homotopy
fibre over $(s,t)\nu_u$ is consequently contractible.

The map $\phi_{\mathfrak G,1}\nu_u$ is one point of this lift space because
the canonical anchor is the formalized source--target map.  The composite
$\phi_1d^{(1)}(s,t)$ is another because $\phi_1$ lies over
$\nu_1:D^{(1)}_{X/S}\to X\times_SX$.  Contractibility gives the stated
homotopy and its coherence.  Since $d^{(1)}(s,t)$ was constructed from
$a_{\mathfrak G}^{*}$, that morphism is precisely the first structured
coefficient of the nonlinear anchor.
\end{proof}

\subsection{Internal-dual tangent anchor}

\begin{theorem}[Spectral-scheme internal-Hom first-order anchor]
\label{thm:represented-internal-dual-anchor}
The cotangent anchor
\[
    a_{\mathfrak G}^{*}:
    L_{X/S}\longrightarrow\mathfrak a_{\mathfrak G}^{*}
\]
induces by precomposition a canonical morphism
\[
\boxed{
    \underline{\operatorname{Hom}}_X
    \bigl(
        \mathfrak a_{\mathfrak G}^{*},
        \mathcal O_X
    \bigr)
    \longrightarrow
    \underline{\operatorname{Hom}}_X
    \bigl(
        L_{X/S},
        \mathcal O_X
    \bigr).
}
\]
No smoothness or dualizability is required for this internal-Hom morphism.

For a Cartesian base change of spectral schemes
\[
    c:S'\longrightarrow S,
    \qquad
    X':=X\times_SS',
    \qquad
    \mathfrak G'_\bullet:=\mathfrak G_\bullet\times_SS',
\]
write $c_X:X'\to X$ for the projection.  There are canonical comparison
morphisms
\[
    c_X^*
    \underline{\operatorname{Hom}}_X
    \bigl(
        \mathfrak a_{\mathfrak G}^{*},
        \mathcal O_X
    \bigr)
    \longrightarrow
    \underline{\operatorname{Hom}}_{X'}
    \bigl(
        \mathfrak a_{\mathfrak G'}^{*},
        \mathcal O_{X'}
    \bigr)
\]
and
\[
    c_X^*
    \underline{\operatorname{Hom}}_X
    \bigl(
        L_{X/S},
        \mathcal O_X
    \bigr)
    \longrightarrow
    \underline{\operatorname{Hom}}_{X'}
    \bigl(
        L_{X'/S'},
        \mathcal O_{X'}
    \bigr),
\]
compatible with the precomposition morphisms.
\end{theorem}

\begin{proof}
The closed presentably symmetric monoidal structure on
$\QCoh(X)$ gives internal Homs for arbitrary quasi-coherent
modules.  Precomposition with
\[
    a_{\mathfrak G}^{*}:
    L_{X/S}\longrightarrow\mathfrak a_{\mathfrak G}^{*}
\]
gives the displayed morphism.

Strong symmetric monoidality of pullback supplies the canonical comparison
\[
    c_X^*\underline{\operatorname{Hom}}_X(M,\mathcal O_X)
    \longrightarrow
    \underline{\operatorname{Hom}}_{X'}
    \bigl(c_X^*M,\mathcal O_{X'}\bigr)
\]
for every $M\in\QCoh(X)$.  Combining this with cotangent base change and
the base-change comparison
\[
    c_X^*\mathfrak a_{\mathfrak G}^{*}
    \xrightarrow{\simeq}
    \mathfrak a_{\mathfrak G'}^{*}
\]
gives the two asserted morphisms.  Naturality of evaluation gives their
compatibility with precomposition.
\end{proof}

\begin{remark}[Internal Hom versus tangent dual]
\label{rem:internal-dual-not-tangent-representing}
For an arbitrary quasi-coherent module $M$, the internal Hom
\[
    \underline{\operatorname{Hom}}_X(M,\mathcal O_X)
\]
exists, but $M$ need not be dualizable.  We therefore do not denote this
object by $M^\vee$ outside the dualizable range.  In particular, the
internal-Hom objects above need not tensor-represent the corresponding
tangent functors.  The finite-locally-free theorem below is the range in
which the established dual notation ${}^\vee$ applies and the
precomposition morphism becomes the usual tangent-module anchor.
\end{remark}

\begin{theorem}[Finite-locally-free spectral-scheme first-order anchor]
\label{thm:represented-first-order-anchor}
Assume that $L_{X/S}$ and
$L_{\mathfrak G_1/X}$ are finite locally free under the standing
spectral-scheme hypotheses.  Then
\[
    \mathfrak a_{\mathfrak G}^{*}
    =
    u^*L_{\mathfrak G_1/X}
\]
is finite locally free.  In the established dualizable notation define
\[
    \mathfrak a_{\mathfrak G}
    :=
    \left(
        \mathfrak a_{\mathfrak G}^{*}
    \right)^\vee.
\]
Then $\mathfrak a_{\mathfrak G}$ and $L_{X/S}^{\vee}$ are finite locally free,
and the precomposition morphism of
Theorem~\ref{thm:represented-internal-dual-anchor} is the canonical tangent
anchor
\[
\boxed{
    a_{\mathfrak G}:
    \mathfrak a_{\mathfrak G}
    \longrightarrow
    L_{X/S}^{\vee}.
}
\]

For every Cartesian base change of spectral schemes
\[
\begin{tikzcd}
X' \arrow[r,"c_X"] \arrow[d] &
X \arrow[d]\\
S' \arrow[r,"c"'] &
S
\end{tikzcd}
\]
put
\[
    \mathfrak G'_\bullet
    :=
    \mathfrak G_\bullet\times_SS'.
\]
Then there are canonical equivalences
\[
    c_X^*\mathfrak a_{\mathfrak G}^{*}
    \xrightarrow{\simeq}
    \mathfrak a_{\mathfrak G'}^{*},
    \qquad
    c_X^*\mathfrak a_{\mathfrak G}
    \xrightarrow{\simeq}
    \mathfrak a_{\mathfrak G'},
\]
and
\[
    c_X^*L_{X/S}^{\vee}
    \xrightarrow{\simeq}
    L_{X'/S'}^{\vee}.
\]
Under these equivalences the pulled-back anchor is canonically the anchor of
the base-changed formal Lie groupoid:
\[
\begin{tikzcd}
c_X^*\mathfrak a_{\mathfrak G}
    \arrow[r,"{c_X^*a_{\mathfrak G}}"]
    \arrow[d,"\simeq"']
&
c_X^*L_{X/S}^{\vee}
    \arrow[d,"\simeq"]
\\
\mathfrak a_{\mathfrak G'}
    \arrow[r,"{a_{\mathfrak G^{\prime}}}"']
&
L_{X'/S'}^{\vee}.
\end{tikzcd}
\]
All these equivalences are coherent under identities, composition, and
iterated Cartesian base change.
\end{theorem}

\begin{proof}
Pullback along $u$ preserves finite locally free quasi-coherent modules.
Hence the finite local freeness of
$L_{\mathfrak G_1/X}$ implies that
\[
    \mathfrak a_{\mathfrak G}^{*}
    =
    u^*L_{\mathfrak G_1/X}
\]
is finite locally free.  Finite locally free quasi-coherent modules are
dualizable, so
\[
    \mathfrak a_{\mathfrak G}
    =
    \left(
        \mathfrak a_{\mathfrak G}^{*}
    \right)^\vee
\]
and $L_{X/S}^{\vee}$ are the ordinary dual tangent modules in this theorem
domain.  The precomposition morphism of
Theorem~\ref{thm:represented-internal-dual-anchor} therefore identifies with
\[
    a_{\mathfrak G}:
    \mathfrak a_{\mathfrak G}
    \longrightarrow
    L_{X/S}^{\vee}.
\]

Now let $c:S'\to S$ be arbitrary and form the displayed Cartesian base
change.  Put
\[
    \mathfrak G'_1
    :=
    \mathfrak G_1\times_SS',
\]
and let
\[
    c_1:\mathfrak G'_1\longrightarrow\mathfrak G_1
\]
be the projection.  The source morphisms form a Cartesian square
\[
\begin{tikzcd}
\mathfrak G'_1 \arrow[r,"c_1"] \arrow[d,"s'"'] &
\mathfrak G_1 \arrow[d,"s"]\\
X' \arrow[r,"c_X"'] &
X.
\end{tikzcd}
\]
The global cotangent base-change theorem of Chapter~2 gives canonical
equivalences
\[
    c_X^*L_{X/S}
    \xrightarrow{\simeq}
    L_{X'/S'},
\]
and
\[
    c_1^*L_{\mathfrak G_1/X}
    \xrightarrow{\simeq}
    L_{\mathfrak G'_1/X'}.
\]
If
\[
    u':X'\longrightarrow\mathfrak G'_1
\]
is the base-changed unit, then
$c_1u'\simeq uc_X$.  Pulling the second equivalence along $u'$ therefore
gives
\[
\begin{aligned}
    c_X^*\mathfrak a_{\mathfrak G}^{*}
    &=
    c_X^*u^*L_{\mathfrak G_1/X}\\
    &\simeq
    u'^*c_1^*L_{\mathfrak G_1/X}\\
    &\simeq
    u'^*L_{\mathfrak G'_1/X'}\\
    &=
    \mathfrak a_{\mathfrak G'}^{*}.
\end{aligned}
\]

Since the relevant cotangent modules are finite locally free, they and their
pullbacks are dualizable.  Strong symmetric monoidality of pullback therefore
identifies pullback of the dual with the dual of the pullback.  Consequently
\[
    c_X^*\mathfrak a_{\mathfrak G}
    \simeq
    \mathfrak a_{\mathfrak G'},
    \qquad
    c_X^*L_{X/S}^{\vee}
    \simeq
    L_{X'/S'}^{\vee}.
\]

Proposition~\ref{prop:first-order-source-target-comparison} constructs the
cotangent anchor functorially from the source--target morphism
\[
    (s,t):\mathfrak G_1\longrightarrow X\times_SX.
\]
Cotangent base change therefore identifies the pullback of
$a_{\mathfrak G}^{*}$ with
$a_{\mathfrak G'}^{*}$.  Dualizing this canonical comparison gives the
commutative square of tangent anchors displayed in the theorem.

The identity, composition, and higher-pasting coherences are those of global
cotangent base change, strong symmetric-monoidal pullback, and functorial
duality on dualizable objects.
\end{proof}

\begin{remark}[Standard smooth recognition]
\label{rem:standard-smooth-recognition}
No unnamed smoothness axiom is used in the preceding theorem:
finite local freeness of
\[
    L_{X/S}
    \qquad\text{and}\qquad
    L_{\mathfrak G_1/X}
\]
is the explicit theorem hypothesis.

On the standard spectral-geometric theorem domain, Lurie,
\emph{Spectral Algebraic Geometry}, Section~11.2 develops formal and
differential smoothness for connective $E_\infty$-rings, gives the relevant
affine cotangent-complex criteria, and subsequently formulates differential
smoothness for spectral Deligne--Mumford stacks.  The broader
infinitesimal-cohesive deformation theory and its lifting formalism are
developed later in Chapter~17, in particular Section~17.3.

Accordingly, a named smoothness condition may replace one or both of the
explicit finite-local-freeness hypotheses above only after invoking a
specific external recognition theorem whose stated hypotheses imply the
required cotangent finiteness on the corresponding morphisms.  No such
smoothness package is imported into the intrinsic definition of a formal Lie
groupoid.
\end{remark}

The chapter has therefore reached the exact interface required for
brave-new differentiation.  The nonlinear formal groupoid, its formal isotropy
kernel, formal orbit disks, crystalline descent, and canonical anchor are already
constructed without linear data.  On spectral schemes the first
structured coefficient of that anchor is the cotangent morphism
\[
    L_{X/S}\longrightarrow\mathfrak a_{\mathfrak G}^{*}.
\]
The next chapter differentiates the entire nonlinear interface.  It will keep
separate the stable shadow of the full arrow geometry, the differentiation
of actual formal isotropy, the spectral Lie algebra obtained from derived
primitives of the isotropy Hopf object, and the partition-Lie algebroid
controller arising from spectral formal-moduli recognition.  None of those
objects is inserted into the present definition of a formal Lie groupoid.

